% Séismes — SVTEEHB 3ème — Cours signé OnBuch+
% Programme officiel MINESEC. Compiler avec : tectonic N12-seismes.tex
\newcommand{\DOCMATIERE}{SVTEEHB}
\newcommand{\DOCNIVEAU}{3ème}
\newcommand{\DOCMODULE}{SVTEEHB – classe de 3ème}
\newcommand{\DOCLECON}{N12}
\newcommand{\DOCTITRE}{Séismes}
% OnBuch+ — préambule « cours signés » ENRICHI (compilé avec tectonic / XeLaTeX)
% Système visuel « Pop » d'OnBuch+ : fond crème (jamais blanc), bordures encre
% épaisses, ombres « dures » décalées, titres Archivo Black en capitales.
% Version MATHÉMATIQUES : graphes (pgfplots/tikz), boîtes pédagogiques
% (définition, propriété, méthode, attention, le savais-tu, exercice résolu,
% exercice, TP, bilan), page de garde et signature OnBuch+ ancrée en bas.
%
% Macros attendues dans main.tex AVANT \input{preamble} :
%   \DOCMATIERE  \DOCNIVEAU  \DOCMODULE  \DOCLECON (numéro)  \DOCTITRE
\documentclass[11pt]{article}

\usepackage{fontspec}
\usepackage[french]{babel}
\usepackage[a4paper,margin=2cm,top=3.4cm,bottom=2.8cm,headheight=1.4cm,footskip=1.3cm]{geometry}
\usepackage{amsmath,amssymb,amsfonts}
\usepackage{mathtools} % \xmapsto, \coloneqq... souvent utilisés par les modèles
\usepackage{cancel} % simplifications barrées (bilans de forces)
% \abs{z} (module, valeur absolue) et \norm{u} (norme), utilisés par les modèles
\providecommand{\abs}{}\let\abs\relax\DeclarePairedDelimiter{\abs}{\lvert}{\rvert}
\providecommand{\norm}{}\let\norm\relax\DeclarePairedDelimiter{\norm}{\lVert}{\rVert}
\usepackage[table]{xcolor} % [table] : fournit aussi \rowcolors (alternance de lignes)
\usepackage{pagecolor}
\usepackage{tikz}
\usetikzlibrary{arrows.meta,positioning,calc,shapes.geometric,shapes.misc,shapes.arrows,decorations.pathmorphing,decorations.pathreplacing,decorations.markings,patterns,shadows,mindmap,trees,fit,backgrounds,matrix,intersections,angles,quotes}
\usepackage{pgfplots}
\usepackage[version=4]{mhchem} % équations nucléaires \ce{^{235}_{92}U}
\usepackage[european,siunitx]{circuitikz} % schémas électriques
\pgfplotsset{compat=1.18}
\usepgfplotslibrary{fillbetween,groupplots} % aires entre courbes (calcul intégral)
\usepackage{enumitem}
\usepackage{fancyhdr}
% [most] charge toutes les bibliothèques tcolorbox (listings, minted, raster,
% documentation...) et gonfle inutilement la mémoire TeX sur un document long
% (~300 boîtes) : on ne charge que ce qui est réellement utilisé.
\usepackage[skins,breakable]{tcolorbox}
\usepackage{graphicx}
\usepackage{colortbl}
\usepackage{booktabs}
\usepackage{tabularx}
\usepackage{multirow}
\usepackage{diagbox} % case d'en-tête diagonale des tableaux à double entrée
% Colonnes extensibles centrée / gauche / droite (C, L, R), souvent utilisées par les modèles
\newcolumntype{C}{>{\centering\arraybackslash}X}
\newcolumntype{L}{>{\raggedright\arraybackslash}X}
\newcolumntype{R}{>{\raggedleft\arraybackslash}X}
\usepackage{array}
\usepackage{multicol}
\usepackage{siunitx}
% parse-numbers=false : accepte aussi \SI{2,5 \times 10^{-3}}{...}
\sisetup{locale=FR,output-decimal-marker={,},group-minimum-digits=4,parse-numbers=false}
\DeclareSIUnit\ohm{\ensuremath{\symup{\Omega}}} % U+2126 absent de la police
\DeclareSIUnit\division{div} % réglages d'oscilloscope (ms/div, V/div)
\DeclareSIUnit\tr{tr} % tours (vitesse de rotation, ex. \SI{1500}{\tr\per\minute})
\DeclareSIUnit\molar{M} % concentration molaire (ex. \SI{3}{\molar}, \SI{1}{\micro\molar})
\DeclareSIUnit\an{an} % année (le modèle confond parfois avec \year, un compteur TeX) — ex. \SI{5}{\kilo\meter\per\an}
\DeclareSIUnit\FCFA{FCFA} % franc CFA (ex. \SI{500}{\FCFA\per\kilo\gram})
\DeclareSIUnit\degreeBrix{\textdegree Brix} % teneur en sucre d'un sirop/jus (réfractométrie)
\DeclareSIUnit\month{mois} % le modèle utilise parfois \month, un compteur TeX, comme unité de durée
\DeclareSIUnit\microliter{\micro\liter} % le modèle écrit parfois \microliter en un seul token
\DeclareSIUnit\million{millions} % dénombrement (ex. \SI{15}{\million\per\mL} spermatozoïdes)
\usepackage{qrcode}
% \degree : commande fréquemment utilisée par les modèles pour un angle ou une
% température en dehors de \SI{}{} (non fournie par les paquets chargés).
\providecommand{\degree}{^{\circ}}
% \qed (fin de démonstration), utilisé par les modèles sans amsthm
\providecommand{\qed}{\hfill$\square$}
% \barre : double barre des valeurs interdites dans un tableau de variations (tkz-tab)
\providecommand{\barre}{\ensuremath{\|}}
% environnement proof (amsthm non chargé) utilisé par les modèles
\ifdefined\proof\else\newenvironment{proof}[1][Démonstration]{\par\noindent\textit{#1.}\ }{\qed\par}\fi
\usepackage{lastpage}
\usepackage{hyperref}
\hypersetup{hidelinks}

% ============================================================
% Palette « Pop » OnBuch+ — mêmes hex que la classe P (plus_ui.dart)
% ============================================================
\definecolor{poporange}{HTML}{F59321}
\definecolor{popdark}{HTML}{E07A0C}
\definecolor{popcream}{HTML}{FFF6E8}
\definecolor{popink}{HTML}{1C1714}
\definecolor{poppurple}{HTML}{7A5AE0}
\definecolor{popgreen}{HTML}{2E9E6A}
\definecolor{poppink}{HTML}{E0457B}
\definecolor{popblue}{HTML}{2F7DE1}
\definecolor{popgold}{HTML}{C98A12}
\definecolor{popmuted}{HTML}{8A7F76}
\definecolor{popline}{HTML}{EADFD2}
\definecolor{popgray}{HTML}{8A7F76}
\colorlet{popgrey}{popgray}
% Alias tolérants (le modèle nomme parfois les couleurs différemment)
\colorlet{popwhite}{white}
\colorlet{popblack}{popink}
\colorlet{popred}{poppink}
\colorlet{popyellow}{popgold}
% Teintes claires (fonds de tableaux/encadrés) — le modèle nomme parfois
% les couleurs avec un suffixe L (ex. popblueL) sans les avoir définies.
\colorlet{poporangeL}{poporange!20}
\colorlet{popdarkL}{popdark!20}
\colorlet{poppurpleL}{poppurple!20}
\colorlet{popgreenL}{popgreen!20}
\colorlet{poppinkL}{poppink!20}
\colorlet{popblueL}{popblue!20}
\colorlet{popgoldL}{popgold!20}
\colorlet{popredL}{popred!20}
\colorlet{popgrayL}{popgray!20}
% Teintes claires pour fonds de tableaux / zones
\colorlet{popblueL}{popblue!10!white}
\colorlet{poporangeL}{poporange!14!white}
\colorlet{popgreenL}{popgreen!12!white}
\colorlet{poppurpleL}{poppurple!12!white}
\colorlet{poppinkL}{poppink!12!white}
\colorlet{popgoldL}{popgold!14!white}

\pagecolor{popcream}

% ============================================================
% Typographie
% ============================================================
\defaultfontfeatures{Ligatures=TeX}
\setmainfont{PlusJakartaSans-Medium.ttf}[Path=../../../fonts/,BoldFont=PlusJakartaSans-Bold.ttf]
% Maths en sans-serif (Fira Math) avec chiffres et lettres droites de Plus
% Jakarta, pour que les formules se fondent dans le texte.
\usepackage[math-style=ISO,bold-style=ISO]{unicode-math}
\setmathfont{FiraMath-Regular.otf}[Path=../../../fonts/]
\setmathfont{PlusJakartaSans-Medium.ttf}[Path=../../../fonts/,range={up/num,up/{Latin,latin},\mathup}]
\usepackage{icomma} % virgule décimale sans espace en mode math
% Caractères mathématiques Unicode absents des polices de texte (Plus Jakarta,
% Archivo Black) : rendus par leur équivalent mathématique, en texte comme en
% formule — sinon ils s'affichent en carré vide (ex. « Arithmétique dans ℤ »).
\usepackage{newunicodechar}
\newunicodechar{ℝ}{\ensuremath{\mathbb{R}}}
\newunicodechar{ℤ}{\ensuremath{\mathbb{Z}}}
\newunicodechar{ℂ}{\ensuremath{\mathbb{C}}}
\newunicodechar{ℕ}{\ensuremath{\mathbb{N}}}
\newunicodechar{ℚ}{\ensuremath{\mathbb{Q}}}
\newunicodechar{𝔻}{\ensuremath{\mathbb{D}}}
\newunicodechar{α}{\ensuremath{\alpha}}
\newunicodechar{β}{\ensuremath{\beta}}
\newunicodechar{γ}{\ensuremath{\gamma}}
\newunicodechar{δ}{\ensuremath{\delta}}
\newunicodechar{ε}{\ensuremath{\varepsilon}}
\newunicodechar{θ}{\ensuremath{\theta}}
\newunicodechar{λ}{\ensuremath{\lambda}}
\newunicodechar{μ}{\ensuremath{\mu}}
\newunicodechar{σ}{\ensuremath{\sigma}}
\newunicodechar{τ}{\ensuremath{\tau}}
\newunicodechar{φ}{\ensuremath{\varphi}}
\newunicodechar{ψ}{\ensuremath{\psi}}
\newunicodechar{ω}{\ensuremath{\omega}}
\newunicodechar{Σ}{\ensuremath{\Sigma}}
% majuscules produites par \MakeUppercase dans les titres (α-aminés -> Α)
\newunicodechar{Α}{\ensuremath{\alpha}}
\newunicodechar{Β}{\ensuremath{\beta}}
\newunicodechar{Γ}{\ensuremath{\Gamma}}
\newunicodechar{Δ}{\ensuremath{\Delta}}
\newunicodechar{Ε}{\ensuremath{\varepsilon}}
\newunicodechar{Θ}{\ensuremath{\Theta}}
\newunicodechar{Λ}{\ensuremath{\Lambda}}
\newunicodechar{Μ}{\ensuremath{\mu}}
\newunicodechar{Τ}{\ensuremath{\tau}}
\newunicodechar{Φ}{\ensuremath{\Phi}}
\newunicodechar{Ψ}{\ensuremath{\Psi}}
\newunicodechar{Ω}{\ensuremath{\Omega}}
\newunicodechar{↦}{\ensuremath{\mapsto}}
\newunicodechar{∈}{\ensuremath{\in}}
\newunicodechar{∉}{\ensuremath{\notin}}
\newunicodechar{∧}{\ensuremath{\wedge}}
\newunicodechar{⇔}{\ensuremath{\Leftrightarrow}}
\newunicodechar{⟺}{\ensuremath{\Longleftrightarrow}}
\newunicodechar{⇒}{\ensuremath{\Rightarrow}}
\newunicodechar{⟹}{\ensuremath{\Longrightarrow}}
\newunicodechar{∘}{\ensuremath{\circ}}
\newunicodechar{∪}{\ensuremath{\cup}}
\newunicodechar{∩}{\ensuremath{\cap}}
\newunicodechar{≡}{\ensuremath{\equiv}}
\newunicodechar{⊂}{\ensuremath{\subset}}
\newunicodechar{⊥}{\ensuremath{\perp}}
\newunicodechar{∃}{\ensuremath{\exists}}
\newunicodechar{∀}{\ensuremath{\forall}}
\newunicodechar{∛}{\ensuremath{\sqrt[3]{\,}}}
\newunicodechar{∜}{\ensuremath{\sqrt[4]{\,}}}
\newunicodechar{⃗}{\ensuremath{{}^{\to}}}
\newunicodechar{ⁿ}{\ifmmode^{n}\else\textsuperscript{\ensuremath{n}}\fi}
\newunicodechar{ˣ}{\ifmmode^{x}\else\textsuperscript{\ensuremath{x}}\fi}
\newunicodechar{ᵇ}{\ifmmode^{b}\else\textsuperscript{\ensuremath{b}}\fi}
\newunicodechar{ʳ}{\ifmmode^{r}\else\textsuperscript{\ensuremath{r}}\fi}
\newunicodechar{ᵘ}{\ifmmode^{u}\else\textsuperscript{\ensuremath{u}}\fi}
\newunicodechar{ʸ}{\ifmmode^{y}\else\textsuperscript{\ensuremath{y}}\fi}
\newunicodechar{ᵃ}{\ifmmode^{a}\else\textsuperscript{\ensuremath{a}}\fi}
\newunicodechar{ᵉ}{\ifmmode^{e}\else\textsuperscript{\ensuremath{e}}\fi}
\newunicodechar{ᵏ}{\ifmmode^{k}\else\textsuperscript{\ensuremath{k}}\fi}
\newunicodechar{ᵖ}{\ifmmode^{p}\else\textsuperscript{\ensuremath{p}}\fi}
\newunicodechar{ᴺ}{\ifmmode^{N}\else\textsuperscript{\ensuremath{N}}\fi}
\newunicodechar{ᶜ}{\ifmmode^{c}\else\textsuperscript{\ensuremath{c}}\fi}
\newunicodechar{ʰ}{\ifmmode^{h}\else\textsuperscript{\ensuremath{h}}\fi}
\newunicodechar{ᵠ}{\ifmmode^{q}\else\textsuperscript{\ensuremath{q}}\fi}
\newunicodechar{ₙ}{\ifmmode_{n}\else\textsubscript{\ensuremath{n}}\fi}
\newunicodechar{ₐ}{\ifmmode_{a}\else\textsubscript{\ensuremath{a}}\fi}
\newunicodechar{ᵢ}{\ifmmode_{i}\else\textsubscript{\ensuremath{i}}\fi}
\newunicodechar{ᵣ}{\ifmmode_{r}\else\textsubscript{\ensuremath{r}}\fi}
\newunicodechar{ₖ}{\ifmmode_{k}\else\textsubscript{\ensuremath{k}}\fi}
\newunicodechar{ᵦ}{\ifmmode_{\beta}\else\textsubscript{\ensuremath{\beta}}\fi}
\newunicodechar{ₓ}{\ifmmode_{x}\else\textsubscript{\ensuremath{x}}\fi}
\newfontfamily\popdisplay{ArchivoBlack-Regular.ttf}[Path=../../../fonts/]
\newfontfamily\poplogo{SpaceGrotesk-Bold.ttf}[Path=../../../fonts/]
\newfontfamily\popsemi{PlusJakartaSans-SemiBold.ttf}[Path=../../../fonts/]
\newfontfamily\popxbold{PlusJakartaSans-ExtraBold.ttf}[Path=../../../fonts/]
\newfontfamily\popxboldls{PlusJakartaSans-ExtraBold.ttf}[Path=../../../fonts/,LetterSpace=3.0]

\setlist[enumerate]{leftmargin=1.7em,itemsep=5pt,topsep=4pt}
\setlist[itemize]{leftmargin=1.7em,itemsep=4pt,topsep=4pt,label={\color{poporange}\textbullet}}
\setlength{\parindent}{0pt}
\setlength{\parskip}{5pt}
\renewcommand{\arraystretch}{1.25}

% pgfplots : style Pop par défaut
\pgfplotsset{
  every axis/.append style={axis line style={popink,line width=1pt},tick style={popink},
    grid style={popline,line width=0.5pt},label style={font=\small},tick label style={font=\footnotesize}},
  popcurve/.style={poporange,line width=1.8pt,no markers,smooth},
}
% Même style disponible dans un \draw TikZ simple (hors pgfplots)
\tikzset{popcurve/.style={poporange,line width=1.8pt}}
\tikzset{popgrid/.style={popline,very thin}}
\colorlet{popcurve}{poporange} % le nom du style est parfois utilisé comme couleur

% ============================================================
% Pill / squiggle
% ============================================================
\newcommand{\poppill}[3]{%
  \tikz[baseline=-0.55ex]\node[draw=popink,line width=1.2pt,rounded corners=2.6pt,%
    fill=#2,inner xsep=6.5pt,inner ysep=3pt]{%
    \popxboldls\fontsize{7}{7}\selectfont\color{#3}\MakeUppercase{#1}};%
}
\newcommand{\popsquiggle}[1]{%
  \tikz[baseline=-0.4ex]{
    \draw[#1,line width=2.6pt,line cap=round]
      plot [smooth] coordinates {(0,0.09) (0.34,0.01) (0.68,0.21) (1.02,0.01) (1.36,0.10)};
  }%
}
% Mot-clé surligné (marqueur orange sous le texte)
\newcommand{\cle}[1]{{\popxbold\color{popdark}#1}}

% ============================================================
% En-tête / pied
% ============================================================
\pagestyle{fancy}
\fancyhf{}
\renewcommand{\headrulewidth}{0pt}
\renewcommand{\footrulewidth}{0pt}
\lhead{\raisebox{-0.15ex}{\poplogo\fontsize{13}{13}\selectfont\color{popink}OnBuch\color{poporange}+}}
\rhead{\poppill{\DOCMATIERE\ \textbullet\ \DOCNIVEAU\ \textbullet\ Leçon \DOCLECON}{white}{popink}}
\lfoot{\popxbold\fontsize{7.6}{7.6}\selectfont\color{popmuted}\MakeUppercase{Cours signé OnBuch+}\ \textbullet\ {\popsemi\DOCTITRE}}
\rfoot{\tikz[baseline=-0.6ex]\node[rounded corners=8.5pt,fill=popink,minimum height=17pt,minimum width=17pt,inner xsep=5pt,inner sep=0]{\popxbold\fontsize{8}{8}\selectfont\color{popcream}\thepage\,/\,\pageref*{LastPage}};}

% ============================================================
% Titres
% ============================================================
\newcommand{\chaptitre}[2]{%
  \begin{tcolorbox}[enhanced,colback=poporange,colframe=popink,boxrule=2.6pt,arc=11pt,
      drop shadow=popink,left=15pt,right=15pt,top=12pt,bottom=11pt,before skip=3pt,after skip=18pt]
    \poppill{#2}{popink}{popcream}\par\vspace{9pt}
    {\raggedright\hyphenpenalty=10000\exhyphenpenalty=10000\popdisplay\fontsize{22}{24}\selectfont\color{popcream}\MakeUppercase{#1}\par}
  \end{tcolorbox}
}
% Section numérotée : pastille orange avec le numéro + titre Archivo
\newcounter{popsec}
\newcommand{\coursec}[1]{%
  \stepcounter{popsec}\setcounter{popsub}{0}%
  \par\vspace{16pt}\noindent
  \begin{minipage}{\linewidth}
  \tikz[baseline=-0.7ex]\node[draw=popink,line width=1.6pt,fill=poporange,rounded corners=4pt,minimum size=22pt,inner sep=0]{\popdisplay\fontsize{12}{12}\selectfont\color{popcream}\thepopsec};%
  \hspace{8pt}{\popdisplay\fontsize{15}{17}\selectfont\color{popink}\MakeUppercase{#1}}\par
  \vspace{4pt}\noindent{\color{poporange}\rule{36pt}{3.6pt}}
  \end{minipage}\par\nopagebreak\vspace{8pt}
}
\newcounter{popsub}
\newcommand{\courssub}[1]{%
  \stepcounter{popsub}%
  \par\vspace{9pt}\noindent{\popxbold\fontsize{11.5}{13}\selectfont\color{popink}{\color{poporange}\thepopsec.\thepopsub}\hspace{6pt}#1}\par\nopagebreak\vspace{4pt}
}

% ============================================================
% Boîtes pédagogiques — même recette : fond blanc, bordure épaisse colorée,
% ombre dure encre, badge plein. Titre optionnel [#1].
% ============================================================
\tcbset{popbase/.style={breakable,enhanced,colback=white,boxrule=2.3pt,arc=9pt,
  shadow={3pt}{-3pt}{0pt}{popink},left=14pt,right=14pt,top=11pt,bottom=11pt,before skip=11pt,after skip=15pt,
  fontupper=\popsemi\fontsize{10.6}{14.5}\selectfont\color{popink},
  fontlower=\popsemi\fontsize{10.6}{14.5}\selectfont\color{popink}}}
% Coche dessinée (le glyphe ✓ n'existe pas dans les polices du document)
\newcommand{\popcheck}[1]{\tikz[baseline=-0.5ex]\draw[#1,line width=1.6pt,line cap=round,line join=round](-0.13,0.0)--(-0.04,-0.09)--(0.13,0.1);}
\providecommand{\coche}{\popcheck{popgreen}}
\providecommand{\resultat}[1]{\par\smallskip\noindent\fcolorbox{popgreen}{popgreenL}{\parbox{\dimexpr\linewidth-2\fboxsep-2\fboxrule\relax}{\textbf{Résultat.} #1}}\par\smallskip}
\newcommand{\popboxhead}[3]{\poppill{#1}{#2}{white}\ifx\relax#3\relax\else\hspace{6pt}{\popxbold\fontsize{10.6}{12}\selectfont\color{popink}#3}\fi\par\vspace{7pt}}

\newtcolorbox{aretenir}[1][]{popbase,colframe=popgreen,before upper={\popboxhead{À retenir}{popgreen}{#1}}}
\newtcolorbox{exemplebox}[1][]{popbase,colframe=poporange,before upper={\popboxhead{Exemple}{poporange}{#1}}}
\newtcolorbox{definition}[1][]{popbase,colframe=popblue,before upper={\popboxhead{Définition}{popblue}{#1}}}
\newtcolorbox{propriete}[1][]{popbase,colframe=poppurple,before upper={\popboxhead{Propriété}{poppurple}{#1}}}
\newtcolorbox{methode}[1][]{popbase,colframe=popgold,before upper={\popboxhead{Méthode}{popgold}{#1}}}
\newtcolorbox{attention}[1][]{popbase,colframe=poppink,before upper={\popboxhead{Attention}{poppink}{#1}}}
\newtcolorbox{experience}[1][]{popbase,colframe=popblue,colback=popblueL,before upper={\popboxhead{Expérience}{popblue}{#1}}}
% « Le savais-tu ? » : carte violette pleine (contexte, culture, Cameroun)
\newtcolorbox{savaistu}[1][]{popbase,colback=poppurple,colframe=popink,
  fontupper=\popsemi\fontsize{10.4}{14.2}\selectfont\color{white},
  before upper={\poppill{Le savais-tu\,?}{white}{poppurple}\ifx\relax#1\relax\else\hspace{6pt}{\popxbold\color{white}#1}\fi\par\vspace{7pt}}}
% Exercice résolu : énoncé en haut, \tcblower puis la solution sur fond crème
\newcounter{popexo}\newcounter{popexr}
\newtcolorbox{exoresolu}[1][]{popbase,colframe=poporange,colbacklower=popcream,
  segmentation style={popink,line width=1.2pt,dashed},
  before upper={\stepcounter{popexr}\popboxhead{Exercice résolu \thepopexr}{poporange}{#1}},
  before lower={\poppill{Solution}{popink}{popcream}\par\vspace{6pt}}}
% Exercice (non résolu en place) — la correction est dans la partie Corrigés
\newtcolorbox{exercice}[1][]{popbase,colframe=popink,
  before upper={\stepcounter{popexo}\popboxhead{Exercice \thepopexo}{popink}{#1}}}
% Corrigé
\newtcolorbox{corrige}[1][]{popbase,colframe=popgreen,colback=popgreenL,
  before upper={\popboxhead{Corrigé}{popgreen}{#1}}}
% Objectifs / prérequis / bilan
\newtcolorbox{objectifs}{popbase,colframe=popink,colback=poporangeL,
  before upper={\popboxhead{Objectifs de la leçon}{popink}{}}}
\newtcolorbox{prerequis}{popbase,colframe=popink,colback=popblueL,
  before upper={\popboxhead{Prérequis}{popblue}{}}}
\newtcolorbox{bilan}{popbase,colframe=popink,colback=white,boxrule=2.8pt,
  before upper={\popboxhead{Fiche bilan}{poporange}{L'essentiel à connaître pour l'examen}}}
% Figure encadrée avec légende
\newtcolorbox{popfigure}[1][]{enhanced,breakable=false,colback=white,colframe=popline,boxrule=1.4pt,arc=8pt,
  left=10pt,right=10pt,top=10pt,bottom=8pt,before skip=10pt,after skip=12pt,halign=center,
  lower separated=false,halign lower=center,
  fontlower=\popsemi\fontsize{9}{11.5}\selectfont\color{popmuted},
  #1}
\newcommand{\legende}[1]{\par\vspace{6pt}{\popsemi\fontsize{9}{11.5}\selectfont\color{popmuted}#1}}

% ============================================================
% Signature OnBuch+
% ============================================================
\newcommand{\poplogoinline}[1]{{\poplogo\fontsize{#1}{#1}\selectfont\color{popink}OnBuch\color{poporange}+}}
\newcommand{\signatureonbuch}{%
  \begin{tcolorbox}[enhanced,colback=popink,colframe=popink,boxrule=2.2pt,arc=10pt,
      drop shadow=poporange,left=14pt,right=14pt,top=10pt,bottom=10pt,before skip=0pt,after skip=0pt]
    \tikz[baseline=-0.7ex]\node[circle,fill=popgreen,draw=popcream,line width=1.3pt,minimum size=22pt,inner sep=0]{\popcheck{white}};%
    \hspace{9pt}\begin{minipage}[c]{0.62\linewidth}
      {\popxbold\fontsize{12}{13}\selectfont\color{popcream}Signé\ \textbullet\ {\poplogo OnBuch\color{poporange}+}}\par\vspace{2pt}
      {\raggedright\popsemi\fontsize{8}{10.5}\selectfont\color{popline}\MakeUppercase{Équipe pédagogique OnBuch+}\\\MakeUppercase{Conforme au programme officiel MINESEC}\par}
    \end{minipage}\hfill
    {\poplogo\fontsize{22}{22}\selectfont\color{popcream}OnBuch\color{poporange}+}
  \end{tcolorbox}%
}

% ============================================================
% Page de garde
% #1 titre · #2 matière · #3 niveau · #4 module · #5 n° leçon · #6 durée ·
% #7 description / objectifs (texte ou itemize)
% ============================================================
\newcommand{\pagegarde}[7]{%
  \thispagestyle{empty}
  \newgeometry{a4paper,margin=1.7cm,top=1.6cm,bottom=1.4cm}%
  \begin{tikzpicture}[remember picture,overlay]
    \fill[poporange!18] ([xshift=-3.2cm,yshift=-3.6cm]current page.north east) circle (4.6cm);
    \fill[poppurple!14] ([xshift=2.4cm,yshift=7.6cm]current page.south west) circle (3.8cm);
  \end{tikzpicture}%
  {\poplogo\fontsize{30}{30}\selectfont\color{popink}OnBuch\color{poporange}+}\hfill
  \raisebox{0.6ex}{\poppill{Cours signé \textbullet\ Premium}{popink}{popcream}}\par
  \vspace{1pt}\popsquiggle{poppurple}\par
  \vspace{8pt}
  \begin{tcolorbox}[enhanced,colback=poporange,colframe=popink,boxrule=2.8pt,arc=13pt,
      drop shadow=popink,left=18pt,right=18pt,top=12pt,bottom=13pt,after skip=10pt]
    \poppill{#2 \textbullet\ #3}{popink}{popcream}\hspace{5pt}\poppill{#4}{white}{popink}\par\vspace{8pt}
    {\popdisplay\fontsize{14}{14}\selectfont\color{popink}LEÇON #5}\par\vspace{4pt}
    {\raggedright\hyphenpenalty=10000\exhyphenpenalty=10000\popdisplay\fontsize{25}{27}\selectfont\color{popcream}\MakeUppercase{#1}\par}
    \vspace{8pt}
    {\popxbold\fontsize{9.5}{9.5}\selectfont\color{popink}\MakeUppercase{Durée conseillée : #6}}
  \end{tcolorbox}
  \begin{tcolorbox}[enhanced,colback=white,colframe=popink,boxrule=2.2pt,arc=10pt,
      drop shadow=popink,left=16pt,right=16pt,top=10pt,bottom=10pt,after skip=8pt]
    \poppill{Dans cette leçon}{poppurple}{white}\par\vspace{5pt}
    \popsemi\fontsize{9.3}{12.6}\selectfont\color{popink}#7
  \end{tcolorbox}
  \noindent\begin{minipage}[t]{0.487\linewidth}
    \begin{tcolorbox}[enhanced,colback=popgreen,colframe=popink,boxrule=2.2pt,arc=10pt,
        drop shadow=popink,left=11pt,right=11pt,top=10pt,bottom=10pt,height=3.1cm,valign=center]
      \tcbox[colback=white,colframe=popink,boxrule=1.3pt,arc=5pt,boxsep=0pt,nobeforeafter,
        left=3pt,right=3pt,top=3pt,bottom=3pt,valign=center]{\qrcode[height=1.9cm]{https://play.google.com/store/apps/details?id=com.luvvix.onbuch&hl=fr}}%
      \hspace{8pt}\begin{minipage}[c]{0.5\linewidth}
        {\popdisplay\fontsize{11}{12.5}\selectfont\color{white}\MakeUppercase{Télécharge OnBuch}}\par\vspace{4pt}
        {\raggedright\popsemi\fontsize{8.4}{11}\selectfont\color{white}Gratuit \textbullet\ Google Play\\Tuteur IA, annales, concours, résultats\par}
      \end{minipage}
    \end{tcolorbox}
  \end{minipage}\hfill
  \begin{minipage}[t]{0.487\linewidth}
    \begin{tcolorbox}[enhanced,colback=poppurple,colframe=popink,boxrule=2.2pt,arc=10pt,
        drop shadow=popink,left=11pt,right=11pt,top=10pt,bottom=10pt,height=3.1cm,valign=center]
      \tikz[baseline=-0.7ex]\node[circle,fill=white,draw=popink,line width=1.3pt,minimum size=1.6cm,inner sep=0]{\popdisplay\fontsize{24}{24}\selectfont\color{poppurple}+};%
      \hspace{8pt}\begin{minipage}[c]{0.6\linewidth}
        {\popdisplay\fontsize{11}{12.5}\selectfont\color{white}\MakeUppercase{Passe à OnBuch+}}\par\vspace{4pt}
        {\raggedright\popsemi\fontsize{8.4}{11}\selectfont\color{white}Tous les cours signés, Tuteur IA illimité, fiches PDF — onglet OnBuch+ de l'app\par}
      \end{minipage}
    \end{tcolorbox}
  \end{minipage}
  \vspace{8pt}
  \popcontenu
  \vfill
  \signatureonbuch
  \restoregeometry
  \newpage
}

% Rangée « Ce cours contient » (page de garde)
\newcommand{\poptile}[3]{% #1 couleur #2 gros texte #3 légende
  \begin{tcolorbox}[enhanced,colback=#1,colframe=popink,boxrule=2pt,arc=9pt,shadow={2.5pt}{-2.5pt}{0pt}{popink},
      left=6pt,right=6pt,top=9pt,bottom=9pt,halign=center,valign=center,height=1.75cm,width=\linewidth,nobeforeafter]
    {\popdisplay\fontsize{12.5}{13}\selectfont\color{white}\MakeUppercase{#2}}\par\vspace{3pt}
    {\centering\popsemi\fontsize{7.8}{9.5}\selectfont\color{white}#3\par}
  \end{tcolorbox}}
\newcommand{\popcontenu}{%
  \par\noindent{\popxboldls\fontsize{7.5}{7.5}\selectfont\color{popmuted}CE COURS SIGNÉ CONTIENT}\par\vspace{7pt}
  \noindent\begin{minipage}[t]{0.235\linewidth}\poptile{popblue}{Cours}{complet, illustré, pas à pas}\end{minipage}\hfill
  \begin{minipage}[t]{0.235\linewidth}\poptile{poporange}{Méthodes}{et exercices résolus}\end{minipage}\hfill
  \begin{minipage}[t]{0.235\linewidth}\poptile{poppink}{Exercices}{type examen + corrigés}\end{minipage}\hfill
  \begin{minipage}[t]{0.235\linewidth}\poptile{popgreen}{Fiche bilan}{l'essentiel à retenir}\end{minipage}\par}

% Sommaire
\newtcolorbox{sommaire}{enhanced,colback=white,colframe=popink,boxrule=2.2pt,arc=10pt,
  drop shadow=popink,left=18pt,right=18pt,top=13pt,bottom=13pt,before skip=6pt,after skip=20pt,
  before upper={\poppill{Sommaire}{poppurple}{white}\par\vspace{9pt}\popsemi\fontsize{10.6}{14.5}\selectfont\color{popink}}}

% Signature de fin de cours — poussée tout en bas de la dernière page
\newcommand{\findecours}{%
  \par\vfill\nopagebreak
  \begin{center}\popsquiggle{poporange}\end{center}\vspace{4pt}
  \signatureonbuch
}
\AtBeginDocument{\def\llbracket{\mathopen{[\mkern-3.5mu[}}\def\rrbracket{\mathclose{]\mkern-3.5mu]}}} % intervalles d'entiers (glyphe ⟦ absent de la police)
% Glyphes absents de Fira Math : redéfinis à partir de glyphes disponibles
\AtBeginDocument{%
  \def\setminus{\mathbin{\backslash}}\let\smallsetminus\setminus
  \def\vdots{\mathord{\vbox{\baselineskip3.2pt\lineskiplimit0pt\kern4pt\hbox{.}\hbox{.}\hbox{.}}}}%
  \def\ddots{\mathinner{\mkern1mu\raise6.5pt\hbox{.}\mkern2mu\raise3.6pt\hbox{.}\mkern2mu\raise0.7pt\hbox{.}\mkern1mu}}%
  \def\triangle{\mathord{\scalebox{0.9}{$\symup{\Delta}$}}}%
  \def\nabla{\mathord{\raisebox{\depth}{\scalebox{1}[-1]{$\symup{\Delta}$}}}}%
  \def\checkmark{\popcheck{popdark}}%
  \def\star{\mathbin{\ast}}\def\bigstar{\mathord{\ast}}%
  \def\diamond{\mathbin{\scalebox{0.6}{$\Diamond$}}}%
  \def\bigcup{\mathop{\mathchoice{\raisebox{-0.3em}{\scalebox{1.7}{$\cup$}}}{\scalebox{1.25}{$\cup$}}{\cup}{\cup}}\displaylimits}%
  \def\bigcap{\mathop{\mathchoice{\raisebox{-0.3em}{\scalebox{1.7}{$\cap$}}}{\scalebox{1.25}{$\cap$}}{\cap}{\cap}}\displaylimits}%
  \def\lesseqqgtr{\mathrel{\lessgtr}}\def\bigoplus{\mathop{\oplus}\displaylimits}%
}
\providecommand{\ssi}{si et seulement si} % abréviation fréquente
\AtBeginDocument{\providecommand{\wideparen}{\overparen}} % arc de cercle

%% FIGFIX-BEGIN v3 — correctifs globaux des figures (docs/onbuch-generation/tools/figfix)
\usepackage{adjustbox}\usepackage{etoolbox}
% toute figure TikZ/pgfplots plus large que la ligne est réduite à la largeur disponible
\BeforeBeginEnvironment{tikzpicture}{\begin{adjustbox}{max width=\linewidth}}
\AfterEndEnvironment{tikzpicture}{\end{adjustbox}}
% légendes pgfplots lisibles par-dessus les courbes
\pgfplotsset{every axis/.append style={legend style={fill=white,fill opacity=0.92,text opacity=1,draw=black!15,inner sep=3pt}}}
% étiquettes d'arêtes (syntaxe edge["..."]) avec fond blanc
\tikzset{every edge quotes/.append style={fill=white,inner sep=1.2pt}}
% bulles de cartes mentales : texte à la ligne dans la bulle, bulle un peu plus grande
\tikzset{every concept/.append style={text width=2.0cm,align=center,minimum size=2.3cm,inner sep=2pt}}
%% FIGFIX-END
\begin{document}

\pagegarde{Séismes}{SVTEEHB}{3ème}{SVTEEHB – classe de 3ème}{N12}{2 séances (fiche de progression harmonisée)}{Cette leçon étudie les séismes : leurs manifestations observables, leur origine liée aux mouvements de l'écorce terrestre et les méthodes permettant d'évaluer leur intensité. Elle situe ensuite les séismes à l'échelle mondiale et aborde les moyens de prévision et de prévention.
\vspace{4pt}
\begin{multicols}{2}\small
\begin{itemize}
\item À la fin de cette leçon, je sais définir un séisme et citer ses manifestations observables.
\item À la fin de cette leçon, je sais distinguer le foyer, l'épicentre et les ondes sismiques.
\item À la fin de cette leçon, je sais expliquer l'origine des séismes à partir de l'accumulation et de la libération des contraintes tectoniques.
\item À la fin de cette leçon, je sais distinguer l'intensité (échelle de Mercalli) de la magnitude (échelle de Richter) et interpréter un sismogramme.
\item À la fin de cette leçon, je sais exploiter une carte mondiale de répartition des séismes pour la relier aux limites des plaques lithosphériques.
\item À la fin de cette leçon, je sais citer les signes précurseurs et les limites de la prévision des séismes.
\end{itemize}\end{multicols}}

\begin{sommaire}
\begin{multicols}{2}
\begin{enumerate}
\item Manifestations des séismes
\item Origine des séismes : foyer, épicentre et ondes
\item Évaluation de l'intensité et de la magnitude
\item Répartition mondiale des séismes
\item Prévision et prévention des séismes
\item Activité d'intégration
\item Méthodes et exercices résolus
\item Exercices
\item Corrigés des exercices
\item Fiche bilan
\end{enumerate}
\end{multicols}
\end{sommaire}

\begin{objectifs}
\begin{itemize}
\item À la fin de cette leçon, je sais définir un séisme et citer ses manifestations observables.
\item À la fin de cette leçon, je sais distinguer le foyer, l'épicentre et les ondes sismiques.
\item À la fin de cette leçon, je sais expliquer l'origine des séismes à partir de l'accumulation et de la libération des contraintes tectoniques.
\item À la fin de cette leçon, je sais distinguer l'intensité (échelle de Mercalli) de la magnitude (échelle de Richter) et interpréter un sismogramme.
\item À la fin de cette leçon, je sais exploiter une carte mondiale de répartition des séismes pour la relier aux limites des plaques lithosphériques.
\item À la fin de cette leçon, je sais citer les signes précurseurs et les limites de la prévision des séismes.
\item À la fin de cette leçon, je sais proposer des mesures de prévention et des comportements adaptés face à un séisme.
\item À la fin de cette leçon, je sais rédiger et justifier une conclusion à partir de documents (sismogrammes, cartes, tableaux de données).
\end{itemize}
\end{objectifs}

\begin{prerequis}
\begin{itemize}
\item Structure interne de la Terre (croûte, manteau, noyau) vue en classe de 4ème
\item Notion de lithosphère et de plaques lithosphériques (début de l'année de 3ème)
\item Notion de faille et de déformation des roches
\item Lecture et interprétation de graphiques et de cartes
\item Les volcans et la ceinture de feu du Pacifique (leçon précédente)
\end{itemize}
\end{prerequis}

\newpage

\coursec{Manifestations des séismes}

Le 28 mars 1999, aux alentours de 3 heures du matin, la région de Monatélé, à une cinquantaine de kilomètres au nord-est de Yaoundé, est brutalement arrachée au sommeil. Le sol ondule, les murs craquent, la vaisselle s'écrase au sol, la panique saisit des milliers d'habitants qui se ruent dehors. Les jours suivants, on compte des maisons fissurées, des cheminées effondrées, une école endommagée. Pourtant, le Cameroun ne figure pas sur la carte des pays à haut risque sismique comme le Japon ou la Turquie. D'où vient cette violence soudaine ? Pourquoi la Terre tremble-t-elle par à-coups ? Comment mesurer la force de ces secousses et, surtout, peut-on les prévoir pour mieux s'en protéger ?

\courssub{Définition et nature du phénomène}

\begin{definition}[Séisme (ou tremblement de terre)]
Un \cle{séisme} est une secousse brusque et passagère de la surface de la Terre, due à une libération soudaine d'énergie accumulée en profondeur dans les roches. Cette énergie se propage sous forme d'\cle{ondes sismiques} qui font vibrer le sol.
\end{definition}

Contrairement aux vibrations continues produites par le passage d'un train lourd ou par une machine d'usine, le séisme est un événement \emph{transitoire} : il dure de quelques secondes à quelques minutes (rarement plus), mais ses effets peuvent être dévastateurs. Le point de départ de la rupture, situé à l'intérieur du Globe, s'appelle le \cle{foyer} (ou hypocentre) ; sa projection verticale à la surface est l'\cle{épicentre}. C'est au voisinage de l'épicentre que les secousses sont les plus violentes.

\begin{experience}[Observation de la propagation d'une onde sur un ressort]
\textbf{Matériel :} Un long ressort (type Slinky) tendu sur une table lisse entre deux élèves.
\textbf{Protocole :} L'un des élèves donne une impulsion brusque et transversale (ou longitudinale) à son extrémité du ressort.
\textbf{Observation :} Une onde se propage le long du ressort ; les spires oscillent autour de leur position d'équilibre mais ne se déplacent pas durablement avec l'onde. L'énergie se déplace, la matière vibre seulement.
\textbf{Interprétation :} Le sol se comporte comme un milieu élastique. La rupture brutale au foyer libère une énergie élastique qui se propage sous forme d'ondes (ondes P et ondes S, détaillées en section 2). La vibration du sol en surface est la manifestation directe du passage de ces ondes.
\end{experience}

\courssub{Manifestations directes sur le milieu naturel}

Lorsque les ondes sismiques atteignent la surface, elles déforment le sol. L'amplitude des déplacements du sol (quelques micromètres à plusieurs mètres) et leur fréquence déterminent l'intensité des dégâts.

\begin{itemize}
    \item \textbf{Secousses et vibrations du sol :} Le sol oscille dans les trois directions de l'espace (verticale et deux horizontales). C'est la manifestation première, ressentie par tous.
    \item \textbf{Fissures et fractures du sol :} La déformation dépasse la limite d'élasticité des roches et des sols meubles : le sol se craquelle, des failles de surface apparaissent, parfois avec un décalage vertical (rejet) ou horizontal (coulissement).
    \item \textbf{Glissements de terrain et éboulements :} Sur les versants instables, la secousse réduit la cohésion des matériaux (effet de \emph{liquéfaction} pour les sols sableux saturés en eau, ou simple dépassement de l'angle de repos). Des pans entiers de colline s'effondrent.
    \item \textbf{Affaissements et liquéfaction :} Dans les plaines alluviales ou les zones portuaires gagnées sur la mer, les sols sableux saturés en eau perdent leur portance sous l'effet des vibrations cycliques : ils se transforment en « soupe » liquide. Les bâtiments s'enfoncent ou basculent.
\end{itemize}

\begin{popfigure}
\begin{tikzpicture}[scale=0.9, every node/.style={font=\small}]
    % Couches géologiques (avant faille - référence fantôme)
    \draw[popmuted, dashed, thick] (-4,0) -- (4,0); % Surface initiale
    \draw[popmuted, dashed] (-4,-1) -- (4,-1); % Couche 1
    \draw[popmuted, dashed] (-4,-2) -- (4,-2); % Couche 2
    \draw[popmuted, dashed] (-4,-3) -- (4,-3); % Couche 3
    
    % Faille normale active (décalage)
    % Bloc de gauche (relativement abaissé)
    \fill[popblueL!50] (-4,0) -- (-0.2,0) -- (-0.2,-1.2) -- (-4,-1.2) -- cycle; % Couche sup gauche
    \fill[popgreenL!50] (-4,-1.2) -- (-0.2,-1.2) -- (-0.2,-2.2) -- (-4,-2.2) -- cycle; % Couche mid gauche
    \fill[poporangeL!50] (-4,-2.2) -- (-0.2,-2.2) -- (-0.2,-3.5) -- (-4,-3.5) -- cycle; % Couche inf gauche
    
    % Bloc de droite (relativement relevé)
    \fill[popblueL!50] (0.2,0) -- (4,0) -- (4,-0.8) -- (0.2,-0.8) -- cycle; % Couche sup droite (plus haute)
    \fill[popgreenL!50] (0.2,-0.8) -- (4,-0.8) -- (4,-1.8) -- (0.2,-1.8) -- cycle; % Couche mid droite
    \fill[poporangeL!50] (0.2,-1.8) -- (4,-1.8) -- (4,-3.0) -- (0.2,-3.0) -- cycle; % Couche inf droite
    
    % Plan de faille
    \draw[popdark, very thick, decorate, decoration={zigzag, segment length=4pt, amplitude=1.5pt}] (-0.2,0) -- (0.2,-3.5);
    \node[popdark, rotate=-55] at (-0.5, -1.8) {\textbf{Plan de faille}};
    
    % Flèches de mouvement (relatif)
    \draw[poporange, thick, ->] (-1.5, -0.5) -- (-0.5, -1.0); % Gauche descend
    \draw[poporange, thick, ->] (1.5, -0.3) -- (0.5, -0.8); % Droite monte
    \node[poporange] at (-2.2, -0.3) {Bloc \emph{toit} (abaissé)};
    \node[poporange] at (2.5, -0.1) {Bloc \emph{mur} (relevé)};
    
    % Surface du sol : fissures
    \draw[popdark, very thick] (-4,0) -- (-0.2,0);
    \draw[popdark, very thick] (0.2,0) -- (4,0);
    % Fissures en surface (tension)
    \draw[popred, thick] (-1.5, 0) -- (-1.2, 0.4);
    \draw[popred, thick] (-0.8, 0) -- (-0.5, 0.5);
    \draw[popred, thick] (0.5, 0) -- (0.8, 0.4);
    \draw[popred, thick] (1.2, 0) -- (1.5, 0.5);
    \node[popred] at (0, 0.8) {\textbf{Fissures de tension}};
    
    % Maison endommagée (sur le bloc de gauche)
    \draw[popdark, thick] (-3.2, 0) -- (-3.2, 0.7) -- (-2.4, 0.7) -- (-2.4, 0); % Murs
    \draw[popdark, thick] (-3.2, 0.7) -- (-2.8, 1.1) -- (-2.4, 0.7); % Toit
    % Fissures sur maison
    \draw[popred, thick] (-3.0, 0.1) -- (-2.8, 0.6);
    \draw[popred, thick] (-2.7, 0.2) -- (-2.5, 0.8);
    \node[popred, anchor=east] at (-3.3, 0.4) {\textbf{Maison fissurée}};
    
    % Légendes couches
    \node[popdark] at (-3.5, -0.6) {Alluvions / Sols meubles};
    \node[popdark] at (-3.5, -1.7) {Roche tendre (grès, argilite)};
    \node[popdark] at (-3.5, -2.8) {Socle dur (granite, gneiss)};
    
    % Échelle
    \draw[|-|, thick] (-4, -3.8) -- (-3, -3.8);
    \node at (-3.5, -4.1) {1 km};
\end{tikzpicture}
\legende{Schéma en coupe d'une faille normale active (type faille de Montobé ou faille bordière du graben de la Sanaga). Le bloc toit (à gauche) s'abaisse par rapport au bloc mur (à droite). Le décalage des couches géologiques et les fissures de tension en surface sont typiques. Une maison construite au-dessus de la zone de rupture subit des dommages structurels majeurs.}
\end{popfigure}

\courssub{Manifestations sur les constructions et infrastructures}

La vulnérabilité des ouvrages dépend de leur conception (parasismique ou non), de la nature du sol (amplification sur sols meubles) et de la durée de la secousse.

\begin{itemize}
    \item \textbf{Fissuration des murs :} Apparition de fissures en escalier (suivant les joints de maçonnerie) ou diagonales (cisaillement) aux angles des ouvertures (portes, fenêtres), points faibles de la structure.
    \item \textbf{Effondrement des bâtiments :} Rupture des poteaux, flambage des murs de refend, pancake (empilement des planchers) pour les structures en béton armé mal ferraillées ; renversement des murs en terre crue ou en parpaings non chaînés.
    \item \textbf{Rupture des réseaux :} Cassure des canalisations d'eau, de gaz (risque d'incendie), de pétrole ; déformation des rails, fissuration du bitume, effondrement de ponts (appuis arrachés, piles cisailées).
\end{itemize}

\begin{popfigure}
\begin{tikzpicture}[scale=0.85, every node/.style={font=\small}]
    % Sol de référence
    \draw[popdark, very thick] (-5,0) -- (5,0);
    
    % Route (vue en coupe transversale : on voit l'épaisseur)
    % Partie gauche
    \fill[popgray] (-4.5, 0) rectangle (-0.5, 0.2);
    \draw[white, dashed] (-4.5, 0.1) -- (-0.5, 0.1); % Marquage central
    % Partie droite (décalée horizontalement vers la droite ET verticalement ?)
    % Simulons un décrochement dextre + rejet.
    \fill[popgray] (0.5, -0.2) rectangle (4.5, 0.0);
    \draw[white, dashed] (0.5, -0.1) -- (4.5, -0.1);
    
    % Faille traversant la route
    \draw[popred, very thick, decorate, decoration={zigzag, segment length=3pt, amplitude=1pt}] (-0.5, 0.2) -- (0.5, -0.2);
    \node[popred, rotate=-25] at (0, 0.3) {\textbf{Faille de surface}};
    \draw[poporange, <->, thick] (-0.5, 0.35) -- (0.5, 0.35);
    \node[poporange] at (0, 0.55) {Décalage horizontal $\approx$ 1,5 m};
    
    % Bâtiment effondré (sur le côté gauche, près faille)
    % État initial (fantôme)
    \draw[popmuted, dashed, thin] (-3.5, 0) rectangle (-2.0, 2.5);
    \draw[popmuted, dashed, thin] (-3.5, 2.5) -- (-2.75, 3.2) -- (-2.0, 2.5);
    % État effondré (pancake)
    \fill[poporangeL!70] (-3.5, 0) rectangle (-2.0, 0.4); % Plancher RDC
    \fill[poporangeL!70] (-3.5, 0.4) rectangle (-2.0, 0.8); % Plancher R+1
    \fill[poporangeL!70] (-3.5, 0.8) rectangle (-2.0, 1.2); % Plancher R+2 (toit)
    \draw[popdark, thick] (-3.5, 0) -- (-3.5, 1.2); % Mur gauche restant
    \draw[popdark, thick] (-2.0, 0) -- (-2.0, 1.2); % Mur droit restant
    \node[popred, align=center] at (-2.75, 1.6) {\textbf{Effondrement}\\\textbf{type « pancake »}};
    
    % Conduite rompue
    \draw[popblue, line width=3pt] (-4.5, -1.0) -- (-0.7, -1.0);
    \draw[popblue, line width=3pt] (0.7, -1.3) -- (4.5, -1.3);
    \draw[popred, thick, fill=popred] (-0.7, -1.0) circle (3pt);
    \draw[popred, thick, fill=popred] (0.7, -1.3) circle (3pt);
    \node[popblue] at (0, -1.6) {\textbf{Conduite rompue / Fuite}};
    
    % Poteau électrique penché
    \draw[popdark, thick] (3.0, 0) -- (3.0, 2.0);
    \draw[popdark, thick] (3.0, 2.0) -- (3.5, 2.0);
    \draw[popdark, thick] (3.0, 1.5) -- (3.5, 1.5);
    \draw[popdark, thick, rotate=-15] (3.0, 0) -- (3.0, 2.0); % Poteau penché
    \node[popmuted] at (3.8, 1.0) {Poteau penché};
    
    % Légende
    \node[anchor=north west, align=left, fill=popcream, draw=popline, rounded corners, inner sep=4pt] at (-4.8, 3.0) {
        \textbf{Légende :}\\
        \textcolor{popgray}{\rule{4pt}{4pt}} Route bitumée (décalée)\\
        \textcolor{poporangeL}{\rule{4pt}{4pt}} Débris bâtiment (béton/parpaing)\\
        \textcolor{popblue}{\rule{4pt}{4pt}} Canalisation eau/gaz\\
        \textcolor{popred}{\rule{4pt}{4pt}} Faille / Rupture nette
    };
\end{tikzpicture}
\legende{Vue schématique (coupe transversale simplifiée) des dégâts d'un séisme de type décrochement (faille transformante) traversant une zone urbaine. On observe le décalage horizontal de la route, l'effondrement par empilement des planchers (« pancake ») d'un bâtiment non parasismique, la rupture des réseaux enterrés et l'instabilité des ouvrages d'art.}
\end{popfigure}

\begin{exemplebox}[Le séisme d'Haïti, 12 janvier 2010]
Le séisme de magnitude 7,0 (Mw), dont l'épicentre était à seulement 25 km de Port-au-Prince (profond de 13 km), a provoqué une catastrophe humanitaire majeure : \textbf{plus de 200 000 morts}, 300 000 blessés et 1,5 million de sans-abri. L'ampleur du drame s'explique par la convergence de trois facteurs :
\begin{enumerate}
    \item \textbf{Proximité de l'épicentre} et faible profondeur $\rightarrow$ fort secouissement (intensité IX-X EMS).
    \item \textbf{Amplification de site} : la capitale est construite sur des sédiments meubles (plaine du Cul-de-Sac) qui amplifient les ondes.
    \item \textbf{Extrême vulnérabilité du bâti} : constructions en parpaings de ciment non armés, sans chaînages horizontaux ni verticaux, dalles lourdes sur murs fragiles. Aucune norme parasismique n'était appliquée.
\end{enumerate}
Cet exemple illustre tragiquement qu'un séisme de magnitude modérée peut être bien plus meurtrier qu'un séisme plus fort (ex: Chili 2010, Mw 8,8, $\sim$ 500 morts) s'il frappe une zone dense et non préparée.
\end{exemplebox}

\courssub{Manifestations secondaires et particulières}

Au-delà des secousses primaires, les séismes déclenchent des phénomènes en cascade :

\begin{itemize}
    \item \textbf{Tsunamis (raz-de-marée) :} Un séisme \emph{sous-marin} (foyer sous l'océan) avec une composante verticale importante (faille inverse ou normale) soulève ou abaisse brutalement une énorme masse d'eau. Cette perturbation se propage à grande vitesse ($\sim$ 700 km/h en haute mer) sous forme d'une vague de très grande longueur d'onde (100--200 km) et de faible hauteur (quelques dizaines de cm). En approchant des côtes, la vague ralentit, sa longueur d'onde diminue et sa hauteur augmente dramatiquement (jusqu'à 30--40 m, ex: Tōhoku 2011), submergeant tout sur son passage.
    \item \textbf{Incendies :} Rupture des conduites de gaz, renversement de réchauds, courts-circuits électriques. Les incendies post-sismiques ont souvent causé plus de dégâts que la secousse elle-même (ex: San Francisco 1906, Tokyo 1923).
    \item \textbf{Modification du cours des rivières :} Soulèvement ou affaissement tectonique du lit, barrages naturels formés par des glissements de terrain (lacs de barrage), détournement du flux.
    \item \textbf{Éruptions volcaniques :} Un séisme peut déclencher l'éruption d'un volcan voisin en modifiant la contrainte dans la chambre magmatique ou en ouvrant des voies de montée pour le magma (ex: Mont Pinatubo 1991 précédé de séismes).
\end{itemize}

\courssub{Manifestations ressenties par les êtres vivants}

\begin{itemize}
    \item \textbf{Êtres humains :} Peur intense, vertiges, nausées (mal des transports dû au conflit visuel/vestibulaire), difficultés à se tenir debout (accélération du sol $> 0,1\,g$), troubles psychologiques post-traumatiques (syndrome de stress post-traumatique).
    \item \textbf{Animaux :} De nombreux témoignages (Chine, Japon, Italie, Haïti) rapportent des comportements anormaux \emph{avant} la secousse principale (chiens qui aboient sans raison, volées d'oiseaux, bétail qui refuse d'entrer dans les étables, serpents qui sortent d'hibernation). L'hypothèse scientifique principale est la détection des \emph{ondes P} (primaires, plus rapides mais moins destructrices) ou de signaux précurseurs (gaz radon, champs électromagnétiques, ultrasons) par une sensibilité sensorielle supérieure à la nôtre. \textbf{Attention :} ce n'est pas un moyen de prévision fiable et reproductible.
\end{itemize}

\courssub{Les répliques (aftershocks)}

\begin{definition}[Réplique]
Une \cle{réplique} est un séisme de magnitude inférieure à celle du séisme principal (\cle{choc principal}), qui se produit dans la même zone (même système de failles) dans les heures, jours, mois, voire années qui suivent.
\end{definition}

Elles résultent de la réorganisation des contraintes dans la croûte terrestre après la rupture principale. Leur fréquence et leur magnitude décroissent avec le temps. Elles sont dangereuses car elles font s'effondrer les bâtiments déjà fragilisés par le choc principal et perturbent les secours.

\begin{aretenir}[Séquence sismique typique]
\begin{itemize}
    \item \textbf{Précurseurs (parfois) :} microséismes, variations niveau nappe, gaz, comportement animal (non systématiques, non fiables).
    \item \textbf{Choc principal (Mainshock) :} plus grande magnitude, libération d'énergie maximale.
    \item \textbf{Répliques (Aftershocks) :} nombre élevé juste après, décroissance temporelle. La plus forte réplique est en moyenne de 1,2 unité de magnitude inférieure au choc principal.
\end{itemize}
\end{aretenir}

\courssub{Exemples de séismes marquants}

\begin{center}
\begin{tabularx}{\linewidth}{|l|X|c|c|X|}\hline
\rowcolor{popblueL}
\textbf{Séisme} & \textbf{Date} & \textbf{Magnitude (Mw)} & \textbf{Profondeur} & \textbf{Conséquences majeures} \\ \hline
Kobe (Japon) & 17 janv. 1995 & 6,9 & 16 km & 6 400 morts. Effondrement autoroutes, port, bâtiments anciens. Réveil sur normes parasismiques. \\ \hline
Haïti & 12 janv. 2010 & 7,0 & 13 km & \textbf{> 200 000 morts}. Catastrophe humanitaire. Vulnérabilité extrême du bâti. \\ \hline
Tōhoku (Japon) & 11 mars 2011 & 9,0 & 29 km & Tsunami géant (40 m). Accident nucléaire Fukushima. 18 000 morts/disparus. \\ \hline
Turquie--Syrie & 6 fév. 2023 & 7,8 / 7,5 & 18 / 10 km & \textbf{> 50 000 morts}. Doublet sismique. Effondrement massif immeubles récents non conformes. \\ \hline
Monatélé (Cameroun) & 28 mars 1999 & 4,5--5,0 & $\sim$ 10 km & Dégâts locaux (fissures, cheminées). Ressenti jusqu'à Yaoundé. Alerte sur risque modéré. \\ \hline
Mont Cameroun & Récurrents & $<$ 4,0 & Faible & Liés à l'activité volcanique (fracturation, migration magmatique). Surveillance OVSM. \\ \hline
\end{tabularx}
\end{center}

\begin{savaistu}[Sismicité au Cameroun : une activité modérée mais réelle]
Le Cameroun n'est pas une zone de forte sismicité interplaque (comme la ceinture de feu du Pacifique), mais il est traversé par la \textbf{Ligne Volcanique du Cameroun} (LVC), une chaîne de volcans active (Mont Cameroun, Manengouba, Bamboutos, Oku, Ngaoundéré) qui s'étend du golfe de Guinée au lac Tchad. Cette structure est associée à un \textbf{rift continental} (fossé d'effondrement) et à des failles transformantes.
\begin{itemize}
    \item \textbf{Mont Cameroun :} Séismes volcaniques fréquents (essaims) précurseurs d'éruptions (1999, 2000). Magnitudes généralement $< 4$.
    \item \textbf{Région de Monatélé / Yaoundé :} Zone de failles anciennes réactivées (faille de Montobé, faille de la Sanaga). Le séisme de 1999 (M $\approx$ 4,8) a été le plus fort ressenti dans la capitale. Des microséismes sont enregistrés régulièrement par le réseau de l'OVSM (Observatoire Volcanologique et Sismologique du Mont Cameroun) et de l'IRGM.
    \item \textbf{Risque :} Faible probabilité de séisme destructeur (M $> 6$), mais \textbf{vulnérabilité élevée} du bâti (parpaings non chaînés, architecture informelle). La prévention passe par l'application des normes de construction (Eurocode 8 adapté / Règlement Parasismique Camerounais) et l'éducation de la population.
\end{itemize}
\end{savaistu}

\begin{attention}[Erreur fréquente : Confondre séisme et volcanisme]
\begin{itemize}
    \item \textbf{Séisme tectonique :} Rupture brutale de roches froides et fragiles par accumulation de contraintes tectoniques (convergence, divergence, cisaillement). Foyer souvent profond (10--700 km). Pas de magma en jeu directement.
    \item \textbf{Séisme volcanique :} Lié à la migration de fluides (magma, gaz, eau) dans l'édifice volcanique. Foyers très superficiels ($< 5$ km), magnitudes faibles ($< 4$), essaims sismiques.
    \item \textbf{Piège :} Dire « le volcan a provoqué un séisme » pour un séisme tectonique loin de tout volcan. Dire « le séisme a provoqué l'éruption » sans preuve de lien causal direct (ouverture de cheminée). Les deux phénomènes coexistent souvent aux limites de plaques (subduction, rift) mais leurs \emph{mécanismes physiques} sont distincts.
\end{itemize}
\end{attention}

\begin{exoresolu}[Analyse de témoignages : Intensité macrosismique]
\textbf{Énoncé :} Lors du séisme de Monatélé (1999), trois témoins rapportent :
\begin{itemize}
    \item Témoin A (Yaoundé, 50 km de l'épicentre) : « J'ai été réveillé, la vaisselle a tinté, j'ai eu peur, mais pas de dégâts chez moi. »
    \item Témoin B (Monatélé, 5 km) : « Les murs de ma maison en parpaings ont fissuré en diagonale aux coins des fenêtres, la cheminée est tombée. »
    \item Témoin C (Mbalmayo, 30 km) : « Difficile de tenir debout, les meubles ont glissé, des fissures dans le sol du jardin. »
\end{itemize}
En utilisant l'échelle macrosismique européenne (EMS-98) simplifiée ci-dessous, estimez l'intensité ressentie par chaque témoin.
\begin{center}
\begin{tabular}{|c|l|}\hline
\textbf{Degré EMS} & \textbf{Description simplifiée} \\ \hline
III & Faible : ressenti par quelques personnes au repos. Objets suspendus oscillent. \\ \hline
IV & Largement ressenti : réveil des dormeurs. Vaisselle tinte. Pas de dégâts. \\ \hline
V & Fort : peur générale. Difficulté à marcher. Petits dégâts (fissures fines, chutes plâtre). \\ \hline
VI & Très fort : dégâts légers à modérés (fissures murs, chutes cheminées). Meubles déplacés. \\ \hline
VII & Fort : dégâts modérés à importants (fissures larges, effondrement partiel). Difficile de se tenir debout. \\ \hline
\end{tabular}
\end{center}

\tcblower
\textbf{Solution détaillée :}
\begin{itemize}
    \item \textbf{Témoin A (Yaoundé) :} « Réveillé, vaisselle tinte, peur, pas de dégâts » $\rightarrow$ Correspond exactement au \textbf{Degré IV} (Largement ressenti). Distance 50 km : atténuation normale.
    \item \textbf{Témoin B (Monatélé) :} « Murs fissurés en diagonale (cisaillement), cheminée tombée » $\rightarrow$ Dégâts modérés sur maçonnerie vulnérable. Correspond au \textbf{Degré VI} (Très fort) voire VII si effondrement partiel. Proximité épicentre (5 km) : fort secouissement.
    \item \textbf{Témoin C (Mbalmayo) :} « Difficile de tenir debout, meubles glissent, fissures sol » $\rightarrow$ Difficulté à se tenir debout = Degré VII. Fissures sol = amplification de site possible (alluvions Nyong). Intensité estimée \textbf{VI--VII}.
\end{itemize}
\textbf{Conclusion :} L'intensité décroît avec la distance à l'épicentre (IV à 50 km, VI-VII à 5-30 km). Elle dépend aussi de la nature du sol (amplification sur alluvions à Mbalmayo) et de la vulnérabilité des bâtiments (parpaings non armés à Monatélé).
\end{exoresolu}

\begin{exercice}[Observation de photos de dégâts]
On vous présente trois photos (non reproduites ici, à imaginer ou chercher) :
\begin{enumerate}
    \item Photo 1 : Une route bitumée coupée net avec un décalage horizontal de 2 m entre les deux bords. Aucune trace de compression verticale.
    \item Photo 2 : Un immeuble de 4 étages dont les 3 premiers étages sont écrasés (hauteur résiduelle $\approx$ 1,5 m), le 4$^{\text{e}}$ posé dessus. Les murs extérieurs sont projetés vers l'extérieur.
    \item Photo 3 : Une zone portuaire où des grues sont basculées, le sol présente des « volcans de sable » (cratères de 50 cm expulsant de l'eau et du sable), des bâtiments penchés de 10--15°.
\end{enumerate}
Pour chaque photo, identifiez : (a) Le type de manifestation (faille de surface, mode de rupture structure, liquéfaction). (b) Le type de sollicitation sismique probable (cisaillement horizontal, vertical, amplification/liquéfaction). (c) Une mesure de prévention constructive adaptée.
\end{exercice}

\begin{corrige}[Exercice n°1]
\begin{enumerate}
    \item \textbf{Photo 1 (Route décalée) :} (a) Rupture de faille de surface en décrochement (transformante). (b) Cisaillement horizontal dominant (ondes S horizontales + déformation permanente). (c) Éviter de construire sur la trace de faille active (zonage) ; fondations profondes traversant la zone de rupture ou isolation sismique à la base.
    \item \textbf{Photo 2 (Immeuble pancake) :} (a) Effondrement progressif par perte de portance des poteaux (rupture fragile béton/acier) $\rightarrow$ mode « pancake ». (b) Fortes accélérations horizontales (période du bâtiment $\approx$ période du site $\rightarrow$ résonance) + absence de chaînages/ferraillage ductile. (c) Conception parasismique : poteaux-poutres ductiles (ferraillage transversal serré), contreventements, dissipation d'énergie, isolation à la base.
    \item \textbf{Photo 3 (Port, volcans de sable, bâtiments penchés) :} (a) Liquéfaction des sols sableux saturés (perte de portance, expulsion eau/sable). (b) Secousses prolongées, basses fréquences, saturation en eau $\rightarrow$ pression interstitielle monte, contrainte efficace $\rightarrow$ 0. (c) Amélioration du sol (drains verticaux, compactage, injection coulis), fondations sur pieux profonds ancrés dans le substratum stable, bâtiments légers ou isolation sismique.
\end{enumerate}
\end{corrige}

\coursec{Origine des séismes : foyer, épicentre et ondes sismiques}

\courssub{Le mécanisme de la rupture élastique (Théorie du rebond élastique)}

\begin{experience}[Modélisation de l'accumulation et de la libération d'énergie]
\textbf{Matériel :} Une planche en bois posée sur une table, un ressort, un dynamomètre, un bloc de bois (représentant un bloc crustal), du papier de verre (rugosité de la faille).
\textbf{Protocole :} On attache le ressort au bloc. On tire lentement sur le dynamomètre (mouvement lent des plaques). Le bloc ne bouge pas d'abord (faille bloquée par frottement), le ressort s'étire (accumulation d'énergie élastique). Brutalement, le bloc glisse (rupture), le ressort se contracte (libération d'énergie).
\textbf{Observation :} Le déplacement du bloc est soudain. La force mesurée chute brutalement.
\textbf{Interprétation :} Les plaques tectoniques bougent lentement (cm/an). Les failles sont verrouillées par le frottement. Les roches se déforment élastiquement (stockage d'énergie). Lorsque la contrainte dépasse la résistance de la roche (ou du frottement), rupture brutale $\rightarrow$ libération de l'énergie élastique stockée $\rightarrow$ ondes sismiques. C'est la \cle{théorie du rebond élastique} (H.F. Reid, 1910).
\end{experience}

\courssub{Les ondes sismiques : ondes P et ondes S}

Au foyer, la rupture génère deux types d'ondes de volume (corps) qui se propagent dans toutes les directions à l'intérieur de la Terre :

\begin{center}
\begin{tabularx}{\linewidth}{|l|X|X|c|}\hline
\rowcolor{popblueL}
\textbf{Caractéristique} & \textbf{Onde P (Primaire / Longitudinale)} & \textbf{Onde S (Secondaire / Transverse)} & \textbf{Onde de surface (L, R)} \\ \hline
\textbf{Mouvement des particules} & Alternance compression/dilatation \textbf{parallèle} à la direction de propagation. & Mouvement \textbf{perpendiculaire} à la direction de propagation (cisaillement). & Mouvement rétrograde (R) ou horizontal (L), confiné près de la surface. \\ \hline
\textbf{Vitesse typique (croûte)} & $\sim$ 5 à 7 km/s & $\sim$ 3 à 4 km/s ($v_P \approx 1,7 \times v_S$) & $<$ vitesse ondes de volume \\ \hline
\textbf{Milieu traversé} & Solides, \textbf{liquides}, gaz. & \textbf{Solides uniquement} (pas de cisaillement dans les fluides). & Surface solide. \\ \hline
\textbf{Dégâts en surface} & Moins destructrices (premières arrivées, secousse verticale). & Très destructrices (amplitude plus grande, secousse horizontale). & \textbf{Les plus destructrices} (amplitude max, durée longue, résonance bâtiments). \\ \hline
\end{tabularx}
\end{center}

\begin{propriete}[Propriétés des ondes sismiques]
\begin{itemize}
    \item Les ondes P arrivent en premier aux stations sismologiques (d'où le nom \emph{Primaire}).
    \item Les ondes S arrivent après (d'où \emph{Secondaire}).
    \item L'intervalle de temps $\Delta t = t_S - t_P$ permet de calculer la \textbf{distance épicentrale} $d$ (loi empirique : $d \approx k \cdot \Delta t$, avec $k \approx 8$ km/s pour la croûte supérieure).
    \item Les ondes de surface (Love et Rayleigh) arrivent en dernier mais causent l'essentiel des dégâts.
\end{itemize}
\end{propriete}

\begin{popfigure}
\begin{tikzpicture}[scale=0.8, every node/.style={font=\small}]
    % Onde P : compression/dilatation
    \node[popdark] at (0, 3.5) {\textbf{Onde P (Longitudinale)}};
    \draw[popblue, thick, ->] (-3, 3) -- (3, 3); % Direction propagation
    \foreach \x in {-2.5,-1.5,-0.5,0.5,1.5,2.5} {
        \draw[popblue, thick] (\x, 2.8) -- (\x, 3.2); % Particules
    }
    % Représentation compression/dilatation
    \draw[popblue, decorate, decoration={snake, amplitude=1pt, segment length=4pt}] (-3, 2.5) -- (3, 2.5);
    \node[popblue] at (0, 2.2) {Compression $\leftrightarrow$ Dilatation};
    
    % Onde S : cisaillement
    \node[popdark] at (0, 1.0) {\textbf{Onde S (Transverse)}};
    \draw[poporange, thick, ->] (-3, 0.5) -- (3, 0.5); % Direction propagation
    \foreach \x in {-2.5,-1.5,-0.5,0.5,1.5,2.5} {
        \draw[poporange, thick] (\x, 0.0) -- (\x, 1.0); % Particules verticales
    }
    \draw[poporange, thick, domain=-3:3, samples=50, smooth] plot (\x, {0.5 + 0.5*sin(\x*180)});
    \node[poporange] at (0, -0.3) {Cisaillement (haut/bas ou gauche/droite)};
    
    % Onde de surface (Rayleigh) - mouvement rétrograde
    \node[popdark] at (0, -1.8) {\textbf{Onde de Rayleigh (Surface)}};
    \draw[popgreen, thick, ->] (-3, -2.2) -- (3, -2.2);
    \foreach \x in {-2.5,-1.5,-0.5,0.5,1.5,2.5} {
        \draw[popgreen, thick, ->] (\x, -2.2) -- (\x + 0.15, -2.5) -- (\x, -2.8) -- (\x - 0.15, -2.5) -- cycle; % Ellipse rétrograde
    }
    \node[popgreen] at (0, -3.2) {Mouvement elliptique rétrograde (solide "roule")};
\end{tikzpicture}
\legende{Comparaison des modes de propagation des ondes P, S et de surface (Rayleigh). Les flèches indiquent le mouvement des particules de roche. L'onde P comprime et dilate ; l'onde S cisaile ; l'onde de Rayleigh fait "rouler" la surface.}
\end{popfigure}

\courssub{Localisation de l'épicentre et détermination de la profondeur du foyer}

\begin{methode}[Localisation de l'épicentre par triangulation (méthode des cercles)]
\begin{enumerate}
    \item Sur un sismogramme (enregistrement d'une station), mesurer l'intervalle de temps $\Delta t = t_S - t_P$ entre l'arrivée de la première onde P et de la première onde S.
    \item Convertir $\Delta t$ en distance épicentrale $d$ grâce à une courbe de marche (ou table de marche) standardisée (ex: $d \approx 8 \times \Delta t$ pour $\Delta t$ en secondes et $d$ en km, pour des distances < 500 km).
    \item Sur une carte, tracer un cercle de rayon $d$ centré sur la station. L'épicentre se trouve sur ce cercle.
    \item Répéter l'opération pour \textbf{au moins trois stations} distinctes.
    \item L'intersection des trois cercles donne la position de l'\textbf{épicentre}.
\end{enumerate}
\end{methode}

\begin{definition}[Profondeur du foyer (Hypocentre)]
La profondeur $h$ du foyer s'obtient en utilisant la relation de Pythagore si on connaît la distance épicentrale $d$ (distance horizontale) et la distance hypocentrale $D$ (distance droite foyer-station, déduite des courbes de marche pour les ondes profondes) : $D^2 = d^2 + h^2$. On classe les séismes : \textbf{superficiels} ($h < 70$ km, les plus destructeurs), \textbf{intermédiaires} ($70 < h < 300$ km), \textbf{profonds} ($h > 300$ km, zones de subduction).
\end{definition}

\coursec{Évaluation de l'intensité et de la magnitude}

\courssub{L'intensité macrosismique (Échelle EMS-98)}

\begin{definition}[Intensité macrosismique]
L'\cle{intensité} d'un séisme en un lieu donné est une mesure \textbf{qualitative} des effets du séisme sur les personnes, les objets, les bâtiments et le milieu naturel. Elle ne dépend pas d'un instrument unique mais d'une enquête de terrain. Elle varie avec la distance à l'épicentre et la nature du sol.
\end{definition}

L'échelle utilisée en Europe et au Cameroun est l'\textbf{EMS-98} (European Macroseismic Scale), à 12 degrés (I à XII). Les degrés usuels au BEPC sont :
\begin{itemize}
    \item \textbf{IV} (Largement ressenti) : Réveil, vaisselle tinte, pas de dégâts.
    \item \textbf{V} (Fort) : Peur, difficulté à marcher, petits dégâts (plâtre, fissures fines).
    \item \textbf{VI} (Très fort) : Dégâts légers à modérés (fissures murs, chutes cheminées), meubles déplacés.
    \item \textbf{VII} (Fort) : Dégâts modérés à importants, difficile de se tenir debout.
    \item \textbf{VIII et +} : Destructions majeures, ruines.
\end{itemize}

\begin{aretenir}[Intensité vs Distance et Sol]
L'intensité \textbf{décroît avec la distance} à l'épicentre (atténuation géométrique + absorption). Elle est \textbf{amplifiée sur sols meubles} (alluvions, remblais) par rapport au roc dur (amplification de site). Une carte d'isosistes (lignes joignant les points de même intensité) est généralement elliptique (allongée dans la direction de la faille).
\end{aretenir}

\courssub{La magnitude (Échelle de Richter / Magnitude de moment Mw)}

\begin{definition}[Magnitude]
La \cle{magnitude} est une mesure \textbf{quantitative} de l'\textbf{énergie libérée} au foyer. Elle est unique pour un séisme donné (ne varie pas avec la distance). Elle se calcule à partir de l'amplitude maximale des ondes enregistrées sur un sismogramme standard (Wood-Anderson pour Richter $M_L$) ou, pour les grands séismes, par la \textbf{Magnitude de moment $M_w$} (basée sur le moment sismique $M_0 = \mu \cdot S \cdot D$ : rigidité $\times$ surface de rupture $\times$ glissement moyen).
\end{definition}

\begin{propriete}[Échelle logarithmique]
L'échelle de magnitude est \textbf{logarithmique} (base 10 pour l'amplitude) :
\begin{itemize}
    \item Une augmentation de 1 unité de magnitude = \textbf{amplitude des ondes multipliée par 10} (sur le sismogramme).
    \item Une augmentation de 1 unité de magnitude = **énergie libérée multipliée par $\approx 32$** ($\approx 10^{1,5}$).
    \item Exemple : Un séisme Mw 7,0 libère $\approx 32 \times$ plus d'énergie qu'un Mw 6,0, et $\approx 1000 \times$ plus qu'un Mw 5,0.
\end{itemize}
\end{propriete}

\begin{attention}[Ne pas confondre Intensité et Magnitude]
\begin{center}
\begin{tabular}{|l|c|c|}\hline
\rowcolor{popblueL}
\textbf{Critère} & \textbf{Intensité (EMS)} & \textbf{Magnitude (Mw, $M_L$)} \\ \hline
\textbf{Nature} & Qualitative (effets observés) & Quantitative (énergie, amplitude) \\ \hline
\textbf{Variabilité} & Change selon le lieu (distance, sol) & Unique pour le séisme (caractère intrinsèque) \\ \hline
\textbf{Unité} & Degrés (I à XII, chiffres romains) & Unités décimales (ex: 6,7) \\ \hline
\textbf{Utilité} & Estimation dégâts, secours, zonage & Comparaison taille séismes, études source \\ \hline
\end{tabular}
\end{center}
\textbf{Piège BEPC :} Dire « l'intensité du séisme est de 7 sur l'échelle de Richter ». \textbf{Faux.} On dit : « La magnitude est de 7,0 (Mw) » ET « L'intensité maximale observée est de VIII (EMS) ».
\end{attention}

\coursec{Répartition mondiale des séismes}

\courssub{Le lien avec la tectonique des plaques}

\begin{propriete}[Répartition des séismes]
La très grande majorité des séismes (\textbf{> 90\%} de l'énergie sismique totale) se concentre le long de \textbf{trois types de limites de plaques lithosphériques} :
\begin{itemize}
    \item \textbf{Zones de subduction (convergence océan/océan ou océan/continent) :} Séismes très fréquents, foyers superficiels à très profonds (jusqu'à 700 km), \textbf{plus grandes magnitudes mondiales} (Mw > 9 possible). Ex: Japon, Chili, Indonésie, Alaska. Plan de Wadati-Benioff incliné.
    \item \textbf{Dorsales médio-océaniques (divergence) :} Séismes fréquents, foyers \textbf{très superficiels} ($< 30$ km), magnitudes modérées (généralement $< 7$). Ex: Dorsale Atlantique, Pacifique Est.
    \item \textbf{Failles transformantes (cisaillement) :} Séismes fréquents, foyers superficiels ($< 20$ km), magnitudes fortes possibles (jusqu'à $\sim$ 8,5). Ex: Faille de San Andreas (USA), Faille de l'Anatolie du Nord (Turquie).
\end{itemize}
\end{propriete}

\begin{itemize}
    \item \textbf{Séismes intraplaques :} Plus rares, mais possibles à l'intérieur des plaques (réactivation de failles anciennes, points chauds). Ex: Séisme de Monatélé (Cameroun), New Madrid (USA), Bassin du Congo. Magnitudes généralement modérées mais dangerosité sous-estimée car bâti non préparé.
\end{itemize}

\begin{savaistu}[La "Ceinture de Feu" du Pacifique]
C'est la zone la plus sismiquement active au monde (environ 80\% des grands séismes). Elle ceinture l'océan Pacifique : Amérique du Nord (Alaska $\rightarrow$ Californie $\rightarrow$ Mexique), Amérique du Sud (Chili, Pérou), Asie (Japon, Philippines, Indonésie), Océanie (Nouvelle-Zélande). Elle correspond aux zones de subduction du Pacifique sous les plaques voisines.
\end{savaistu}

\courssub{La sismicité au Cameroun : contexte local}

Comme vu précédemment, le Cameroun se situe sur la \textbf{Ligne Volcanique du Cameroun (LVC)}, structure de rift intracontinental.
\begin{itemize}
    \item \textbf{Zones actives :} Mont Cameroun (volcanique), Région de l'Adamaoua / Ngaoundéré (rift), Région de Monatélé/Yaoundé (faille de Montobé, faille de la Sanaga - failles de cisaillement/rejet réactivées).
    \item \textbf{Surveillance :} L'OVSM (Observatoire Volcanologique et Sismologique du Mont Cameroun, basé à Buéa) et l'IRGM (Institut de Recherche Géologique et Minière) gèrent un réseau de stations sismiques (Buéa, Yaoundé, Ngaoundéré, Maroua, etc.).
    \item \textbf{Risque sismique = Aléa $\times$ Vulnérabilité $\times$ Enjeux.} Aléa modéré, mais Vulnérabilité forte (bâti informel, parpaings non chaînés) et Enjeux croissants (urbanisation rapide de Yaoundé, Douala).
\end{itemize}

\coursec{Prévision et prévention des séismes}

\courssub{La prévision : un défi scientifique non résolu}

\begin{definition}[Prévision vs Prédiction]
\begin{itemize}
    \item \textbf{Prévision (Long terme / Probabiliste) :} Estimation de la probabilité qu'un séisme d'une certaine magnitude se produise dans une zone donnée sur une période donnée (ex: 30 ans). Basée sur l'historique sismique, la vitesse de déplacement des plaques (GPS), la charge de contrainte. \textbf{C'est possible et opérationnel (cartes d'aléa).}
    \item \textbf{Prédiction (Court terme / Déterministe) :} Annoncer \textbf{la date, le lieu et la magnitude} d'un séisme à venir avec une certitude suffisante pour évacuer. \textbf{Actuellement impossible de façon fiable et reproductible.}
\end{itemize}
\end{definition}

\begin{attention}[Les "précurseurs" non fiables]
Aucun précurseur (gaz radon, niveau nappe phréatique, comportement animal, signaux électromagnétiques, "lumières sismiques") n'a montré une corrélation \textbf{systématique, spécifique et suffisamment précoce} pour servir de base à une évacuation. De fausses alertes provoquent panique, pertes économiques et discrédit des autorités. La stratégie officielle mondiale (et au Cameroun) est la \textbf{prévention}, pas l'évacuation préventive basée sur une prédiction.
\end{attention}

\courssub{La prévention : réduire la vulnérabilité}

C'est le seul levier d'action efficace. Elle repose sur trois piliers :

\begin{enumerate}
    \item \textbf{Connaissance de l'aléa et Zonage sismique :} Cartes d'aléa (probabilité d'accélération du sol dépassée pour une période de retour, ex: 475 ans). Définition de zones de réglementation parasismique (au Cameroun : Zones 1 à 4 selon le Règlement Parasismique Camerounais, inspiré de l'Eurocode 8). \textbf{Interdiction de construire sur failles actives connues.}
    \item \textbf{Construction parasismique (Règles de l'art) :}
    \begin{itemize}
        \item \textbf{Simplicité et régularité :} Plan symétrique, peu d'encorbellements, masses réparties (éviter torsion).
        \item \textbf{Chaînages :} Chaînages horizontaux (linteaux, ceintures d'étage) et verticaux (poteaux) en béton armé \textbf{obligatoires} pour lier les murs (parpaings, briques) et éviter l'écartement.
        \item \textbf{Ductilité :} Ferraillage transversal serré (étriers) dans les poteaux et poutres pour confiner le béton et permettre déformation plastique sans rupture fragile.
        \item \textbf{Fondations :} Semelles filantes ou radier, ancrées dans le sol sain. Isolation à la base (appuis en caoutchouc/acier) pour les bâtiments stratégiques (hôpitaux, écoles, tours).
        \item \textbf{Matériaux locaux améliorés :} Terre stabilisée (ciment/chaux), briques compressées (BTC) avec chaînages, bois/bambou (flexibilité) pour l'habitat rural.
    \end{itemize}
    \item \textbf{Éducation, préparation et gestion de crise :}
    \begin{itemize}
        \item Exercices de simulation (Drop, Cover, Hold On / "Baisez-vous, Protégez-vous, Attendez").
        \item Sac d'urgence (eau, lampe, radio, papiers, médicaments).
        \item Plan familial de rassemblement.
        \item Après le choc : ne pas rallumer gaz/électricité si fuite suspectée, écouter radio officielle, ne pas utiliser téléphone (saturer réseaux secours), aider voisins, ne pas rentrer dans bâtiments fissurés sans expertise.
    \end{itemize}
\end{enumerate}

\begin{methode}[Conduite à tenir pendant un séisme (Mémo BEPC)]
\begin{enumerate}
    \item \textbf{À l'intérieur :} S'abriter sous une table solide (protection chute objets), loin des fenêtres/miroirs/armoires. Ne pas courir vers la sortie (risque chutes, vitres, foule). Ne pas prendre l'ascenseur.
    \item \textbf{À l'extérieur :} S'éloigner des bâtiments (chutes tuiles, corniches, fils électriques), poteaux, arbres. Se mettre en place ouverte.
    \item \textbf{En voiture :} S'arrêter côté route (loin ponts, tunnels, lignes haute tension), rester dans le véhicule, frein à main.
    \item \textbf{Après la secousse :} Vérifier blessures, couper gaz/électricité/eau si dégâts, sortir calmement avec sac urgence, rejoindre point de rassemblement, écouter consignes autorités.
\end{enumerate}
\end{methode}

\begin{exemplebox}[Le Règlement Parasismique Camerounais (RPC) et l'Eurocode 8]
Le Cameroun dispose d'un cadre réglementaire (Arrêté ministériel) définissant les règles de calcul et de construction parasismique, largement inspiré de l'Eurocode 8 (norme européenne). Il classe le territoire en zones de sismicité (Faible, Modérée, Forte) et impose des règles de conception (coefficient de comportement $q$, accélération de référence $a_{gR}$) pour les ouvrages neufs. Le défi majeur reste l'application effective sur le terrain (contrôle des bureaux d'études, des entreprises, auto-construction).
\end{exemplebox}

\begin{exoresolu}[Calcul de distance épicentrale et magnitude]
\textbf{Énoncé :} Une station sismologique enregistre un séisme. L'onde P arrive à $t_P = 10$ h $12$ min $05$ s. L'onde S arrive à $t_S = 10$ h $12$ min $37$ s. L'amplitude maximale de l'onde S (composante horizontale) sur le sismogramme (simulation Wood-Anderson) est de $A = 25$ mm. La période de l'onde est $T = 0,5$ s. La station est à $d = 200$ km de l'épicentre (vérifié par triangulation).
\begin{enumerate}
    \item Calculer l'intervalle $\Delta t = t_S - t_P$ et vérifier la distance $d$ avec la formule approximative $d \approx 8 \times \Delta t$ ($\Delta t$ en s, $d$ en km).
    \item Calculer la magnitude locale $M_L$ (Richter) avec la formule simplifiée : $M_L = \log_{10}(A) - \log_{10}(A_0)$.
    On donne $\log_{10}(A_0)$ pour $d=200$ km et $T=0,5$ s : $\log_{10}(A_0) \approx -1,3$ (valeur tabulée).
\end{enumerate}

\tcblower
\textbf{Solution détaillée :}
\begin{enumerate}
    \item $\Delta t = 37\,\text{s} - 05\,\text{s} = 32\,\text{s}$.
    $d_{\text{calc}} \approx 8 \times 32 = 256\,\text{km}$.
    \textbf{Commentaire :} La formule $d \approx 8 \Delta t$ est une approximation pour la croûte supérieure. La distance réelle (200 km) diffère car la vitesse moyenne dépend de la structure géologique traversée. La triangulation (3 stations) donne la valeur fiable (200 km).
    \item $A = 25\,\text{mm}$. $\log_{10}(25) \approx 1,398$.
    $M_L = 1,398 - (-1,3) = 1,398 + 1,3 = 2,698 \approx \textbf{2,7}$.
    \textbf{Interprétation :} C'est un séisme de faible magnitude (micro-séisme), ressenti localement près de l'épicentre (Intensité III-IV EMS) mais sans dégâts.
\end{enumerate}
\end{exoresolu}

\begin{exercice}[Synthèse : Le séisme de Turquie-Syrie du 6 février 2023]
Ce doublet sismique (Mw 7,8 puis Mw 7,5, 9 heures plus tard) a frappé la faille de l'Anatolie Orientale (failles transformantes).
\begin{enumerate}
    \item Rappeler le type de limite de plaque et le mécanisme de rupture (failles) associé.
    \item Expliquer pourquoi la profondeur des foyers ($\sim$ 10-18 km) a aggravé les secousses en surface.
    \item L'intensité maximale observée a atteint XII (EMS) par endroits. Expliquer ce degré par rapport à la magnitude (7,8) et à la vulnérabilité locale (bâtiments en béton non ductiles, "pancake").
    \item Citer trois mesures de prévention constructive qui auraient pu limiter le nombre d'effondrements.
    \item Pourquoi la seconde secousse (Mw 7,5) a-t-elle causé des effondrements supplémentaires alors qu'elle était de magnitude inférieure ?
\end{enumerate}
\end{exercice}

\begin{corrige}[Exercice n°2]
\begin{enumerate}
    \item \textbf{Limite :} Frontière de plaques transformante (cisaillement) entre la plaque Arabique et la plaque Anatolienne (micro-plaque eurasienne). \textbf{Mécanisme :} Rupture en décrochement (principalement dextre) le long de segments de faille. Libération brutale de contraintes de cisaillement accumulées.
    \item \textbf{Profondeur superficielle :} Moins d'atténuation géométrique et anélastique des ondes avant d'atteindre la surface. Accélérations du sol (PGA) très élevées ($> 1\,g$ par endroits) au plus près de la rupture. Ondes de surface très énergétiques.
    \item \textbf{Intensité XII (Destruction totale) :} Magnitude grande (grande surface de rupture $\approx$ 300 km pour le 1er, durée longue $\approx$ 60-80s) $\rightarrow$ fort secouissement sur vaste zone. + Vulnérabilité extrême : Immeubles en béton armé "poteaux-poutres" non conçus pour le séisme (pas d'étriers serrés, acier lisse, joints faibles, étages souples rez-de-chaussée commerciaux) $\rightarrow$ Rupture fragile des poteaux $\rightarrow$ Mécanisme "pancake" (empilement planchers). Résonance possible (période site $\approx$ période bâtiment).
    \item \textbf{Mesures :} (1) Chaînages verticaux/horizontaux rigoureux (étriers rapprochés aux nœuds poteau-poutre). (2) Contreventements (voiles en béton armé) pour rigidifier/assurer ductilité. (3) Isolation sismique à la base (appuis pendulaires/élastomères) pour bâtiments stratégiques. (4) Respect des règles de régularité en plan/élévation (pas d'étage souple).
    \item \textbf{Répliques / Second choc :} Le 2ème séisme (Mw 7,5) n'est pas une réplique classique (sur même faille) mais un séisme déclenché sur faille adjacente (faille de Sürgü/Cardak) par transfert de contrainte (Coulomb stress transfer). Les bâtiments déjà fissurés/affaiblis par le 1er choc (résidu de résistance nul) s'effondrent sous une accélération encore forte (proximité épicentre 2ème choc). Secours piégés.
\end{enumerate}
\end{corrige}

\coursec{Origine des séismes : foyer, épicentre et ondes}

Dans la section précédente, nous avons décrit les \cle{manifestations} des séismes : secousses du sol, fissures, effondrements de constructions, glissements de terrain, parfois raz-de-marée. Mais une question fondamentale reste posée : \textbf{d'où vient un séisme ?} Où naît-il, et comment les vibrations parviennent-elles jusqu'à nous ? C'est ce que nous allons découvrir maintenant, en suivant la démarche des scientifiques : observer, émettre des hypothèses, puis conclure.

\courssub{Le foyer : là où naît le séisme}

\textbf{Observation.} Après un séisme, on constate souvent à la surface du sol un \cle{décalage des couches rocheuses} : une route, une clôture ou une rivière se retrouvent coupées en deux parties déplacées l'une par rapport à l'autre, parfois de plusieurs mètres. Ce décalage se fait le long d'une fracture appelée \cle{faille}.

\textbf{Hypothèse.} Les roches de la lithosphère ne sont pas immobiles : elles sont entraînées par les mouvements des plaques. Elles subissent alors des \cle{contraintes} qui les déforment progressivement.

\textbf{Interprétation.} Les roches se déforment d'abord de façon \emph{élastique}, comme un bâton que l'on plie sans le casser : elles accumulent de l'énergie élastique de déformation. Mais lorsque la contrainte dépasse la résistance mécanique des roches, celles-ci \textbf{se rompent brutalement} le long d'une faille préexistante ou nouvelle. L'énergie accumulée est libérée en un instant sous forme d'ondes mécaniques : c'est le séisme. Ce mécanisme est appelé \cle{rebond élastique}.

\begin{definition}[Foyer et épicentre]
Le \cle{foyer} (ou \cle{hypocentre}) est la zone située \textbf{en profondeur} où les roches se rompent brutalement : c'est là que \textbf{naît} le séisme. L'\cle{épicentre} est le point de la \textbf{surface} terrestre situé \textbf{à la verticale du foyer}.
\end{definition}

\begin{attention}[Foyer ou épicentre ?]
L'épicentre n'est \textbf{pas} le lieu où le séisme naît ! La rupture des roches a lieu au \textbf{foyer}, en profondeur. L'épicentre n'est que sa \emph{projection verticale à la surface}. C'est pourtant à l'épicentre que les dégâts sont généralement \textbf{maximaux}, car c'est le point de la surface le plus proche du foyer.
\end{attention}

\courssub{La profondeur du foyer}

Les sismologues ont mesuré la profondeur des foyers de très nombreux séismes. On distingue trois catégories :

\begin{center}
\begin{tabularx}{\linewidth}{|l|c|X|}
\hline
\rowcolor{popblueL} \textbf{Type de séisme} & \textbf{Profondeur du foyer} & \textbf{Remarque} \\
\hline
Séisme \cle{superficiel} & \SI{0}{\kilo\meter} à \SI{70}{\kilo\meter} & Les plus fréquents (75\,\%) et les plus destructeurs \\
\hline
Séisme \cle{intermédiaire} & \SI{70}{\kilo\meter} à \SI{300}{\kilo\meter} & Moins fréquents, liés aux zones de subduction \\
\hline
Séisme \cle{profond} & \SI{300}{\kilo\meter} à \SI{700}{\kilo\meter} & Ressentis sur de vastes régions mais dégâts souvent atténués en surface \\
\hline
\end{tabularx}
\end{center}

\begin{aretenir}
Plus le foyer est \textbf{superficiel}, plus les vibrations arrivent fortes à la surface (peu d'atténuation géométrique et anélastique) : les séismes superficiels sont donc les plus dangereux pour les populations.
\end{aretenir}

\courssub{Les ondes sismiques : des vibrations qui voyagent}

Au moment de la rupture, l'énergie libérée se propage à partir du foyer \textbf{dans toutes les directions}, sous forme de vibrations appelées \cle{ondes sismiques}. On distingue deux grandes familles d'ondes.

\textbf{Les ondes de volume}, qui traversent l'intérieur de la Terre :
\begin{itemize}
\item les \cle{ondes P} (ondes \emph{primaires} ou de compression) : ce sont les plus \textbf{rapides} ($v_P \approx \SI{5}{\kilo\meter/\second}$ à \SI{8}{\kilo\meter/\second} dans la croûte), elles arrivent donc en premier aux stations d'enregistrement. Elles compriment puis dilatent les roches dans le sens de la propagation (compression--dilatation), comme un ressort que l'on pousse et que l'on relâche. Elles traversent les solides \textbf{et} les liquides ;
\item les \cle{ondes S} (ondes \emph{secondaires} ou de cisaillement) : plus \textbf{lentes} ($v_S \approx \SI{3}{\kilo\meter/\second}$ à \SI{5}{\kilo\meter/\second} dans la croûte), elles arrivent en second. Elles font vibrer les roches par \emph{cisaillement} (mouvement perpendiculaire à la direction de propagation). Elles \textbf{ne traversent pas les liquides} (module de cisaillement nul).
\end{itemize}

\textbf{Les ondes de surface} (ondes \cle{L}, ondes de Love et de Rayleigh) : elles se propagent le long de la surface terrestre et font onduler le sol (mouvements horizontaux et/ou verticaux). Leur amplitude diminue moins vite avec la distance que les ondes de volume. Ce sont elles qui sont responsables de \textbf{l'essentiel des dégâts} aux constructions.

\begin{propriete}[Vitesse des ondes (admise)]
Les ondes P se propagent \textbf{plus vite} que les ondes S ($v_P > v_S$). En un point donné, les ondes P arrivent donc \emph{avant} les ondes S : c'est l'\textbf{écart de temps d'arrivée} $\Delta t = t_S - t_P$ entre les ondes P et S qui permet, à l'aide d'une courbe temps-distance (hodochrone), de calculer la distance $d$ entre la station sismique et l'épicentre.
\end{propriete}

\begin{popfigure}
\begin{center}
\begin{tikzpicture}[scale=1.0]
% Surface
\fill[popgreenL] (-5,0) rectangle (5,0.35);
\draw[popdark, thick] (-5,0) -- (5,0);
% Sous-sol
\fill[popgoldL] (-5,-3.2) rectangle (5,0);
\draw[popdark, thick] (-5,0) -- (5,0);
% Faille inclinée
\draw[popink, very thick] (-1.3,0) -- (0.4,-3.2);
\node[popink, rotate=68] at (-0.9,-1.9) {\small faille};
% Foyer
\fill[poporange] (0.2,-2.4) circle (0.14);
\node[poporange, below right] at (0.25,-2.55) {\small \textbf{Foyer}};
% Ondes concentriques (représentation schématique)
\draw[poporange, dashed] (0.2,-2.4) circle (0.55);
\draw[poporange, dashed] (0.2,-2.4) circle (1.05);
\draw[poporange, dashed] (0.2,-2.4) circle (1.55);
% Épicentre
\fill[popblue] (0.2,0.35) circle (0.12);
\node[popblue, above] at (0.2,0.5) {\small \textbf{Épicentre}};
% Verticale foyer-épicentre
\draw[popdark, densely dotted] (0.2,-2.4) -- (0.2,0.35);
% Profondeur
\draw[<->, popdark] (-4.6,0) -- (-4.6,-2.4) node[midway, left] {\small profondeur $h$};
% Flèches ondes
\draw[->, poppurple, thick] (0.2,-2.4) -- (2.6,-1.2) node[right] {\small ondes sismiques};
\draw[->, poppurple, thick] (0.2,-2.4) -- (-2.6,-1.0);
\end{tikzpicture}
\end{center}
\legende{Coupe du sous-sol lors d'un séisme. Le séisme naît au \textbf{foyer}, par rupture des roches le long d'une \textbf{faille} inclinée. Les ondes sismiques (cercles en pointillés) se propagent dans toutes les directions. L'\textbf{épicentre} est le point de la surface situé à la verticale du foyer. La profondeur $h$ est la distance foyer-épicentre.}
\end{popfigure}

\courssub{Le mécanisme du rebond élastique}

Comment les roches passent-elles de la déformation à la rupture ? Le schéma ci-dessous résume les trois étapes.

\begin{popfigure}
\begin{center}
\begin{tikzpicture}[scale=0.95]
% Étape 1 : Repos
\begin{scope}[shift={(0,0)}]
\draw[popdark, thick] (-1.4,-1) rectangle (1.4,1);
\draw[popink, thick] (0,-1) -- (0,1); % Faille
\draw[popmuted] (-1.4,0.5) -- (1.4,0.5); % Repères couches
\draw[popmuted] (-1.4,-0.5) -- (1.4,-0.5);
\node at (0,1.35) {\small \textbf{1. Roches non déformées}};
\end{scope}
% Étape 2 : Déformation élastique
\begin{scope}[shift={(5,0)}]
\draw[popdark, thick] (-1.4,-1) -- (1.4,-1) -- (1.0,1) -- (-1.8,1) -- cycle; % Bloc déformé
\draw[popink, thick] (-0.2,-1) -- (-0.6,1); % Faille déformée
\draw[popmuted] (-1.5,0.5) -- (1.3,0.5); % Couches courbées
\draw[popmuted] (-1.3,-0.5) -- (1.1,-0.5);
\draw[->, poporange, thick] (-2.2,0) -- (-1.7,0);
\draw[->, poporange, thick] (2.2,0) -- (1.7,0);
\node at (-0.2,1.35) {\small \textbf{2. Déformation élastique}};
\node[poporange] at (0,-1.5) {\small contraintes $\sigma$};
\end{scope}
% Étape 3 : Rupture
\begin{scope}[shift={(10,0)}]
\draw[popdark, thick] (-1.4,-1) -- (-0.1,-1) -- (0.5,1) -- (-1.8,1) -- cycle;
\draw[popdark, thick] (0.3,-1) -- (1.4,-1) -- (1.0,1) -- (0.9,1) -- cycle;
\draw[popink, very thick] (-0.1,-1) -- (0.5,1); % Faille rompue
\draw[popink, very thick] (0.3,-1) -- (0.9,1);
\draw[->, poppurple, thick] (0.7,0.3) -- (1.5,0.3);
\draw[->, poppurple, thick] (0.0,-0.3) -- (-0.8,-0.3);
\node at (0,1.35) {\small \textbf{3. Rupture et décalage}};
\node[poppurple] at (0.4,-1.5) {\small libération énergie};
\end{scope}
\end{tikzpicture}
\end{center}
\legende{Le rebond élastique en trois étapes : (1) les couches rocheuses sont au repos ; (2) les contraintes liées aux mouvements des plaques les déforment de façon élastique, l'énergie s'accumule ; (3) la rupture brutale le long de la faille décale les blocs et libère l'énergie sous forme d'ondes sismiques.}
\end{popfigure}

\begin{savaistu}[La faille de San Andreas]
La théorie du rebond élastique a été proposée par le géologue américain \textbf{H. F. Reid} après le grand séisme de \textbf{San Francisco (1906)}. En observant la région, il a constaté que des clôtures et des routes qui traversaient la faille de San Andreas avaient été décalées de \textbf{plusieurs mètres} (jusqu'à \SI{6}{\meter}). Il en a déduit que les roches s'étaient d'abord déformées lentement, puis avaient « rebondi » brutalement en se rompant : le séisme ! Cette faille transformante marque la limite entre la plaque Pacifique et la plaque Nord-Américaine.
\end{savaistu}

\courssub{Localiser l'épicentre : la méthode des trois stations}

\begin{methode}[Localisation d'un épicentre par triangulation]
\begin{enumerate}
\item Sur le sismogramme de chaque station, on repère l'arrivée des ondes P ($t_P$) et des ondes S ($t_S$). On mesure l'\textbf{écart de temps} $\Delta t = t_S - t_P$.
\item À l'aide de la \textbf{courbe temps-distance (hodochrone)} de référence (ou d'un tableau de correspondance), on déduit la \textbf{distance épicentrale} $d$ correspondant à cet écart $\Delta t$.
\item Sur une carte, autour de chaque station, on trace un \textbf{cercle} de rayon $d$ (à l'échelle de la carte).
\item Avec les enregistrements d'\textbf{au moins trois stations}, les trois cercles se coupent en un point unique (ou une toute petite zone) : c'est \textbf{l'épicentre}.
\end{enumerate}
\end{methode}

\begin{exoresolu}[Trouver l'épicentre à partir des distances épicentrales]
Trois stations sismiques A, B et C enregistrent un même séisme. L'analyse des sismogrammes (écart P-S et hodochrone) donne les distances épicentrales suivantes : station A : \SI{300}{\kilo\meter} ; station B : \SI{450}{\kilo\meter} ; station C : \SI{200}{\kilo\meter}. Comment trouver la position de l'épicentre ?
\tcblower
\textbf{Solution.}
\begin{enumerate}
\item On reporte les trois stations A, B, C sur une carte à l'échelle connue.
\item Autour de A, on trace un cercle de rayon \SI{300}{\kilo\meter} (échelle respectée) ; autour de B, un cercle de rayon \SI{450}{\kilo\meter} ; autour de C, un cercle de rayon \SI{200}{\kilo\meter}.
\item Le \textbf{point d'intersection des trois cercles} est l'épicentre du séisme.
\end{enumerate}
\textbf{Interprétation.} Une seule station ne donnerait qu'un cercle (l'épicentre peut être n'importe où sur ce cercle) ; deux stations donneraient deux points d'intersection (ambiguïté) ; il faut donc \textbf{au moins trois stations} non alignées pour obtenir un point unique par intersection des trois cercles.
\end{exoresolu}

\begin{exercice}[À toi de jouer]
Lors d'un séisme, une station enregistre l'arrivée des ondes P à \SI{14}{\hour} \SI{02}{\minute} \SI{10}{\second} et l'arrivée des ondes S à \SI{14}{\hour} \SI{02}{\minute} \SI{50}{\second}.
\begin{enumerate}
\item Quelle onde est la plus rapide ? Justifie ta réponse.
\item Calcule l'écart de temps d'arrivée $\Delta t$ entre les deux ondes.
\item En utilisant l'hodochrone ci-dessous, détermine la distance épicentrale approximative.
\item Explique pourquoi cette seule station ne suffit pas à localiser l'épicentre.
\end{enumerate}

\begin{center}
\begin{tikzpicture}
\begin{axis}[
    width=0.8\linewidth, height=5cm,
    xlabel={Distance épicentrale $d$ (km)}, ylabel={Temps de parcours $t$ (s)},
    domain=0:1000, samples=100,
    xmin=0, xmax=1000, ymin=0, ymax=200,
    grid=major, legend pos=north west
]
\addplot[popblue, thick] {x/6}; % Onde P approx 6 km/s
\addplot[poporange, thick] {x/3.5}; % Onde S approx 3.5 km/s
\legend{Ondes P, Ondes S}
\end{axis}
\end{tikzpicture}
\end{center}
\emph{Indication : lire sur l'axe des abscisses la distance pour laquelle l'écart vertical entre les deux courbes vaut $\Delta t$.}
\end{exercice}

\begin{corrige}[Exercice]
\begin{enumerate}
\item Les \textbf{ondes P} sont les plus rapides : elles arrivent en premier à la station (\SI{14}{\hour} \SI{02}{\minute} \SI{10}{\second}, avant les ondes S à \SI{14}{\hour} \SI{02}{\minute} \SI{50}{\second}).
\item $\Delta t = t_S - t_P = \SI{14}{\hour} \SI{02}{\minute} \SI{50}{\second} - \SI{14}{\hour} \SI{02}{\minute} \SI{10}{\second} = \textbf{\SI{40}{\second}}$.
\item Sur l'hodochrone, on cherche la distance $d$ pour laquelle $t_S(d) - t_P(d) = \SI{40}{\second}$. Graphiquement (ou par calcul $d(1/v_S - 1/v_P) = 40$), on trouve environ \textbf{\SI{350}{\kilo\meter}} (avec $v_P=\SI{6}{\kilo\meter/\second}, v_S=\SI{3.5}{\kilo\meter/\second}$ : $d(1/3.5 - 1/6) = 40 \Rightarrow d \approx 350$).
\item Cet écart $\Delta t$ ne donne que la \emph{distance} $d$ entre la station et l'épicentre : l'épicentre peut se trouver n'importe où sur le cercle de centre la station et de rayon $d$. Il faut les enregistrements d'au moins trois stations pour obtenir un point unique par intersection des trois cercles.
\end{enumerate}
\end{corrige}

\begin{aretenir}[L'essentiel]
Un séisme naît au \cle{foyer}, en profondeur, par \textbf{rupture brutale des roches} le long d'une faille, après une déformation élastique (rebond élastique). L'\cle{épicentre} est le point de la surface à la verticale du foyer, où les dégâts sont maximaux. L'énergie se propage par des \cle{ondes sismiques} : ondes P (rapides, compression, traversent solides et liquides), ondes S (plus lentes, cisaillement, arrêtées par les liquides) et ondes de surface L (les plus destructrices). L'écart d'arrivée $\Delta t = t_S - t_P$, mesuré par au moins trois stations et converti en distance via une hodochrone, permet de localiser l'épicentre par triangulation.
\end{aretenir}

\coursec{Évaluation de l'intensité et de la magnitude}

Dans la section précédente, nous avons vu qu'un séisme naît en un point de la profondeur appelé \cle{foyer}, et que les vibrations produites se propagent sous forme d'\cle{ondes} (ondes P, ondes S, ondes de surface) jusqu'à la surface, où le point situé à la verticale du foyer est l'\cle{épicentre}. Mais une question se pose immédiatement : comment dire qu'un séisme est « fort » ou « faible » ? Comment comparer le séisme de Kobe (Japon, 1995) à celui du Chili (1960) ? Pour répondre à ces questions, les scientifiques utilisent \textbf{deux approches complémentaires} : l'une fondée sur les \textbf{effets observés} sur place (c'est l'\cle{intensité}), l'autre fondée sur l'\textbf{énergie mesurée} par des appareils (c'est la \cle{magnitude}). C'est ce que nous allons étudier maintenant.

\courssub{Deux approches complémentaires : effets observés et énergie mesurée}

Imaginons deux villages touchés par le même séisme. Dans le premier, proche de l'épicentre et construit en torchis (banco), les maisons s'effondrent. Dans le second, situé à \SI{200}{km} et construit en parpaings, on a seulement senti trembler les meubles. Le séisme est pourtant \textbf{le même} ! Il faut donc distinguer :

\begin{itemize}
\item ce que l'on \textbf{observe} en chaque lieu (dégâts, effets sur les personnes et les constructions) : c'est l'\cle{intensité} ;
\item ce que l'on \textbf{mesure} avec des appareils (l'énergie libérée au foyer) : c'est la \cle{magnitude}.
\end{itemize}

\begin{definition}[Intensité et magnitude]
L'\cle{intensité} d'un séisme mesure les \textbf{effets} du séisme en un lieu donné (effets sur les personnes, les objets, les constructions). Elle s'exprime en \textbf{degrés} sur l'\textbf{échelle de Mercalli} (échelle MSK), de I à XII.\par
La \cle{magnitude} d'un séisme mesure l'\textbf{énergie libérée au foyer}. Elle est calculée à partir de l'amplitude des vibrations enregistrées par un sismographe et s'exprime sur l'\textbf{échelle de Richter}, \textbf{sans unité ni degré}.
\end{definition}

\begin{aretenir}[L'idée essentielle]
Un même séisme possède \textbf{une seule magnitude} (une seule énergie libérée au foyer), mais des \textbf{intensités variables} selon les lieux : l'intensité diminue en général quand on s'éloigne de l'épicentre, et elle dépend aussi de la nature des constructions et du sol.
\end{aretenir}

\courssub{L'échelle de Mercalli : évaluer les effets observés}

\textbf{Présentation de l'échelle.}\par
L'échelle de Mercalli (améliorée sous le nom d'échelle \textbf{MSK}) est une échelle \textbf{qualitative} : elle ne repose sur aucun appareil de mesure, mais sur l'\textbf{observation} des effets du séisme. Elle comporte \textbf{12 degrés}, notés en chiffres romains de \textbf{I à XII}. Pour attribuer un degré à un lieu, on interroge la population et on examine les dégâts : les personnes ont-elles senti le séisme ? Les objets ont-ils bougé ? Les bâtiments sont-ils endommagés ou détruits ?

\begin{exemplebox}[Quelques degrés repères de l'échelle de Mercalli]
\begin{center}
\begin{tabularx}{\linewidth}{|c|X|}
\hline
\rowcolor{popblueL} \textbf{Degré} & \textbf{Effets observés} \\
\hline
I & Séisme \textbf{non ressenti} par la population ; détecté seulement par les appareils. \\
\hline
IV & Séisme ressenti par beaucoup de personnes ; la \textbf{vaisselle s'entrechoque}, les portes et fenêtres vibrent. \\
\hline
VII & Dégâts dans les constructions ordinaires : \textbf{cheminées brisées}, fissures dans les murs, chutes de plâtre. \\
\hline
X & \textbf{Destructions générales} : nombreux bâtiments effondrés, rails tordus, glissements de terrain. \\
\hline
XII & \textbf{Catastrophe totale} : tout est détruit, le paysage est modifié, des ondulations sont visibles à la surface du sol. \\
\hline
\end{tabularx}
\end{center}
\end{exemplebox}

\begin{attention}[Ne confonds pas degrés et magnitude !]
Erreur très fréquente au BEPC : écrire « un séisme de 7 degrés sur l'échelle de Richter ». C'est \textbf{faux} ! Retiens bien :
\begin{itemize}
\item l'échelle de \textbf{Mercalli} s'exprime en \textbf{degrés} (I à XII) et décrit les \textbf{effets} ;
\item l'échelle de \textbf{Richter} s'exprime en \textbf{magnitude}, \textbf{sans degré ni unité}, et mesure l'\textbf{énergie}.
\end{itemize}
On dit donc : « un séisme de magnitude 7 sur l'échelle de Richter » ou « un séisme d'intensité VII sur l'échelle de Mercalli ».
\end{attention}

\courssub{Le sismographe et le sismogramme : mesurer les vibrations du sol}

\textbf{Le sismographe : un appareil enregistreur.}\par
Pour mesurer objectivement un séisme, on utilise un appareil appelé \cle{sismographe}.

\begin{definition}[Sismographe et sismogramme]
Le \cle{sismographe} est un appareil qui \textbf{enregistre les vibrations du sol} au passage des ondes sismiques. Le document obtenu, qui traduit les vibrations du sol en fonction du temps, s'appelle un \cle{sismogramme}.
\end{definition}

Le principe est simple : une masse lourde suspendue reste presque immobile par inertie pendant que le sol (et le support de l'appareil) vibre ; un stylet solidaire de la masse trace sur un cylindre tournant (ou un capteur électronique enregistre) les mouvements relatifs du sol. Le tracé obtenu est d'autant plus ample que les vibrations sont fortes.

\textbf{Lecture d'un sismogramme.}\par

\begin{experience}[Observation et interprétation d'un sismogramme]
\textbf{Document :} on observe le sismogramme enregistré par une station après un séisme lointain (voir figure ci-dessous).\par
\textbf{Observations :}
\begin{enumerate}
\item Au début, le tracé est presque plat : le sol ne vibre pas encore.
\item Apparaissent d'abord des vibrations de \textbf{faible amplitude} : ce sont les \textbf{ondes P}, les plus rapides.
\item Arrivent ensuite des vibrations de \textbf{plus forte amplitude} : ce sont les \textbf{ondes S}, plus lentes.
\item Enfin, les vibrations les \textbf{plus amples} correspondent aux \textbf{ondes de surface}, responsables de l'essentiel des dégâts.
\end{enumerate}
\textbf{Interprétation :} les ondes ne se propagent pas à la même vitesse (les ondes P sont les plus rapides), donc elles n'arrivent pas en même temps à la station. L'écart de temps entre l'arrivée des ondes S et des ondes P, noté \textbf{écart S--P}, est d'autant plus grand que la station est \textbf{éloignée de l'épicentre}.
\end{experience}

\begin{popfigure}
\begin{center}
\begin{tikzpicture}
\begin{axis}[
width=0.95\linewidth, height=6.2cm,
xlabel={Temps (min)},
ylabel={Amplitude des vibrations},
xmin=0, xmax=10, ymin=-3.2, ymax=3.2,
xtick={1.5,3.5,6},
xticklabels={arrivée des ondes P, arrivée des ondes S, ondes de surface},
x tick label style={font=\small, align=center},
ytick=\empty,
axis lines=left,
]
\addplot[popcurve, domain=0:1.5, samples=60] {0.05*sin(deg(20*x))};
\addplot[popcurve, domain=1.5:3.5, samples=120] {0.5*sin(deg(25*x))};
\addplot[popcurve, domain=3.5:6, samples=150] {1.2*sin(deg(22*x))};
\addplot[popcurve, domain=6:10, samples=200] {2.6*sin(deg(9*x))*exp(-0.25*(x-6))};
\draw[dashed, popmuted] (axis cs:1.5,-3.2) -- (axis cs:1.5,3.2);
\draw[dashed, popmuted] (axis cs:3.5,-3.2) -- (axis cs:3.5,3.2);
\draw[<->, poporange, thick] (axis cs:1.5,2.6) -- (axis cs:3.5,2.6) node[midway, above, font=\small\bfseries]{écart S--P};
\end{axis}
\end{tikzpicture}
\end{center}
\legende{Sismogramme simplifié enregistré par une station sismique. On distingue l'arrivée successive des ondes P (faible amplitude), des ondes S (amplitude moyenne), puis des ondes de surface (amplitude maximale). L'écart temporel S--P augmente avec la distance de la station à l'épicentre.}
\end{popfigure}

\begin{methode}[Exploiter un sismogramme]
Pour exploiter un sismogramme au BEPC, procède par étapes :
\begin{enumerate}
\item \textbf{Repérer} sur le tracé l'arrivée des ondes P (premières vibrations, faible amplitude) puis l'arrivée des ondes S (vibrations plus amples).
\item \textbf{Mesurer} l'écart temporel S--P (en secondes ou en minutes) : il est \textbf{proportionnel à la distance} entre la station et l'épicentre. Plus l'écart est grand, plus la station est loin.
\item \textbf{Mesurer} l'amplitude maximale du tracé (hauteur de la plus grande vibration) : elle sert au calcul de la magnitude.
\item \textbf{Comparer} plusieurs stations si nécessaire : la station où l'écart S--P est le plus petit est la plus proche de l'épicentre.
\end{enumerate}
\end{methode}

\courssub{L'échelle de Richter : mesurer l'énergie libérée}

L'échelle de Richter est une échelle \textbf{quantitative} : elle repose sur la \textbf{mesure de l'amplitude maximale} enregistrée sur les sismogrammes. La magnitude s'exprime par un nombre, en général \textbf{de 1 à 9 et plus}, sans unité.

\begin{propriete}[Magnitude et énergie]
Sur l'échelle de Richter, chaque \textbf{unité supplémentaire} de magnitude correspond à une énergie libérée environ \textbf{30 fois plus grande}. Ainsi, un séisme de magnitude 6 libère environ 30 fois plus d'énergie qu'un séisme de magnitude 5, et environ $30 \times 30 = 900$ fois plus d'énergie qu'un séisme de magnitude 4.
\end{propriete}

\begin{exemplebox}[Quelques séismes célèbres et leur magnitude]
\begin{itemize}
\item Le séisme de \textbf{Kobe} (Japon, 1995) avait une magnitude de \textbf{6,9} : il a détruit une grande partie de la ville et fait plus de \num{5000} morts.
\item Le séisme du \textbf{Chili} (1960), le plus fort jamais enregistré, a atteint une magnitude de \textbf{9,5}.
\item Le séisme de \textbf{Haïti} (2010), de magnitude \textbf{7,0}, a montré qu'une magnitude « moyenne » peut être catastrophique lorsque les constructions sont fragiles.
\end{itemize}
\end{exemplebox}

\begin{savaistu}[Des centaines de milliers de séismes chaque année]
Les séismes de magnitude \textbf{inférieure à 3} sont généralement \textbf{imperceptibles} : seuls les sismographes les détectent. Il en survient des \textbf{centaines de milliers chaque année} dans le monde ! Heureusement, les séismes destructeurs (magnitude supérieure à 7) sont beaucoup plus rares : quelques-uns par an seulement. Au Cameroun, la région du Mont Cameroun, volcan actif, connaît parfois de petites secousses enregistrées par les sismographes.
\end{savaistu}

\courssub{Distinction fondamentale : une seule magnitude, des intensités variables}

Reprenons notre exemple des deux villages. Le séisme libère au foyer une quantité d'énergie déterminée : sa \textbf{magnitude est unique}, par exemple 6,0. Mais les effets observés varient d'un lieu à l'autre :

\begin{itemize}
\item près de l'épicentre, les vibrations sont fortes, les dégâts importants : l'\textbf{intensité est élevée} (par exemple VIII) ;
\item loin de l'épicentre, les vibrations sont affaiblies : l'\textbf{intensité est plus faible} (par exemple IV) ;
\item à distance égale, une ville aux constructions fragiles (torchis, bâtiments non renforcés) subit des dégâts plus graves qu'une ville aux constructions parasismiques : l'\textbf{intensité dépend aussi de la nature des constructions et du sol}.
\end{itemize}

\begin{aretenir}[À retenir absolument]
Un même séisme a \textbf{une seule magnitude} (une seule énergie libérée au foyer) mais des \textbf{intensités différentes} selon les lieux. L'intensité dépend de la \textbf{distance à l'épicentre} et de la \textbf{nature des constructions et du sol}.
\end{aretenir}

Le tableau ci-dessous résume la comparaison entre les deux échelles.

\begin{center}
\begin{tabularx}{\linewidth}{|l|X|X|}
\hline
\rowcolor{popblueL} \textbf{Critère} & \textbf{Échelle de Mercalli (MSK)} & \textbf{Échelle de Richter} \\
\hline
Grandeur évaluée & Intensité : les \textbf{effets} observés en un lieu & Magnitude : l'\textbf{énergie} libérée au foyer \\
\hline
Nature de l'échelle & \textbf{Qualitative} (observations, enquêtes) & \textbf{Quantitative} (mesures sur sismogrammes) \\
\hline
Expression & En \textbf{degrés}, de I à XII & En \textbf{magnitude}, de 1 à 9 et plus, sans unité \\
\hline
Valeur pour un même séisme & \textbf{Variable} selon les lieux & \textbf{Unique} \\
\hline
Exemples & I : non ressenti ; VII : cheminées brisées ; XII : catastrophe totale & Kobe 1995 : 6,9 ; Chili 1960 : 9,5 \\
\hline
\end{tabularx}
\end{center}

\courssub{Exploitation de données : comparer des sismogrammes}

\begin{exoresolu}[Quelle station est la plus proche de l'épicentre ?]
\textbf{Énoncé.} Un séisme est enregistré par trois stations A, B et C. On mesure sur chaque sismogramme l'écart temporel entre l'arrivée des ondes P et des ondes S :
\begin{center}
\begin{tabular}{|l|c|c|c|}
\hline
\rowcolor{popblueL} \textbf{Station} & A & B & C \\
\hline
Écart S--P (s) & 20 & 55 & 35 \\
\hline
\end{tabular}
\end{center}
1. Rappeler ce que mesure l'écart S--P.\par
2. Classer les stations de la plus proche à la plus éloignée de l'épicentre. Justifier.\par
3. Ces trois stations donneront-elles la même magnitude pour ce séisme ? Même intensité ? Justifier.
\tcblower
\textbf{Solution détaillée.}\par
1. L'écart S--P est le temps qui sépare l'arrivée des ondes S de celle des ondes P sur le sismogramme. Comme les ondes P sont plus rapides que les ondes S, cet écart est \textbf{proportionnel à la distance} entre la station et l'épicentre.\par
2. Plus l'écart S--P est petit, plus la station est proche de l'épicentre. Or $20 < 35 < 55$. Le classement de la plus proche à la plus éloignée est donc : \textbf{A, puis C, puis B}.\par
3. Les trois stations trouveront (après calcul) la \textbf{même magnitude}, car la magnitude mesure l'énergie libérée au foyer : elle est \textbf{unique} pour un séisme donné. En revanche, les \textbf{intensités seront différentes} : la station A, la plus proche de l'épicentre, observera les effets les plus importants (intensité la plus élevée), et la station B, la plus éloignée, les effets les plus faibles.
\end{exoresolu}

\begin{exercice}[Pour t'entraîner]
Un séisme de magnitude 5,2 est ressenti dans deux villes : Douala, proche de l'épicentre, et Bafoussam, plus éloignée.\par
1. La magnitude du séisme est-elle différente dans les deux villes ? Justifier.\par
2. Dans quelle ville l'intensité (échelle de Mercalli) est-elle a priori la plus élevée ? Justifier.\par
3. Citer un autre facteur que la distance qui peut faire varier l'intensité entre deux villes situées à la même distance de l'épicentre.
\end{exercice}

\begin{corrige}[Exercice]
1. Non : la magnitude est \textbf{unique} pour un séisme donné, car elle mesure l'énergie libérée au foyer. Elle vaut 5,2 dans les deux villes.\par
2. C'est à \textbf{Douala}, la ville la plus proche de l'épicentre, que l'intensité est a priori la plus élevée, car l'intensité diminue en général lorsque la distance à l'épicentre augmente.\par
3. L'autre facteur est la \textbf{nature des constructions et du sol} : des bâtiments fragiles ou un sol meuble amplifient les dégâts, donc l'intensité observée.
\end{corrige}

\begin{attention}[Récapitulatif des pièges du BEPC]
\begin{itemize}
\item Ne dis jamais « degrés Richter » : la magnitude de Richter s'exprime \textbf{sans degré ni unité}.
\item Ne confonds pas sismographe (l'\textbf{appareil}) et sismogramme (le \textbf{document enregistré}).
\item N'affirme pas qu'un séisme a « plusieurs magnitudes » : il a \textbf{une seule magnitude} et \textbf{plusieurs intensités}.
\item Sur un sismogramme, ce sont les ondes \textbf{P} qui arrivent en premier, pas les ondes de surface.
\end{itemize}
\end{attention}

En résumé, l'évaluation d'un séisme repose sur deux outils complémentaires : l'échelle de Mercalli, qui décrit les effets observés en degrés, et l'échelle de Richter, qui mesure l'énergie libérée grâce aux sismographes. Mais où ces séismes se produisent-ils sur la planète ? Sont-ils répartis au hasard ? C'est ce que nous découvrirons dans la section suivante, consacrée à la répartition mondiale des séismes.

\coursec{Répartition mondiale des séismes}

Dans la section précédente, nous avons appris à mesurer un séisme : son intensité (effets ressentis, échelle de Mercalli) et sa magnitude (énergie libérée, échelle de Richter). Mais une question demeure : les séismes se produisent-ils n'importe où sur la planète, au hasard ? Pour le savoir, il suffit de reporter sur une carte du monde l'emplacement de tous les épicentres enregistrés par les sismographes pendant plusieurs années. Le résultat est spectaculaire — et c'est l'une des plus belles découvertes des sciences de la Terre.

\courssub{Observation d'une carte mondiale : une répartition qui n'est pas due au hasard}

\begin{experience}[Observation d'une carte mondiale de répartition des épicentres]
\textbf{Document :} on dispose d'une carte du monde sur laquelle chaque point représente l'épicentre d'un séisme enregistré au cours des dernières décennies (plusieurs dizaines de milliers de points).

\textbf{Protocole d'observation :}
\begin{enumerate}
\item Décrire la distribution globale des points : sont-ils dispersés uniformément ou regroupés ?
\item Repérer les zones où les points sont très nombreux et celles où ils sont absents.
\item Comparer cette carte avec une carte des plaques lithosphériques.
\end{enumerate}

\textbf{Observations :}
\begin{itemize}
\item Les épicentres ne sont \textbf{pas répartis au hasard} : ils forment des \cle{alignements} étroits et allongés, comme des « coutures » à la surface du globe.
\item De vastes régions (intérieur de l'Afrique centrale, de l'Australie, du Brésil, du Canada) sont presque \textbf{dépourvues} de séismes.
\item Les alignements d'épicentres coïncident avec les \textbf{limites des plaques lithosphériques}.
\end{itemize}

\textbf{Interprétation :} la répartition des séismes est \textbf{localisée} et \textbf{significative} : elle dessine le découpage de la lithosphère en plaques.
\end{experience}

\begin{attention}[Piège fréquent]
Ne jamais écrire que « les séismes ont lieu partout sur la Terre » ou qu'ils « frappent au hasard ». C'est faux : leur répartition est \textbf{localisée} en zones étroites et allongées, et cette localisation a une \textbf{signification} : elle révèle les limites des plaques lithosphériques. L'intérieur des plaques est, lui, relativement stable.
\end{attention}

\courssub{Les grandes zones sismiques du globe}

L'observation détaillée de la carte permet de distinguer quatre grandes catégories de zones sismiques.

\textbf{a) La ceinture circum-pacifique, ou « ceinture de feu »}

C'est la zone sismique la plus active du monde. Elle fait le tour de l'océan Pacifique : elle longe les côtes du \textbf{Chili}, du Pérou, de l'Amérique centrale et de la \textbf{Californie} d'un côté, et celles du \textbf{Japon}, des Philippines et de l'\textbf{Indonésie} de l'autre. À elle seule, elle concentre environ \textbf{80 \% de l'énergie sismique mondiale}. C'est là que se produisent les séismes les plus violents et les plus profonds (exemple : le séisme du Japon de 2011, de magnitude 9,0).

\textbf{b) La ceinture alpine-himalayenne}

Deuxième grande zone sismique, elle s'étend de la \textbf{Méditerranée} (Italie, Grèce, \textbf{Turquie}) jusqu'à l'Indonésie, en passant par l'\textbf{Iran} et la chaîne de l'\textbf{Himalaya}. Elle correspond au rapprochement entre les plaques africaine, arabe et indienne d'une part, et la plaque eurasiatique d'autre part. Le séisme de Turquie de 2023 (magnitude 7,8) en est une tragique illustration.

\textbf{c) Les dorsales océaniques}

Au milieu des océans, on observe des alignements de séismes de \textbf{faible magnitude} et de \textbf{faible profondeur}, qui suivent exactement le tracé des \cle{dorsales océaniques} (chaînes de montagnes sous-marines), comme la dorsale médio-atlantique qui coupe l'océan Atlantique en deux, de l'Islande jusqu'au sud de l'Afrique.

\textbf{d) Les rifts continentaux}

Certains continents sont traversés par des \cle{rifts} : des fossés d'effondrement bordés de failles, où la croûte terrestre s'amincit et s'écarte. Le \textbf{rift est-africain} (Éthiopie, Kenya, Tanzanie) en est l'exemple type : il est marqué par des séismes fréquents et par du volcanisme.

\begin{popfigure}
\begin{center}
\begin{tikzpicture}[scale=0.72]
% Océan de fond
\fill[popblueL] (-6,-3.2) rectangle (6,3.2);
% Continents simplifiés
\fill[popgreenL,rounded corners=2pt] (-5.6,0.6) -- (-4.2,1.6) -- (-3.0,1.4) -- (-2.8,0.2) -- (-3.6,-0.4) -- (-4.6,-0.2) -- cycle; % Amérique N
\fill[popgreenL,rounded corners=2pt] (-3.4,-0.5) -- (-2.4,-0.6) -- (-2.2,-2.6) -- (-3.0,-2.8) -- (-3.4,-1.6) -- cycle; % Amérique S
\fill[popgreenL,rounded corners=2pt] (-0.6,0.8) -- (0.8,1.6) -- (2.4,1.4) -- (2.8,0.4) -- (1.6,0.0) -- (0.2,0.2) -- cycle; % Eurasie
\fill[popgreenL,rounded corners=2pt] (0.2,0.1) -- (1.4,0.0) -- (1.6,-1.8) -- (0.8,-2.4) -- (0.2,-1.2) -- cycle; % Afrique
\fill[popgreenL,rounded corners=2pt] (2.6,0.3) -- (4.6,0.8) -- (5.4,0.2) -- (4.4,-0.6) -- (3.0,-0.4) -- cycle; % Asie SE
\fill[popgreenL,rounded corners=2pt] (4.0,-1.4) -- (5.0,-1.2) -- (5.2,-2.0) -- (4.2,-2.2) -- cycle; % Australie
% Ceinture circum-pacifique (points rouges)
\foreach \p in {(-5.3,1.2),(-4.9,0.9),(-4.5,0.6),(-4.1,0.35),(-3.7,0.1),(-3.3,-0.3),(-2.9,-0.7),(-2.6,-1.2),(-2.45,-1.8),(-2.4,-2.4)} {\fill[popink] \p circle (1.6pt);}
\foreach \p in {(5.5,1.4),(5.2,1.0),(4.9,0.6),(4.7,0.1),(4.5,-0.5),(4.3,-1.0),(5.6,0.5),(5.4,-0.3)} {\fill[popink] \p circle (1.6pt);}
% Ceinture alpine-himalayenne (points orange)
\foreach \p in {(-0.2,0.55),(0.4,0.5),(1.0,0.45),(1.6,0.4),(2.2,0.35),(2.8,0.3),(3.4,0.25),(4.0,0.15)} {\fill[poporange] \p circle (1.6pt);}
% Dorsale médio-atlantique (points bleus)
\foreach \p in {(-1.6,1.6),(-1.5,1.1),(-1.55,0.6),(-1.5,0.1),(-1.45,-0.4),(-1.5,-0.9),(-1.55,-1.4),(-1.5,-1.9),(-1.45,-2.4)} {\fill[popblue] \p circle (1.6pt);}
% Rift est-africain (points violets)
\foreach \p in {(1.35,-0.3),(1.4,-0.7),(1.45,-1.1),(1.5,-1.5)} {\fill[poppurple] \p circle (1.6pt);}
% Légende
\node[anchor=west,font=\small] at (-5.9,-3.7) {\textcolor{popink}{$\bullet$} Ceinture circum-pacifique (\og ceinture de feu \fg)};
\node[anchor=west,font=\small] at (-0.4,-3.7) {\textcolor{poporange}{$\bullet$} Ceinture alpine-himalayenne};
\node[anchor=west,font=\small] at (-5.9,-4.2) {\textcolor{popblue}{$\bullet$} Dorsale médio-atlantique};
\node[anchor=west,font=\small] at (-0.4,-4.2) {\textcolor{poppurple}{$\bullet$} Rift est-africain};
\end{tikzpicture}
\end{center}
\legende{Carte simplifiée de la répartition mondiale des épicentres. Les points ne sont pas dispersés au hasard : ils forment des alignements étroits qui dessinent les limites des plaques lithosphériques (ceinture circum-pacifique, ceinture alpine-himalayenne, dorsales océaniques, rift est-africain).}
\end{popfigure}

\courssub{Interprétation : les séismes marquent les limites des plaques lithosphériques}

Superposons maintenant la carte des épicentres et la carte des \cle{plaques lithosphériques} (les grands morceaux rigides qui constituent l'enveloppe externe de la Terre et qui se déplacent lentement, de quelques centimètres par an). Le constat est frappant : les alignements sismiques coïncident presque exactement avec les \textbf{frontières des plaques}.

\begin{propriete}[Répartition des séismes et plaques lithosphériques]
Les zones sismiques coïncident avec les \textbf{limites des plaques lithosphériques}. L'intérieur des plaques est relativement \textbf{stable} (peu ou pas de séismes). Les séismes traduisent donc les \textbf{mouvements relatifs des plaques} : rapprochement, écartement ou coulissage.
\end{propriete}

On distingue trois types de limites de plaques, chacune associée à un type de séisme :

\begin{itemize}
\item \textbf{Limite divergente} (dorsales, rifts) : deux plaques \textbf{s'écartent}. Les roches se fissurent et se déchirent, provoquant des séismes \textbf{fréquents mais peu violents} et superficiels.
\item \textbf{Limite convergente} (ceintures circum-pacifique et alpine-himalayenne) : deux plaques \textbf{se rapprochent}. L'une peut plonger sous l'autre (\cle{subduction}) ou les deux se compriment et se plissent (Himalaya). Les contraintes accumulées se libèrent brutalement : séismes \textbf{très violents}, parfois très profonds.
\item \textbf{Limite transformante} (faille de San Andreas) : deux plaques \textbf{coulissent} horizontalement l'une contre l'autre, sans créer ni détruire de croûte. Le frottement bloque puis relâche brutalement les roches : séismes \textbf{fréquents et parfois destructeurs}, toujours superficiels.
\end{itemize}

\begin{popfigure}
\begin{center}
\begin{tikzpicture}[scale=0.9]
% --- Limite divergente ---
\begin{scope}[shift={(-4.6,0)}]
\fill[popgreenL] (-1.6,0) rectangle (-0.15,0.5);
\fill[popgreenL] (0.15,0) rectangle (1.6,0.5);
\fill[poporangeL] (-0.15,-0.1) rectangle (0.15,0.5);
\draw[->,thick,popdark] (-0.8,0.7) -- (-1.5,0.7);
\draw[->,thick,popdark] (0.8,0.7) -- (1.5,0.7);
\draw[popink,thick] (0,0.5) -- (0,0.95);
\draw[popink] (0,1.05) node[font=\footnotesize] {$\bigstar$};
\node[font=\small\bfseries,align=center] at (0,-0.5) {Divergente\\ (dorsale, rift)};
\node[font=\footnotesize,align=center,text=popmuted] at (0,-1.15) {séismes superficiels,\\ peu violents};
\end{scope}
% --- Limite convergente ---
\begin{scope}[shift={(0,0)}]
\fill[popgreenL] (-1.6,0) -- (-0.1,0) -- (-0.1,-0.7) -- (-1.6,-0.7) -- cycle;
\fill[popblueL] (1.6,0) -- (0.1,0) -- (-0.35,-0.7) -- (1.6,-0.7) -- cycle;
\draw[->,thick,popdark] (-1.5,0.35) -- (-0.5,0.35);
\draw[->,thick,popdark] (1.5,0.35) -- (0.5,0.35);
\draw[popink] (-0.5,-0.45) node[font=\footnotesize] {$\bigstar$};
\draw[popink] (-0.25,-0.15) node[font=\footnotesize] {$\bigstar$};
\node[font=\small\bfseries,align=center] at (0,-1.2) {Convergente\\ (subduction, collision)};
\node[font=\footnotesize,align=center,text=popmuted] at (0,-1.85) {séismes violents,\\ parfois profonds};
\end{scope}
% --- Limite transformante ---
\begin{scope}[shift={(4.6,0)}]
\fill[popgreenL] (-1.6,0) rectangle (0,0.5);
\fill[popblueL] (0,0) rectangle (1.6,0.5);
\draw[thick,popink,dashed] (0,0) -- (0,0.5);
\draw[->,thick,popdark] (-0.8,0.62) -- (-1.5,0.62);
\draw[->,thick,popdark] (0.8,-0.12) -- (1.5,-0.12);
\draw[popink] (0,0.25) node[font=\footnotesize] {$\bigstar$};
\node[font=\small\bfseries,align=center] at (0,-0.5) {Transformante\\ (faille de San Andreas)};
\node[font=\footnotesize,align=center,text=popmuted] at (0,-1.15) {séismes superficiels,\\ fréquents};
\end{scope}
\end{tikzpicture}
\end{center}
\legende{Les trois types de limites de plaques et les séismes associés. Les flèches noires indiquent le sens de déplacement des plaques ; les étoiles rouges figurent les foyers sismiques. À gauche : écartement (divergence). Au centre : rapprochement (convergence, avec subduction d'une plaque sous l'autre). À droite : coulissage horizontal (limite transformante).}
\end{popfigure}

\begin{exemplebox}[La faille de San Andreas (Californie)]
La faille de San Andreas, longue d'environ \SI{1300}{km}, est une limite de plaque \textbf{transformante} : la plaque Pacifique et la plaque Nord-Américaine \textbf{coulissent} l'une contre l'autre à raison de quelques centimètres par an. Lorsque le frottement bloque les roches, la contrainte s'accumule pendant des années, puis se libère brutalement : c'est un séisme. La Californie connaît ainsi des séismes fréquents (San Francisco, 1906 ; Loma Prieta, 1989), ce qui explique les normes de construction parasismiques très strictes de cette région.
\end{exemplebox}

\begin{savaistu}[Des séismes à 700 km de profondeur]
Les séismes les plus profonds de la planète — jusqu'à environ \SI{700}{km} de profondeur — se produisent dans les \textbf{zones de subduction}, là où une plaque océanique plonge sous une autre plaque et s'enfonce dans le manteau. En s'enfonçant, la plaque se déforme et se rompt : les foyers sismiques dessinent alors un plan incliné, véritable « photographie » de la plaque plongeante. C'est une preuve directe de l'existence de la subduction.
\end{savaistu}

\begin{methode}[Exploiter une carte de répartition des séismes]
Face à une carte mondiale des épicentres (document fréquent au BEPC), procède toujours en trois étapes :
\begin{enumerate}
\item \textbf{Décrire} la distribution : les points forment-ils des alignements ? Quelles zones sont vides ? (Ne conclus pas encore !)
\item \textbf{Comparer} avec une carte des plaques lithosphériques : faire correspondre les alignements sismiques avec les limites de plaques.
\item \textbf{Conclure} : les séismes se produisent aux limites des plaques et traduisent leurs mouvements relatifs (écartement, rapprochement, coulissage) ; l'intérieur des plaques est stable.
\end{enumerate}
\end{methode}

\courssub{Et le Cameroun dans tout cela ?}

Le Cameroun n'est pas situé sur une limite de plaque active : il se trouve à l'intérieur de la plaque africaine. Pourtant, notre pays connaît une \textbf{sismicité modérée}, qui s'explique par deux facteurs :

\begin{itemize}
\item La \textbf{ligne volcanique du Cameroun}, alignement volcanique qui s'étend du golfe de Guinée (île de Bioko) au Mont Cameroun (\SI{4100}{m}), volcan encore actif (éruptions en 1999 et 2000). L'activité volcanique s'accompagne de séismes de faible à moyenne magnitude.
\item Des \textbf{failles anciennes réactivées} : de vieilles cassures de la croûte terrestre, héritées de l'histoire géologique de l'Afrique, peuvent rejouer faiblement sous l'effet des contraintes actuelles, provoquant des secousses ressenties par la population (par exemple dans l'ouest du pays).
\end{itemize}

Ainsi, même loin des limites de plaques, quelques séismes peuvent survenir, mais ils restent \textbf{moins fréquents et moins violents} que ceux de la ceinture de feu.

\begin{exoresolu}[Exploiter une carte de répartition]
\textbf{Énoncé.} Sur une carte du monde, on a reporté les épicentres des séismes enregistrés entre 2000 et 2020. On constate que : (1) les épicentres forment un anneau presque continu autour de l'océan Pacifique ; (2) un alignement de points suit le milieu de l'océan Atlantique ; (3) le centre de l'Afrique est presque dépourvu de séismes. Décrire puis interpréter ces observations.
\tcblower
\textbf{Solution détaillée.}

\emph{Étape 1 — Description.} Les épicentres ne sont pas répartis au hasard : ils forment des alignements étroits et allongés (anneau circum-pacifique, alignement médio-atlantique), tandis que de vastes régions (centre de l'Afrique) en sont vides.

\emph{Étape 2 — Comparaison.} En superposant la carte des plaques lithosphériques, on constate que l'anneau circum-pacifique correspond aux limites convergentes autour de la plaque Pacifique (subductions), et que l'alignement atlantique suit la dorsale médio-atlantique, limite divergente entre les plaques américaines d'une part, africaine et eurasiatique d'autre part. Le centre de l'Afrique se situe à l'intérieur de la plaque africaine.

\emph{Étape 3 — Conclusion.} Les séismes se produisent aux \textbf{limites des plaques lithosphériques} et traduisent leurs mouvements relatifs : rapprochement (ceinture circum-pacifique) ou écartement (dorsale atlantique). L'intérieur des plaques, comme le centre de l'Afrique, est relativement stable.
\end{exoresolu}

\begin{aretenir}[L'essentiel de la section]
\begin{itemize}
\item Les séismes ne sont \textbf{pas répartis au hasard} : leurs épicentres forment des alignements étroits et allongés.
\item Les grandes zones sismiques sont : la \textbf{ceinture circum-pacifique} (environ 80 \% de l'énergie sismique mondiale), la \textbf{ceinture alpine-himalayenne}, les \textbf{dorsales océaniques} et les \textbf{rifts continentaux}.
\item Les zones sismiques coïncident avec les \textbf{limites des plaques lithosphériques} ; l'intérieur des plaques est relativement stable.
\item Les séismes traduisent les \textbf{mouvements relatifs des plaques} : écartement (divergence), rapprochement (convergence, subduction) ou coulissage (limite transformante, ex. faille de San Andreas).
\item Le Cameroun connaît une sismicité \textbf{modérée}, liée à la ligne volcanique du Cameroun et à des failles anciennes réactivées.
\end{itemize}
\end{aretenir}

\coursec{Prévision et prévention des séismes}

Nous avons vu que les séismes se concentrent dans des zones bien précises du globe (limites de plaques, grandes failles). Une question se pose alors naturellement : \textbf{peut-on annoncer un séisme à l'avance pour sauver des vies ?} C'est tout l'enjeu de la \cle{prévision} et de la \cle{prévention} des séismes. Nous allons voir que la prévision exacte reste impossible, mais que la prévention, elle, permet de sauver des milliers de vies.

\courssub{La prévision des séismes : un défi non résolu}

\begin{definition}[Prévision sismique]
La \cle{prévision sismique} consiste à annoncer \textbf{à l'avance} le \textbf{lieu}, la \textbf{date} et la \textbf{magnitude} d'un séisme futur, avec assez de certitude pour déclencher une alerte ou une évacuation.
\end{definition}

\begin{aretenir}[Savoir admis : l'impossibilité actuelle]
À ce jour, \textbf{aucune méthode ne permet de prévoir de manière fiable la date exacte d'un séisme}. Les sismologues savent seulement :
\begin{itemize}
    \item identifier les \textbf{zones à risque} (les régions proches des limites de plaques et des failles actives) ;
    \item estimer des \textbf{probabilités} sur de longues durées (par exemple : « forte probabilité d'un grand séisme sur telle faille dans les 30 prochaines années »).
\end{itemize}
Mais ils ne peuvent pas dire : « Demain à 14 h, un séisme de magnitude 7 frappera telle ville. »
\end{aretenir}

\courssub{Les signes précurseurs : des pistes étudiées, pas des certitudes}

Depuis longtemps, les scientifiques recherchent des \cle{signes précurseurs}, c'est-à-dire des phénomènes anormaux qui annonceraient un séisme. Certains ont été observés avant certains séismes, mais \textbf{aucun n'est assez régulier et précis pour servir d'alerte fiable}.

\begin{experience}[Un cas célèbre : le séisme de Haicheng (1975)]
\textbf{Observation :} En février 1975, à Haicheng (Chine), les autorités font évacuer la ville quelques heures avant un séisme de magnitude 7,3. Des dizaines de milliers de vies sont sauvées.
\textbf{Signes observés avant le séisme :}
\begin{itemize}
    \item de nombreux \cle{microséismes} (petites secousses) pendant les jours précédents ;
    \item des variations anormales du niveau de l'eau dans les puits ;
    \item des comportements étranges des animaux (serpents sortant de leur terrier en plein hiver, bétail agité).
\end{itemize}
\textbf{Interprétation :} la réunion de ces signes a conduit les scientifiques chinois à émettre une alerte.
\textbf{Conclusion :} c'est le \textbf{seul cas connu d'évacuation réussie} avant un grand séisme. Mais ce succès n'a \textbf{jamais été reproduit} : en 1976, le séisme de Tangshan (magnitude 7,8, environ 240 000 morts) a frappé \textbf{sans aucun signe précurseur détecté}.
\end{experience}

\begin{center}
\begin{tabularx}{\linewidth}{|l|X|X|}\hline
\rowcolor{popblueL}
\textbf{Signe précurseur étudié} & \textbf{Observation} & \textbf{Fiabilité} \\ \hline
Microséismes & Petites secousses précédant parfois le choc principal. & \textbf{Faible} : beaucoup de grands séismes n'en ont pas ; on ne peut pas les distinguer d'une activité normale. \\ \hline
Niveau de l'eau des puits & La déformation des roches peut modifier le niveau des nappes d'eau souterraines. & \textbf{Faible} : le niveau varie aussi avec les pluies et le pompage. \\ \hline
Déformations du sol & Le sol se déforme lentement quand les contraintes s'accumulent (mesures GPS). & \textbf{Moyenne} : montre \emph{où} le danger s'accumule, mais pas \emph{quand} la rupture aura lieu. \\ \hline
Comportement des animaux & Agitation anormale de certains animaux avant certains séismes. & \textbf{Nulle} : observations non mesurables, trop de fausses alertes. \\ \hline
\end{tabularx}
\end{center}

\begin{attention}[Piège : les animaux « prévisionnistes »]
Les médias rapportent souvent des comportements bizarres d'animaux avant un séisme. \textbf{Attention :} il s'agit d'anecdotes, pas de science. Aucune étude rigoureuse n'a montré que les animaux permettent de prévoir un séisme. Se fier à eux pour évacuer une population serait dangereux (fausses alertes, panique inutile).
\end{attention}

\courssub{Localiser les zones à risque : la seule « prévision » possible}

Puisqu'on ne peut pas prévoir \emph{quand} un séisme aura lieu, les scientifiques cherchent à déterminer \emph{où} le risque est le plus grand.

\begin{definition}[Faille active]
Une \cle{faille active} est une faille qui a déjà produit des séismes dans les temps géologiques récents et qui peut en produire à nouveau. Le long d'une faille active, un segment qui n'a \textbf{pas rompu depuis longtemps} alors que les segments voisins ont déjà eu des séismes accumule les contraintes : c'est une zone de \textbf{risque élevé}.
\end{definition}

\begin{exemplebox}[Et au Cameroun ?]
Le Cameroun connaît une sismicité modérée mais réelle, liée notamment à la \textbf{faille de la Sanaga} et à la \textbf{Ligne du Cameroun} (zone volcanique comprenant le \textbf{Mont Cameroun}, volcan actif). Des secousses de magnitude généralement inférieure à 5 y sont enregistrées. Les villes comme Limbé, Buéa ou Douala doivent donc tenir compte de ce risque dans la construction.
\end{exemplebox}

\courssub{La prévention parasismique : trois piliers pour sauver des vies}

Puisque la prévision est impossible, la \cle{prévention} est la seule réponse efficace. Elle repose sur trois piliers.

\begin{methode}[Les trois piliers de la prévention sismique]
\begin{enumerate}
    \item \textbf{Construire parasismique} : adapter les constructions au risque sismique de la région (bâtiments solides et souples).
    \item \textbf{Aménager le territoire} : éviter de construire sur les failles actives, les sols meubles saturés d'eau et les versants instables.
    \item \textbf{Éduquer et préparer la population} : exercices d'évacuation, plans d'urgence, trousse de secours (eau, lampe, radio, médicaments).
\end{enumerate}
\end{methode}

\courssub{Construire parasismique : un bâtiment qui se déforme sans s'effondrer}

L'effondrement des bâtiments est la \textbf{première cause de mort} lors des séismes. Un bâtiment parasismique ne doit pas être trop rigide (il casserait net) : il doit être \textbf{solide mais souple}, capable de se déformer sans s'effondrer.

\begin{popfigure}
\begin{tikzpicture}[scale=0.75, every node/.style={font=\small}]
    % --- BÂTIMENT ORDINAIRE ---
    \begin{scope}[xshift=-6.5cm]
        \fill[popmuted!30] (-2,-0.5) rectangle (2,0);
        \draw[thick] (-2,0) -- (2,0);
        \node[below=2pt] at (0,-0.5) {Sol};
        % Poteaux
        \foreach \x in {-1.5, 1.3} {
            \draw[fill=poporange!20, thick] (\x,0) rectangle (\x+0.2, 4);
        }
        % Poutres
        \foreach \y in {1.3, 2.6, 3.9} {
            \draw[fill=poporange!20, thick] (-1.5,\y) rectangle (1.5,\y+0.2);
        }
        % Murs en maçonnerie
        \foreach \y in {0.3, 1.6, 2.9} {
            \draw[fill=poporange!40] (-1.3,\y) rectangle (1.3,\y+0.8);
        }
        % Fissures
        \draw[very thick, poppink, decorate, decoration={zigzag, segment length=3mm, amplitude=1mm}] (-1.3,0.3) -- (1.3,1.1);
        \draw[very thick, poppink, decorate, decoration={zigzag, segment length=3mm, amplitude=1mm}] (1.3,1.6) -- (-1.3,2.4);
        \node[poppink, align=center, font=\footnotesize] at (0,4.6) {Fissures, effondrement};
        \node[below=8pt, font=\bfseries, align=center] at (0,-0.5) {Bâtiment ordinaire\\(maçonnerie rigide)};
    \end{scope}
    % --- BÂTIMENT PARASISMIQUE ---
    \begin{scope}[xshift=6.5cm]
        \fill[popmuted!30] (-2,-0.5) rectangle (2,0);
        \draw[thick] (-2,0) -- (2,0);
        \node[below=2pt] at (0,-0.5) {Sol};
        % Ossature
        \foreach \x in {-1.5, 1.3} {
            \draw[fill=popblue!10, thick, pattern=north east lines, pattern color=popblue] (\x,0) rectangle (\x+0.2, 4);
        }
        \foreach \y in {1.3, 2.6, 3.9} {
            \draw[fill=popblue!10, thick, pattern=north east lines, pattern color=popblue] (-1.5,\y) rectangle (1.5,\y+0.2);
        }
        % Contreventements (croix)
        \draw[thick, popblue] (-1.5,0) -- (1.5,1.3);
        \draw[thick, popblue] (1.5,0) -- (-1.5,1.3);
        \draw[thick, popblue] (-1.5,1.5) -- (1.5,2.6);
        \draw[thick, popblue] (1.5,1.5) -- (-1.5,2.6);
        \draw[thick, popblue] (-1.5,2.8) -- (1.5,3.9);
        \draw[thick, popblue] (1.5,2.8) -- (-1.5,3.9);
        \node[popblue, align=center, font=\footnotesize] at (0,4.6) {Se déforme sans s'effondrer};
        \node[below=8pt, font=\bfseries, align=center] at (0,-0.5) {Bâtiment parasismique\\(ossature souple renforcée)};
    \end{scope}
\end{tikzpicture}
\legende{Comparaison schématique : à gauche, un bâtiment ordinaire en maçonnerie rigide se fissure et s'effondre. À droite, un bâtiment parasismique possède une ossature en béton armé ou en acier avec des \textbf{contreventements} (croix de renfort) qui lui permettent de se déformer sans s'effondrer.}
\end{popfigure}

\begin{aretenir}[Principes de la construction parasismique]
\begin{itemize}
    \item Des \textbf{fondations profondes} ancrées dans le sol dur.
    \item Une \textbf{ossature en béton armé ou en acier}, solide et souple, avec des \textbf{contreventements} (renforts en croix).
    \item Des matériaux \textbf{légers} pour les toitures.
    \item Des formes \textbf{simples et régulières} (éviter les bâtiments trop hauts et déséquilibrés).
\end{itemize}
\end{aretenir}

\begin{exemplebox}[La magnitude ne tue pas, la vulnérabilité tue]
En 2010, deux séismes frappent l'Amérique : Haïti (magnitude 7,0, environ 250 000 morts) et le Chili (magnitude 8,8, environ 500 morts). Pourtant, le séisme chilien était bien plus puissant ! La différence s'explique par la \textbf{prévention} : au Chili, les normes de construction parasismique sont strictes et appliquées, et la population est entraînée ; en Haïti, les constructions étaient fragiles et la population non préparée. \textbf{Ce n'est pas le séisme qui tue, ce sont les constructions inadaptées.}
\end{exemplebox}

\courssub{Les comportements qui sauvent : pendant et après le séisme}

La prévention, c'est aussi savoir \textbf{réagir en quelques secondes}. Ces gestes s'apprennent et se répètent lors d'exercices.

\begin{center}
\begin{tabularx}{\linewidth}{|l|X|}\hline
\rowcolor{popblueL}
\textbf{Situation} & \textbf{Conduite à tenir} \\ \hline
\textbf{À l'intérieur} (maison, salle de classe) & Se mettre \textbf{sous une table solide} et s'y accrocher, en protégeant sa tête. S'éloigner des \textbf{fenêtres}, des miroirs et des meubles hauts non fixés. \textbf{Ne pas courir} vers la sortie pendant les secousses. \textbf{Ne jamais prendre l'ascenseur}. \\ \hline
\textbf{À l'extérieur} & S'écarter des \textbf{bâtiments} (chute de façades, de vitres), des murs, des poteaux électriques et des arbres. Rejoindre un \textbf{espace dégagé} (place, terrain de sport). \\ \hline
\textbf{Après les secousses} & \textbf{Couper} le gaz, l'électricité et l'eau (risques d'incendie et d'inondation). Sortir calmement par les escaliers vers le point de rassemblement. Se méfier des \textbf{répliques} (secousses qui suivent le séisme principal). \textbf{Ne pas réintégrer} un bâtiment fissuré. Écouter la radio pour les consignes officielles. \\ \hline
\end{tabularx}
\end{center}

\begin{popfigure}
\begin{tikzpicture}[scale=0.9, every node/.style={font=\small, inner sep=2pt}]
    \def\boxw{4.5}
    \def\boxh{3.4}
    \def\gap{0.6}
    \tikzset{
        picbox/.style={draw=popdark, thick, rounded corners=5pt, minimum width=\boxw cm, minimum height=\boxh cm, fill=popcream},
        ictitle/.style={font=\bfseries\footnotesize, text width=4cm, align=center, anchor=north, yshift=-0.2cm},
        icdesc/.style={font=\scriptsize, text width=4cm, align=center, anchor=south, yshift=0.2cm}
    }
    % 1. Sous table
    \node[picbox] (b1) at (0,0) {};
    \node[ictitle] at (b1.north) {À L'INTÉRIEUR};
    \begin{scope}[shift={(b1.center)}]
        \draw[fill=poporange!30, thick] (-1.2,-0.3) rectangle (1.2,0);
        \foreach \x in {-1.1, 1.0} \draw[thick, fill=poporange!30] (\x,-0.3) rectangle (\x+0.1,-1);
        \draw[fill=popblue!20, thick] (0,-1.3) circle (0.25);
        \draw[thick] (0,-1.05) -- (0,-0.4);
        \draw[thick] (-0.35,-0.7) -- (0.35,-0.7);
    \end{scope}
    \node[icdesc] at (b1.south) {Se mettre sous une table solide et s'y accrocher.};
    % 2. Loin fenêtres
    \node[picbox] (b2) at (\boxw+\gap, 0) {};
    \node[ictitle] at (b2.north) {LOIN DES VITRES};
    \begin{scope}[shift={(b2.center)}]
        \draw[thick, fill=popblue!10] (-1.5,0.2) rectangle (-0.3,1.3);
        \draw[thick] (-1.5,0.75) -- (-0.3,0.75);
        \draw[thick] (-0.9,0.2) -- (-0.9,1.3);
        \draw[poppink, thick, decorate, decoration={zigzag, amplitude=1pt, segment length=3pt}] (-1.5,0.75) -- (-0.3,0.75);
        \draw[fill=popblue!20, thick] (0.9,-0.6) circle (0.25);
        \draw[thick] (0.9,-0.35) -- (0.9,-1);
        \draw[thick] (0.55,-0.7) -- (1.25,-0.7);
        \draw[->, thick, popgreen] (0.2,-0.6) -- (0.55,-0.6);
    \end{scope}
    \node[icdesc] at (b2.south) {S'éloigner des fenêtres, miroirs et meubles hauts.};
    % 3. Pas d'ascenseur
    \node[picbox] (b3) at ({2*(\boxw+\gap)}, 0) {};
    \node[ictitle] at (b3.north) {PAS D'ASCENSEUR};
    \begin{scope}[shift={(b3.center)}]
        \draw[thick] (-0.8,1.3) rectangle (0.8,-1.3);
        \draw[fill=popmuted!40, thick] (-0.6,0) rectangle (0.6,0.7);
        \node[poppink, font=\bfseries\Large] at (0,1.05) {$\times$};
        \draw[->, thick, popgreen] (1.6,-1) -- (1.6,1) node[above, font=\scriptsize] {Escaliers};
    \end{scope}
    \node[icdesc] at (b3.south) {Ne jamais prendre l'ascenseur ; utiliser les escaliers après les secousses.};
    % 4. Dehors
    \node[picbox] (b4) at (0, -\boxh-\gap) {};
    \node[ictitle] at (b4.north) {À L'EXTÉRIEUR};
    \begin{scope}[shift={(b4.center)}]
        \draw[fill=poporange!20, thick] (-1.8,-1) rectangle (-0.7,0.8);
        \draw[fill=poporange!20, thick] (0.7,-1) rectangle (1.8,0.8);
        \draw[poppink, thick] (-1.2,0.8) -- (-1.0,-0.2);
        \draw[poppink, thick] (1.2,0.8) -- (1.4,-0.2);
        \draw[fill=popgreen!20, thick] (0,-0.5) circle (0.25);
        \draw[thick] (0,-0.25) -- (0,-0.9);
        \draw[thick] (-0.3,-0.6) -- (0.3,-0.6);
    \end{scope}
    \node[icdesc] at (b4.south) {S'écarter des bâtiments et poteaux ; rejoindre un espace dégagé.};
    % 5. Couper gaz/élec
    \node[picbox] (b5) at (\boxw+\gap, -\boxh-\gap) {};
    \node[ictitle] at (b5.north) {APRÈS : COUPER};
    \begin{scope}[shift={(b5.center)}]
        \draw[thick] (-1.2,0.3) -- (0.2,0.3);
        \draw[thick, fill=poppink!30] (-0.6,0.05) rectangle (-0.2,0.55);
        \draw[thick] (-0.4,0.05) -- (-0.4,-0.35);
        \node[font=\scriptsize] at (-0.5,-0.7) {Gaz};
        \draw[thick, fill=popmuted!40] (0.7,-0.4) rectangle (1.4,0.6);
        \draw[thick] (1.05,0.1) -- (1.05,-0.9);
        \node[font=\scriptsize] at (1.05,-1.15) {Électricité};
    \end{scope}
    \node[icdesc] at (b5.south) {Couper le gaz, l'électricité et l'eau pour éviter incendies et inondations.};
    % 6. Kit
    \node[picbox] (b6) at ({2*(\boxw+\gap)}, -\boxh-\gap) {};
    \node[ictitle] at (b6.north) {TROUSSE DE SECOURS};
    \begin{scope}[shift={(b6.center)}]
        \draw[thick, fill=popblue!10, rounded corners] (-1.1,-0.9) rectangle (1.1,0.4);
        \draw[thick] (-0.4,0.4) -- (-0.4,0.7) -- (0.4,0.7) -- (0.4,0.4);
        \node[font=\scriptsize, align=center] at (0,-0.25) {Eau, radio, lampe,\\médicaments, documents};
    \end{scope}
    \node[icdesc] at (b6.south) {Préparer une trousse de secours et écouter la radio.};
\end{tikzpicture}
\legende{Les comportements vitaux en cas de séisme : (1) se protéger sous une table, (2) s'éloigner des vitres, (3) ne pas prendre l'ascenseur, (4) à l'extérieur, rejoindre un espace dégagé, (5) après les secousses, couper le gaz et l'électricité, (6) disposer d'une trousse de secours.}
\end{popfigure}

\courssub{Éducation et sensibilisation : le pilier le plus efficace}

La meilleure construction parasismique ne suffit pas si la population ne sait pas réagir. L'\textbf{éducation au risque sismique} (exercices d'évacuation, affichage des consignes, plans d'urgence) est le moyen de prévention le moins coûteux et le plus efficace.

\begin{experience}[Démarche d'investigation : évaluer la préparation de son établissement]
\textbf{Question :} notre collège est-il prêt face à un séisme ?
\textbf{Protocole (par groupes) :}
\begin{enumerate}
    \item \textbf{Observer les bâtiments :} type de construction (béton armé, parpaings ?), nombre d'étages, fissures éventuelles.
    \item \textbf{Observer les salles :} les armoires et tableaux sont-ils fixés ? Les issues de secours sont-elles dégagées et signalées ?
    \item \textbf{Vérifier l'organisation :} existe-t-il un plan d'évacuation affiché ? Des exercices d'évacuation ont-ils lieu ? Un point de rassemblement est-il prévu, loin des bâtiments ?
\end{enumerate}
\textbf{Analyse et conclusion :} rédiger un court rapport avec des propositions concrètes (fixer les meubles hauts, dégager les issues, organiser un exercice d'évacuation par trimestre) et le présenter au chef d'établissement.
\end{experience}

\begin{exoresolu}[Exercice résolu : Haïti 2010 contre Chili 2010]
\textbf{Énoncé :} En 2010, un séisme de magnitude 7,0 frappe Haïti (environ 250 000 morts) et un séisme de magnitude 8,8 frappe le Chili (environ 500 morts). Le séisme chilien était pourtant beaucoup plus puissant. Expliquer cette différence de bilan humain.
\tcblower
\textbf{Solution détaillée :}
\begin{itemize}
    \item \textbf{Le danger (aléa) :} il était plus fort au Chili (magnitude 8,8 contre 7,0).
    \item \textbf{La fragilité des constructions (vulnérabilité) :} en Haïti, les bâtiments étaient en maçonnerie fragile, sans normes parasismiques, souvent sur des pentes instables. Au Chili, les normes de construction parasismique sont strictes et respectées (béton armé, ossatures souples).
    \item \textbf{La préparation de la population (prévention) :} au Chili, la population est entraînée (exercices, culture du risque) et les secours sont organisés ; en Haïti, la population n'était pas préparée.
\end{itemize}
\textbf{Conclusion :} le bilan d'un séisme ne dépend pas seulement de sa magnitude, mais surtout de la \textbf{qualité des constructions} et de la \textbf{préparation de la population}. C'est pourquoi on dit : « Ce ne sont pas les séismes qui tuent, ce sont les bâtiments. »
\end{exoresolu}

\begin{aretenir}[Bilan : la prévention prime sur la prévision]
\begin{itemize}
    \item \textbf{Prévision :} impossible à ce jour d'annoncer la date, le lieu et la magnitude d'un séisme. Les signes précurseurs (microséismes, niveau des puits, comportement des animaux) ne sont \textbf{pas fiables}. On sait seulement localiser les \textbf{zones à risque} (failles actives, limites de plaques).
    \item \textbf{Prévention (3 piliers) :} 1) \textbf{construire parasismique} (fondations profondes, ossature souple, contreventements) ; 2) \textbf{aménager le territoire} (éviter failles et sols instables) ; 3) \textbf{éduquer la population} (gestes qui sauvent, exercices d'évacuation, trousse de secours).
    \item \textbf{Au Cameroun :} le risque est modéré mais réel (faille de la Sanaga, Ligne du Cameroun, Mont Cameroun). Priorités : respecter les règles de construction et former les élèves aux bons réflexes.
\end{itemize}
\end{aretenir}

\begin{exercice}[Entraînement BEPC : prévention sismique]
\begin{enumerate}
    \item Un journaliste affirme : « Grâce à l'observation des animaux, nous pourrons bientôt prévoir les séismes une semaine à l'avance. » En t'appuyant sur le cours, explique pourquoi cette affirmation n'est pas scientifiquement fondée.
    \item Cite \textbf{trois} caractéristiques d'un bâtiment parasismique et explique brièvement le rôle de chacune.
    \item Un séisme survient alors que tu es en salle de classe au deuxième étage. Décris, dans l'ordre, les actions à effectuer pendant et après les secousses, en justifiant chaque geste.
    \item Deux villes A et B subissent un séisme de même magnitude. La ville A compte 1 000 morts, la ville B seulement 10. Propose deux explications possibles à cette différence.
\end{enumerate}
\end{exercice}

\begin{corrige}[Exercice : prévention sismique]
\begin{enumerate}
    \item Les comportements anormaux des animaux sont des observations \emph{après coup}, non mesurables et non reproductibles ; ils peuvent avoir d'autres causes (météo, maladie, prédateurs). Aucune étude rigoureuse n'a montré qu'ils permettent de prévoir un séisme : cette piste ne peut pas servir à déclencher une alerte fiable.
    \item \textbf{1. Fondations profondes} : ancrent le bâtiment dans le sol dur et limitent les déformations. \textbf{2. Ossature en béton armé ou en acier, souple} : le bâtiment se déforme sans rompre brutalement, ce qui évite l'effondrement. \textbf{3. Contreventements (renforts en croix)} : rigidifient l'ensemble et absorbent une partie de l'énergie des secousses.
    \item \textbf{Pendant les secousses :} se mettre sous une table solide et s'y accrocher (protection contre la chute d'objets et de plâtre), loin des fenêtres (bris de verre) ; ne pas courir vers la sortie (bousculade, chutes d'objets). \textbf{Après les secousses :} évacuer calmement par les escaliers (jamais l'ascenseur : risque de blocage), rejoindre le point de rassemblement à l'écart des bâtiments et des poteaux (chutes de façades, de câbles), se méfier des répliques et ne pas réintégrer le bâtiment.
    \item \textbf{Explication 1 (constructions) :} la ville A possède des bâtiments fragiles (maçonnerie non renforcée, sans normes), la ville B des bâtiments parasismiques. \textbf{Explication 2 (préparation) :} la population de la ville A n'était ni informée ni entraînée (pas d'exercices, pas de plans d'urgence), alors que celle de la ville B connaissait les bons réflexes et disposait de secours organisés.
\end{enumerate}
\end{corrige}

\newpage

\coursec{Activité d'intégration}

\coursec{Leçon 12 : Les séismes — Activité d'investigation : Localisation et caractérisation d'un séisme dans l'Ouest camerounais}

\courssub{Objectifs de l'activité}
\begin{itemize}
    \item Mobiliser les connaissances sur les ondes sismiques (ondes P et S), l'épicentre, le foyer, la magnitude et l'intensité pour résoudre un problème concret.
    \item Exploiter des sismogrammes pour déterminer l'intervalle de temps S-P et en déduire la distance épicentrale à l'aide d'un abaque.
    \item Localiser un épicentre par la méthode des cercles (trilatération) sur une carte du Cameroun.
    \item Estimer la magnitude d'un séisme à partir de l'amplitude maximale enregistrée et de la distance épicentrale (nomogramme de Richter).
    \item Évaluer l'intensité macrosismique attendue dans les localités proches et proposer des mesures de prévention adaptées au contexte camerounais.
    \item Rédiger un rapport scientifique structuré, argumenté et adressé aux autorités locales.
\end{itemize}

\courssub{Situation : Enquête sismologique dans la région de l'Ouest}
\label{sit:seisme}
Le \textbf{14 mars 2025 à 08 h 12 min 04 s TU} (09 h 12 min 04 s heure locale), une secousse tellurique est fortement ressentie dans la ville de Bafoussam et ses environs (région de l'Ouest, Cameroun). Des dégâts légers sont signalés : chutes d'objets, fissures sur des murs en banco, affolement des populations. Aucune victime n'est à déplorer pour l'instant.

Le Réseau Sismologique National (RSNC), dont les stations de \textbf{Bafoussam (BFS)}, \textbf{Douala (DLA)} et \textbf{Yaoundé (YAO)} font partie, a enregistré l'événement. Les trois sismogrammes verticaux (composante Z) sont mis à votre disposition (documents 1, 2, 3). L'épicentre est supposé se situer dans la région de l'Ouest, à faible profondeur (foyer superficiel, $h < 20$ km).

\vspace{0.3cm}
\noindent\textbf{Données géographiques des stations (approximatives) :}
\begin{center}
\begin{tabularx}{\linewidth}{|c|X|X|X|}\hline
\textbf{Station} & \textbf{Code} & \textbf{Latitude (°N)} & \textbf{Longitude (°E)} \\ \hline
Bafoussam & BFS & 5,47 & 10,42 \\ \hline
Douala & DLA & 4,05 & 9,70 \\ \hline
Yaoundé & YAO & 3,87 & 11,52 \\ \hline
\end{tabularx}
\end{center}

\vspace{0.3cm}
\noindent\textbf{Document 1 : Sismogramme de la station Bafoussam (BFS) — Proche de l'épicentre}
\begin{popfigure}
\begin{tikzpicture}[scale=0.8]
% Axes
\draw[->] (0,0) -- (14,0) node[right] {Temps (s) après 08:12:04 TU};
\draw[->] (0,-2.5) -- (0,2.5) node[above] {Amplitude (mm)};
% Grid
\foreach \x in {0,2,...,14} \draw[gray!30] (\x,-2.5) -- (\x,2.5);
\foreach \y in {-2,-1,0,1,2} \draw[gray!30] (0,\y) -- (14,\y);
% Time labels
\foreach \x/\t in {0/0, 2/2, 4/4, 6/6, 8/8, 10/10, 12/12, 14/14} \node[below] at (\x, -2.7) {\small \t};
% Signal P wave (arrival ~1.5s)
\draw[popblue, thick] (0,0) -- (1.5,0);
\draw[popblue, thick, decorate, decoration={snake, amplitude=0.5mm, segment length=3mm}] (1.5,0) -- (3.5,0) node[midway, above=2mm] {\small \textbf{Onde P}};
% Signal S wave (arrival ~4.5s)
\draw[popblue, thick] (3.5,0) -- (4.5,0);
\draw[poporange, thick, decorate, decoration={snake, amplitude=1.5mm, segment length=5mm}] (4.5,0) -- (8,0) node[midway, above=4mm] {\small \textbf{Onde S}};
% Coda
\draw[popmuted, thick, decorate, decoration={snake, amplitude=0.8mm, segment length=4mm}] (8,0) -- (14,0);
% Max Amplitude marker on S wave
\draw[dashed, popdark] (6, -1.8) -- (6, 1.8);
\node[popdark, right] at (6, 1.8) {$A_{\text{max}} = 15,0$ mm};
% Arrow for P arrival
\draw[<-, popblue] (1.5, 2.2) -- (1.5, 1.5) node[above] {$t_P = 1,5$ s};
% Arrow for S arrival
\draw[<-, poporange] (4.5, -2.2) -- (4.5, -1.5) node[below] {$t_S = 4,5$ s};
\end{tikzpicture}
\legende{Sismogramme simulé station BFS (distance $\approx$ 25 km). L'arrivée P est nette à 1,5 s ; l'arrivée S à 4,5 s. Amplitude max sur l'onde S : 15,0 mm.}
\end{popfigure}

\vspace{0.3cm}
\noindent\textbf{Document 2 : Sismogramme de la station Douala (DLA) — Distance intermédiaire}
\begin{popfigure}
\begin{tikzpicture}[scale=0.8]
\draw[->] (0,0) -- (14,0) node[right] {Temps (s) après 08:12:04 TU};
\draw[->] (0,-1.5) -- (0,1.5) node[above] {Amplitude (mm)};
\foreach \x in {0,2,...,14} \draw[gray!30] (\x,-1.5) -- (\x,1.5);
\foreach \y in {-1,0,1} \draw[gray!30] (0,\y) -- (14,\y);
% Time labels: window 25s to 60s
\foreach \x/\t in {0/25, 2/27, 4/29, 6/31, 8/33, 10/35, 12/37, 14/39} \node[below] at (\x, -1.7) {\small \t};
\foreach \x/\t in {0/41, 2/43, 4/45, 6/47, 8/49, 10/51, 12/53, 14/55} \node[below, gray] at (\x, -2.0) {\small \t};
% P wave at ~30s (x=2.5)
\draw[popblue, thick] (0,0) -- (2.5,0);
\draw[popblue, thick, decorate, decoration={snake, amplitude=0.3mm, segment length=2mm}] (2.5,0) -- (4.5,0) node[midway, above=2mm] {\small \textbf{Onde P}};
% S wave at ~51s (x=13)
\draw[popblue, thick] (4.5,0) -- (13,0);
\draw[poporange, thick, decorate, decoration={snake, amplitude=0.8mm, segment length=4mm}] (13,0) -- (14,0) node[midway, above=3mm] {\small \textbf{Onde S}};
% Markers
\draw[<-, popblue] (2.5, 1.7) -- (2.5, 1.0) node[above] {$t_P \approx 30$ s};
\draw[<-, poporange] (13, -1.7) -- (13, -1.0) node[below] {$t_S \approx 51$ s};
\draw[dashed, popdark] (13.5, -1.2) -- (13.5, 1.2);
\node[popdark, right] at (13.5, 1.2) {$A_{\text{max}} = 2,5$ mm};
\end{tikzpicture}
\legende{Sismogramme simulé station DLA (distance $\approx$ 176 km). Fenêtre temporelle 25--55 s. Arrivée P vers 30 s, arrivée S vers 51 s. Amplitude plus faible qu'à BFS.}
\end{popfigure}

\vspace{0.3cm}
\noindent\textbf{Document 3 : Sismogramme de la station Yaoundé (YAO) — Distance intermédiaire}
\begin{popfigure}
\begin{tikzpicture}[scale=0.8]
\draw[->] (0,0) -- (14,0) node[right] {Temps (s) après 08:12:04 TU};
\draw[->] (0,-1.5) -- (0,1.5) node[above] {Amplitude (mm)};
\foreach \x in {0,2,...,14} \draw[gray!30] (\x,-1.5) -- (\x,1.5);
\foreach \y in {-1,0,1} \draw[gray!30] (0,\y) -- (14,\y);
% Time labels: window 25s to 55s
\foreach \x/\t in {0/25, 2/27, 4/29, 6/31, 8/33, 10/35, 12/37, 14/39} \node[below] at (\x, -1.7) {\small \t};
\foreach \x/\t in {0/41, 2/43, 4/45, 6/47, 8/49, 10/51, 12/53, 14/55} \node[below, gray] at (\x, -2.0) {\small \t};
% P wave at ~29s (x=2.0)
\draw[popblue, thick] (0,0) -- (2.0,0);
\draw[popblue, thick, decorate, decoration={snake, amplitude=0.3mm, segment length=2mm}] (2.0,0) -- (4.0,0) node[midway, above=2mm] {\small \textbf{Onde P}};
% S wave at ~47s (x=11)
\draw[popblue, thick] (4.0,0) -- (11,0);
\draw[poporange, thick, decorate, decoration={snake, amplitude=0.9mm, segment length=4mm}] (11,0) -- (14,0) node[midway, above=3mm] {\small \textbf{Onde S}};
% Markers
\draw[<-, popblue] (2.0, 1.7) -- (2.0, 1.0) node[above] {$t_P \approx 29$ s};
\draw[<-, poporange] (11, -1.7) -- (11, -1.0) node[below] {$t_S \approx 47$ s};
\draw[dashed, popdark] (11.5, -1.2) -- (11.5, 1.2);
\node[popdark, right] at (11.5, 1.2) {$A_{\text{max}} = 3,0$ mm};
\end{tikzpicture}
\legende{Sismogramme simulé station YAO (distance $\approx$ 151 km). Fenêtre temporelle 25--55 s. Arrivée P vers 29 s, arrivée S vers 47 s.}
\end{popfigure}

\vspace{0.3cm}
\noindent\textbf{Données numériques extraites des sismogrammes (Documents 1 à 3) :}
\begin{center}
\begin{tabularx}{\linewidth}{|c|X|X|X|X|}\hline
\textbf{Station} & \textbf{Heure arrivée onde P ($t_P$)} & \textbf{Heure arrivée onde S ($t_S$)} & \textbf{Intervalle S-P ($\Delta t$)} & \textbf{Amplitude max $A_{\text{max}}$ (mm)} \\ \hline
Bafoussam (BFS) & 08 h 12 min 05,5 s & 08 h 12 min 08,5 s & \textbf{3,0 s} & \textbf{15,0} \\ \hline
Douala (DLA) & 08 h 12 min 34,0 s & 08 h 12 min 55,0 s & \textbf{21,0 s} & \textbf{2,5} \\ \hline
Yaoundé (YAO) & 08 h 12 min 29,0 s & 08 h 12 min 47,0 s & \textbf{18,0 s} & \textbf{3,0} \\ \hline
\end{tabularx}
\end{center}

\vspace{0.3cm}
\noindent\textbf{Document 4 : Abaque distance épicentrale en fonction de l'intervalle S-P (modèle crustal moyen Cameroun)}
\begin{popfigure}
\begin{tikzpicture}
\begin{axis}[
    width=\linewidth, height=6cm,
    xlabel={Intervalle S-P $\Delta t$ (s)}, ylabel={Distance épicentrale $D$ (km)},
    xmin=0, xmax=30, ymin=0, ymax=260,
    grid=major, grid style={gray!30},
    tick label style={font=\small},
    label style={font=\small},
    legend style={at={(0.95,0.05)}, anchor=south east, font=\small, fill=white, fill opacity=0.8}
]
\addplot[popcurve, thick, domain=0:30, samples=50] {8.4*x};
\addlegendentry{$D \approx 8,4 \times \Delta t$ (modèle théorique)}
\addplot[only marks, mark=*, mark size=3, popblue] coordinates {(3, 25.2)}; \addlegendentry{BFS (3,0 s)};
\addplot[only marks, mark=square*, mark size=3, poporange] coordinates {(21, 176.4)}; \addlegendentry{DLA (21,0 s)};
\addplot[only marks, mark=triangle*, mark size=3, popgreen] coordinates {(18, 151.2)}; \addlegendentry{YAO (18,0 s)};
\draw[dashed, popmuted] (3,0) -- (3,25.2) -- (0,25.2);
\draw[dashed, popmuted] (21,0) -- (21,176.4) -- (0,176.4);
\draw[dashed, popmuted] (18,0) -- (18,151.2) -- (0,151.2);
\end{axis}
\end{tikzpicture}
\legende{Abaque $\Delta t \rightarrow D$ utilisé pour l'activité. La relation affine $D = k \cdot \Delta t$ avec $k \approx 8,4$ km/s est adaptée au contexte crustal camerounais (vitesse P $\approx 6,0$ km/s, vitesse S $\approx 3,5$ km/s).}
\end{popfigure}

\vspace{0.3cm}
\noindent\textbf{Document 5 : Nomogramme simplifié de magnitude locale $M_L$ (type Richter) pour le Cameroun}
\begin{popfigure}
\begin{tikzpicture}
\begin{axis}[
    width=\linewidth, height=6cm,
    xlabel={Distance épicentrale $D$ (km)}, ylabel={Amplitude maximale $A$ (mm)},
    xmode=log, ymode=log,
    xmin=10, xmax=300, ymin=0.1, ymax=100,
    grid=both, grid style={gray!30},
    tick label style={font=\small},
    label style={font=\small},
    legend style={at={(0.05,0.95)}, anchor=north west, font=\small, fill=white, fill opacity=0.8}
]
\pgfplotsforeachungrouped \M/\Y in {3.0/1, 3.5/3.1623, 4.0/10, 4.5/31.623, 5.0/100} {
    \addplot[popmuted, thin, dashed, domain=10:300, samples=50] ({x}, {pow(10, \M - log10(x) - 1.0)});
    \node[popmuted, font=\tiny, rotate=-30] at (axis cs: 100, \Y) {$M_L = \M$};
}
\addplot[only marks, mark=*, mark size=3, popblue] coordinates {(25.2, 15.0)}; \addlegendentry{BFS};
\addplot[only marks, mark=square*, mark size=3, poporange] coordinates {(176.4, 2.5)}; \addlegendentry{DLA};
\addplot[only marks, mark=triangle*, mark size=3, popgreen] coordinates {(151.2, 3.0)}; \addlegendentry{YAO};
\end{axis}
\end{tikzpicture}
\legende{Nomogramme magnitude-distance-amplitude (échelles logarithmiques). Les droites pointillées correspondent aux magnitudes $M_L$ constantes. La magnitude se lit à l'intersection du point (Distance, Amplitude) avec les isolignes de magnitude.}
\end{popfigure}

\vspace{0.3cm}
\noindent\textbf{Document 6 : Carte de travail pour la trilatération (Repère orthonormé, 1 unité = 50 km)}
Les coordonnées des stations dans ce repère local (origine arbitraire, axes Est/Nord) sont choisies pour respecter les distances inter-stations réelles (BFS-DLA $\approx$ 160 km, BFS-YAO $\approx$ 150 km, DLA-YAO $\approx$ 200 km) et permettre un tracé précis au compas.
\begin{center}
\begin{tabularx}{\linewidth}{|c|X|X|}\hline
\textbf{Station} & \textbf{Abscisse $x$ (unités)} & \textbf{Ordonnée $y$ (unités)} \\ \hline
Bafoussam (BFS) & 20,0 & 20,0 \\ \hline
Douala (DLA) & 16,8 & 20,0 \\ \hline
Yaoundé (YAO) & 19,5 & 23,0 \\ \hline
\end{tabularx}
\end{center}

\begin{popfigure}
\begin{tikzpicture}[scale=0.6]
% Grid
\draw[gray!20, step=1] (14,14) grid (27,27);
\draw[->, thin, gray] (14,14) -- (27,14) node[right] {Est (x 50 km)};
\draw[->, thin, gray] (14,14) -- (14,27) node[above] {Nord (x 50 km)};
% Stations
\node[popblue, circle, fill, inner sep=2.5pt, label=above right:\textbf{BFS (20;20)}] (BFS) at (20,20) {};
\node[poporange, circle, fill, inner sep=2.5pt, label=below left:\textbf{DLA (16,8;20)}] (DLA) at (16.8,20) {};
\node[popgreen, circle, fill, inner sep=2.5pt, label=above left:\textbf{YAO (19,5;23)}] (YAO) at (19.5,23) {};
% Epicenter (Solution - faint)
\node[popred, star, fill, inner sep=3pt, label=right:\textbf{E}] (EPI) at (20.5,20) {};
% Circles (Solution - faint)
\draw[popblue, dashed, thick] (BFS) circle (0.5); % 25 km = 0.5 units
\draw[poporange, dashed, thick] (DLA) circle (3.53); % 176 km = 3.53 units
\draw[popgreen, dashed, thick] (YAO) circle (3.02); % 151 km = 3.02 units
% Scale
\draw[popdark, thick] (24, 14.5) -- (26, 14.5) node[midway, below] {100 km};
\end{tikzpicture}
\legende{Carte de travail (repère 1 unité = 50 km). Les cercles solution sont tracés en pointillés pour le corrigé ; l'élève doit les tracer au compas. L'épicentre E est au Nord-Est immédiat de BFS.}
\end{popfigure}

\vspace{0.3cm}
\noindent\textbf{Document 7 : Échelle d'intensité macrosismique simplifiée (EMS-98 adaptée)}
\begin{center}
\begin{tabularx}{\linewidth}{|c|X|}\hline
\textbf{Degré} & \textbf{Description des effets observés} \\ \hline
I -- II & Non ressenti ou ressenti seulement par quelques personnes au repos. \\ \hline
III -- IV & Ressenti à l'intérieur par plusieurs personnes ; objets qui oscillent ; bruits. \\ \hline
V -- VI & Ressenti par tous ; peur ; chutes d'objets ; fissures légères (degré VI). \\ \hline
VII -- VIII & Dégâts légers à modérés sur bâtiments vulnérables (banco, parpaings mal ferraillés) ; panique. \\ \hline
IX + & Dégâts importants, effondrements. \\ \hline
\end{tabularx}
\end{center}
\noindent\textit{Règle empirique (approximation pour foyer superficiel $h \approx 10$ km) :}
\[ I_{\text{épicentre}} \approx M_L + 1 \quad \text{et} \quad I(D) \approx I_{\text{épicentre}} - 3 \log_{10}(D/10) \quad (D \text{ en km}) \]

\courssub{Consignes}
Travail en groupes de 3 ou 4 élèves. Chaque groupe rend un rapport écrit unique.
\begin{enumerate}
    \item \textbf{Analyse des sismogrammes :} Pour chaque station (BFS, DLA, YAO), relever l'intervalle de temps $\Delta t = t_S - t_P$ (donné dans le tableau) et, à l'aide de l'abaque (Document 4), déterminer la distance épicentrale $D$ (en km). Compléter le tableau de résultats.
    \item \textbf{Localisation de l'épicentre :} Sur la carte fournie (Document 6), tracer au compas trois cercles centrés sur chaque station et de rayon égal à la distance épicentrale calculée (échelle : 1 unité = 50 km). Déterminer les coordonnées cartographiques approximatives $(x_E ; y_E)$ de l'épicentre (intersection des cercles). Reporter l'épicentre sur la carte.
    \item \textbf{Estimation de la magnitude :} Pour chaque station, reporter le point de coordonnées $(D ; A_{\text{max}})$ sur le nomogramme (Document 5). Lire la magnitude locale $M_L$ correspondante. Donner la valeur moyenne de $M_L$ retenue pour ce séisme.
    \item \textbf{Évaluation de l'intensité :} En utilisant la règle empirique du Document 7, calculer l'intensité macrosismique attendue à l'épicentre ($I_0$) et dans les villes de Bafoussam, Douala et Yaoundé. Classer les effets attendus selon l'échelle EMS-98.
    \item \textbf{Rédaction du rapport d'alerte :} Rédiger un rapport adressé au \textbf{Préfet du Département de la Mifi} (chef-lieu Bafoussam), structuré ainsi :
    \begin{itemize}
        \item \textit{Objet :} Localisation et caractérisation du séisme du 14 mars 2025.
        \item \textit{Localisation :} Coordonnées géographiques approximatives (Lat, Long) de l'épicentre et distance aux villes principales.
        \item \textit{Caractéristiques :} Magnitude $M_L$, Intensité épicentrale $I_0$, Intensités attendues à Bafoussam, Douala, Yaoundé.
        \item \textit{Recommandations (3 mesures) :} Proposer trois mesures concrètes de prévention/protection immédiate ou à moyen terme, justifiées par la vulnérabilité locale (habitat en banco, densité de population, infrastructures).
    \end{itemize}
    \item \textbf{Restitution orale :} Un porte-parole par groupe présente les résultats en 3 minutes. Confrontation des localisations et magnitudes trouvées. Discussion sur les incertitudes.
\end{enumerate}

\courssub{Ressources à mobiliser}
\begin{definition}[Ondes sismiques et distances]
Dans un milieu homogène, les ondes P (primaires, longitudinales) propagent plus vite ($v_P \approx \SI{6,0}{km/s}$ dans la croûte supérieure) que les ondes S (secondaires, transversales, $v_S \approx \SI{3,5}{km/s}$). L'intervalle S-P $\Delta t$ croît avec la distance épicentrale $D$ :
\[ D \approx \frac{v_P v_S}{v_P - v_S} \cdot \Delta t = k \cdot \Delta t \]
Au Cameroun, le coefficient $k$ vaut environ \textbf{8,4 km/s} (abaque Document 4).
\end{definition}

\begin{definition}[Magnitude locale $M_L$ (Richter)]
La magnitude quantifie l'énergie libérée à la source. Pour un sismographe de type Wood-Anderson (simulé ici), la formule pratique est :
\[ M_L = \log_{10}(A_{\text{max}}) + \alpha \log_{10}(D) + \beta \]
où $A_{\text{max}}$ est l'amplitude maximale en mm (crête-crête/2), $D$ la distance en km. Le nomogramme (Doc 5) permet une lecture graphique directe sans calcul logarithmique complexe.
\end{definition}

\begin{definition}[Intensité macrosismique (EMS-98)]
L'intensité qualifie les effets en surface. Elle décroît avec la distance $D$ (en km) par rapport à l'intensité épicentrale $I_0$ :
\[ I(D) \approx I_0 - 3 \log_{10}\left(\frac{D}{10}\right) \]
Pour un foyer superficiel, on admet souvent $I_0 \approx M_L + 1$ (valeur arrondie à l'entier supérieur).
\end{definition}

\begin{methode}[Trilatération (Localisation par intersection de cercles)]
\begin{enumerate}
    \item Convertir les distances $D$ (km) en rayon $R$ sur la carte : $R = D / 50$ (échelle 1 unité = 50 km).
    \item Tracer les trois cercles au compas (pointe sur la station, ouverture $R$).
    \item L'épicentre est la zone de convergence des trois cercles (idéalement un point, pratiquement un petit triangle d'erreur).
    \item Lire les coordonnées $(x_E ; y_E)$ du centroïde de cette zone.
    \item Convertir en coordonnées géographiques si nécessaire (en se référant aux coordonnées réelles des stations du tableau initial).
\end{enumerate}
\end{methode}

\courssub{Démarche et corrigé commenté}
\begin{exoresolu}{Corrigé détaillé de l'enquête sismologique}
\tcblower

\textbf{1. Détermination des distances épicentrales (Consigne 1)}
\medskip
On utilise l'abaque (Doc 4) ou la formule $D = 8,4 \times \Delta t$.
\begin{center}
\begin{tabular}{|c|c|c|c|}\hline
\textbf{Station} & $\Delta t$ (s) & \textbf{Calcul} $D = 8,4 \times \Delta t$ & \textbf{Distance $D$ (km)} \\ \hline
Bafoussam (BFS) & 3,0 & $8,4 \times 3,0$ & \textbf{25,2 km} \\ \hline
Douala (DLA) & 21,0 & $8,4 \times 21,0$ & \textbf{176,4 km} \\ \hline
Yaoundé (YAO) & 18,0 & $8,4 \times 18,0$ & \textbf{151,2 km} \\ \hline
\end{tabular}
\end{center}
\noindent\textit{Justification :} L'intervalle S-P est proportionnel à la distance pour un modèle de vitesse constant. La station la plus proche (BFS) a le plus petit $\Delta t$, ce qui est cohérent avec le ressenti fort sur place.

\medskip
\textbf{2. Localisation de l'épicentre par trilatération (Consigne 2)}
\medskip
Échelle carte : 1 unité = 50 km. Rayons des cercles à tracer :
\begin{itemize}
    \item $R_{\text{BFS}} = 25,2 / 50 = \mathbf{0,50}$ unité.
    \item $R_{\text{DLA}} = 176,4 / 50 = \mathbf{3,53}$ unités.
    \item $R_{\text{YAO}} = 151,2 / 50 = \mathbf{3,02}$ unités.
\end{itemize}
\noindent\textit{Tracé sur la carte (Doc 6) :}
\begin{itemize}
    \item Cercle BFS : Centre $(20,0; 20,0)$, rayon $0,50$ (très petit cercle autour de Bafoussam).
    \item Cercle DLA : Centre $(16,8; 20,0)$, rayon $3,53$.
    \item Cercle YAO : Centre $(19,5; 23,0)$, rayon $3,02$.
\end{itemize}
Les trois cercles se coupent dans une zone réduite autour du point de coordonnées \textbf{$(20,5 ; 20,0)$}.
\noindent\textbf{Coordonnées de l'épicentre E (carte) :} $x_E \approx 20,5$ ; $y_E \approx 20,0$.
\noindent\textit{Conversion en coordonnées géographiques (approximatif, à partir des coordonnées réelles des stations) :}
L'épicentre se situe à \textbf{25 km au Nord-Est de Bafoussam}.
$\text{Lat}_E \approx 5,47^\circ + (25/111) \approx \mathbf{5,7^\circ \text{ N}}$
$\text{Long}_E \approx 10,42^\circ + (25/110) \approx \mathbf{10,6^\circ \text{ E}}$
L'épicentre se situe dans la zone des hauts plateaux (proche du Mont Bamboutos), zone connue pour sa sismicité modérée.

\medskip
\textbf{3. Estimation de la magnitude $M_L$ (Consigne 3)}
\medskip
Lecture graphique sur le nomogramme (Doc 5) pour chaque station $(D ; A_{\text{max}})$ :
\begin{itemize}
    \item \textbf{BFS} $(25,2 ; 15,0)$ : Le point se situe sur l'isoligne $M_L \approx \mathbf{4,3}$.
    \item \textbf{DLA} $(176,4 ; 2,5)$ : Le point se situe entre $M_L = 4,0$ et $4,5$, proche de \textbf{4,2}.
    \item \textbf{YAO} $(151,2 ; 3,0)$ : Le point se situe sur l'isoligne $M_L \approx \mathbf{4,3}$.
\end{itemize}
\noindent\textbf{Magnitude moyenne retenue :} $M_L = \dfrac{4,3 + 4,2 + 4,3}{3} \approx \mathbf{4,3}$.
\noindent\textit{Commentaire :} La cohérence des trois estimations (écart < 0,2) valide la qualité des mesures et l'hypothèse d'un foyer unique. C'est un séisme de magnitude \textbf{modérée} (ressenti fort localement, dégâts légers possibles).

\medskip
\textbf{4. Évaluation de l'intensité macrosismique (Consigne 4)}
\medskip
Intensité épicentrale estimée : $I_0 \approx M_L + 1 = 4,3 + 1 = 5,3 \rightarrow \mathbf{VI}$ (arrondi au degré supérieur pour la sécurité).
\medskip
Calcul de l'intensité à distance $D$ : $I(D) \approx 5,3 - 3 \log_{10}(D/10)$.
\begin{center}
\begin{tabularx}{\linewidth}{|c|c|c|c|X|}\hline
\textbf{Localité} & \textbf{Distance $D$ (km)} & \textbf{Calcul $I(D)$} & \textbf{Degré EMS} & \textbf{Effets attendus} \\ \hline
Épicentre (E) & 0 (foyer $\approx$ 10 km) & $I_0 \approx 5,3$ & \textbf{VI} & Fissures légères, chutes d'objets, peur. \\ \hline
Bafoussam & 25,2 & $5,3 - 3 \log_{10}(2,52) \approx 5,3 - 1,2 = 4,1$ & \textbf{IV-V} & Ressenti par tous, objets oscillants, bruit. \\ \hline
Douala & 176,4 & $5,3 - 3 \log_{10}(17,6) \approx 5,3 - 3,7 = 1,6$ & \textbf{II} & Ressenti par quelques personnes au repos. \\ \hline
Yaoundé & 151,2 & $5,3 - 3 \log_{10}(15,1) \approx 5,3 - 3,5 = 1,8$ & \textbf{II} & Ressenti par quelques personnes au repos. \\ \hline
\end{tabularx}
\end{center}
\noindent\textit{Analyse :} L'intensité forte (VI) est très localisée près de l'épicentre. Bafoussam (IV-V) subit un secouement modéré compatible avec les témoignages (objets chutent, fissures sur banco). Douala et Yaoundé ne ressentent qu'une vibration légère (II), souvent imperçue en activité.

\medskip
\textbf{5. Rapport d'alerte au Préfet de la Mifi (Consigne 5)}
\medskip
\noindent\textbf{RAPPORT D'INFORMATION SISMIQUE URGENT}\\
\textbf{Objet :} Séisme du 14 mars 2025 – Région de l'Ouest (Département de la Mifi).\\
\textbf{De :} Cellule de veille sismologique – RSNC / Lycée de Bafoussam (Exercice pédagogique).\\
\textbf{À :} Monsieur le Préfet du Département de la Mifi.\\
\textbf{Date :} 14 mars 2025 – 10 h 00 locales.

\medskip
\noindent\textbf{1. Localisation de l'épicentre}\\
L'épicentre est situé aux coordonnées géographiques approximatives \textbf{5,7° N ; 10,6° E}, soit à \textbf{25 km au Nord-Est de Bafoussam}, dans le secteur des Hauts-Plateaux (proximité Mont Bamboutos). Le foyer est superficiel (profondeur estimée $< 15$ km).

\medskip
\noindent\textbf{2. Caractéristiques du séisme}\\
\begin{itemize}
    \item \textbf{Magnitude locale $M_L$ :} \textbf{4,3} (échelle de Richter). Séisme modéré.
    \item \textbf{Intensité épicentrale $I_0$ :} \textbf{VI} (EMS-98) – Dégâts légers sur bâtiments vulnérables.
    \item \textbf{Intensités attendues :}
    \begin{itemize}
        \item Bafoussam (25 km) : \textbf{Degré IV-V} (Secousse modérée, chutes d'objets, fissures légères sur habitat précaire).
        \item Douala (176 km) : \textbf{Degré II} (Ressenti faiblement au repos).
        \item Yaoundé (151 km) : \textbf{Degré II} (Ressenti faiblement au repos).
    \end{itemize}
\end{itemize}

\medskip
\noindent\textbf{3. Recommandations immédiates (Prévention / Protection)}\\
\noindent\textbf{Mesure 1 – Inspection et évacuation préventive des bâtiments publics vulnérables :}\\
Ordonner l'inspection immédiate par le Génie Militaire ou la Protection Civile des écoles, hôpitaux (Hôpital Régional de Bafoussam) et bâtiments administratifs construits en banco ou parpaings non ferraillés (type "parpaing creux" sans chaînages). Évacuer préventivement les locaux présentant des fissures traversantes récentes jusqu'à expertise structurelle. \textit{Justification : L'intensité VI-V cause des dégâts significatifs sur l'habitat informel majoritaire en zone urbaine périphérique.}

\noindent\textbf{Mesure 2 – Information et consignes de sécurité à la population (Médias locaux + CRTV) :}\\
Diffuser en langues locales (Bamileke, Français) les consignes "Se protéger (table solide), S'éloigner des façades/vitres, Couper gaz/électricité si possible, Ne pas utiliser les ascenseurs". Interdire l'accès aux zones de falaise instables (carrières, bords de route RN6) suite aux risques d'éboulements déclenchés par la secousse. \textit{Justification : La panique cause plus de blessés que la secousse elle-même ; l'habitat en hauteur (étages) sans ascenseur est à risque.}

\noindent\textbf{Mesure 3 – Renforcement de la surveillance sismique et préparation post-crise :}\\
Demander au RSNC/IRGM le déploiement d'une station mobile temporaire (sismomètre large bande) sur l'épicentre pour enregistrer les répliques (probables dans les 48 h, magnitude $M_L - 1 \approx 3,3$). Pré-positionner des tentes d'urgence (stocks MINAS/Protection Civile) à Bafoussam et Batcham pour d'éventuels sinistrés si une réplique plus forte ($M_L > 4,5$) survient. \textit{Justification : La séquence sismique n'est pas terminée ; la loi d'Omori prévoit des répliques. La vulnérabilité sociale (logements précaires) impose une anticipation humanitaire.}

\medskip
\noindent\textbf{Signature :} Chef de la Cellule de Veille / Enseignant SVT.
\end{exoresolu}

\courssub{Bilan de l'activité}
\begin{aretenir}[Ce que cette activité a permis de mettre en évidence]
\begin{itemize}
    \item \textbf{La méthode des cercles (trilatération)} est la technique fondamentale de localisation d'un épicentre : elle nécessite au minimum 3 stations pour lever l'ambiguïté (intersection de 2 cercles = 2 points possibles).
    \item \textbf{L'intervalle S-P} est l'observable direct sur un sismogramme qui permet de remonter à la distance, grâce à la différence de vitesse entre ondes P et S ($v_P > v_S$). La relation est linéaire pour des distances locales ($D = k \Delta t$).
    \item \textbf{Magnitude vs Intensité :} La magnitude ($M_L = 4,3$) est une \textbf{unique valeur} caractéristique de la source (énergie). L'intensité (VI à l'épicentre, II à Douala) est une \textbf{valeur variable dans l'espace} qui décrit les effets sur l'environnement et les humains ; elle décroît avec la distance.
    \item \textbf{Le contexte camerounais} (foyers superficiels sur la chaîne volcanique du Cameroun / ligne du Cameroun, habitat en banco, forte densité urbaine) amplifie la vulnérabilité : un séisme modéré ($M_L \approx 4,5$) peut causer des dégâts locaux significatifs (degré VI-VII) sur les constructions non parasismiques.
    \item \textbf{La prévention} ne se limite pas à la prévision (impossible à court terme) mais repose sur : la \textbf{réglementation parasismique} (normes de construction Eurocode 8 adaptées), l'\textbf{éducation aux gestes qui sauvent} (exercices d'évacuation), et la \textbf{surveillance instrumentale} (réseau RSNC/IRGM) pour l'alerte rapide et le suivi des répliques.
\end{itemize}
\end{aretenir}

\begin{attention}[Pièges fréquents et erreurs à éviter]
\begin{itemize}
    \item \textbf{Confondre $t_P$ et $t_S$ :} L'onde P arrive \textbf{toujours} en premier. Sur le sismogramme, c'est la première rupture de pente nette (souvent faible amplitude). L'onde S arrive plus tard, avec une amplitude plus grande et un signal plus complexe.
    \item \textbf{Unités sur l'abaque/nomogramme :} Vérifier que la distance $D$ est en \textbf{km} et l'amplitude $A$ en \textbf{mm} (crête-crête/2). Ne pas confondre amplitude en mm et en µm (facteur 1000 sur la magnitude !).
    \item \textbf{Échelle de la carte :} Ne pas oublier de convertir la distance réelle (km) en rayon graphique (cm ou unités de repère) avant de tracer le compas. Erreur classique : tracer un cercle de rayon 176 cm au lieu de 3,5 cm !
    \item \textbf{Précision de l'intersection :} Les cercles ne se coupent jamais en un point mathématique parfait (incertitudes sur les vitesses, lecture $\Delta t$, profondeur du foyer). On localise l'épicentre au \textbf{centre du triangle d'erreur} formé par les trois arcs de cercle.
    \item \textbf{Magnitude moyenne :} Ne pas faire la moyenne des magnitudes si l'une est aberrante (ex: station saturée). Vérifier la cohérence (écart < 0,3-0,5). Ici, la cohérence est excellente.
\end{itemize}
\end{attention}

\begin{savaistu}[Le contexte sismique du Cameroun : La "Ligne du Cameroun"]
Le Cameroun n'est pas une zone de forte sismicité comme le Japon ou la Californie, mais il présente une \textbf{sismicité modérée intraplaque} concentrée le long de la \textbf{Ligne du Cameroun} (alignement volcanique Mont Cameroun – Manengouba – Bamboutos – Adamaoua – Lac Tchad).
\begin{itemize}
    \item \textbf{Séismes historiques :} 1986 (Lac Nyos, dégazage + secousse), 2001 (Mont Cameroun, éruption + séisme $M_L \approx 4,5$), 2018 (Bafia, $M_L \approx 4,0$), 2023 (Ngaoundéré, $M_L \approx 4,2$).
    \item \textbf{Mécanisme :} Failles normales et décrochantes liées à l'extension lithosphérique (rift) et au magmatisme. Foyers souvent superficiels ($< 15$ km), ce qui maximise l'intensité en surface pour une magnitude donnée.
    \item \textbf{Surveillance :} L'Institut de Recherche Géologique et Minière (IRGM) gère le Réseau Sismologique National (RSNC) avec une dizaine de stations (dont Bafoussam, Douala, Yaoundé, Ngaoundéré, Garoua, Buéa). Les données sont échangées en temps réel avec le réseau international (ISC, EMSC).
    \item \textbf{Risque majeur :} La croissance urbaine non planifiée (Yaoundé, Douala, Bafoussam, Bamenda) sur des sols meubles (alluvions, latérites) qui amplifient le secouement (effet de site), couplée à un parc immobilier largement non conforme aux normes parasismiques (Eurocode 8 zone 1-2).
\end{itemize}
\end{savaistu}

\newpage

\coursec{Méthodes et exercices résolus}

\courssub{Méthode 1 : Décrire et comparer les manifestations d'un séisme}

\begin{definition}[Séisme]
Un \cle{séisme} (ou tremblement de terre) est une série de vibrations brèves et brusques du sol, provoquées par la libération soudaine de l'énergie accumulée dans les roches de la lithosphère.
\end{definition}

\begin{methode}[Analyser les manifestations d'un séisme]
Pour décrire les manifestations d'un séisme à partir d'un document (témoignage, photographie, article de presse) :
\begin{enumerate}
\item \textbf{Repérer les effets sur les constructions} : fissures dans les murs, effondrement de maisons, destruction de ponts et de routes.
\item \textbf{Repérer les effets sur le sol} : fissures du sol, glissements de terrain, affaissements, jaillissement d'eau.
\item \textbf{Repérer les effets sur les personnes et les biens} : blessés, morts, dégâts matériels, incendies.
\item \textbf{Classer les manifestations} en deux catégories : effets \cle{directs} (vibrations, fissures, effondrements) et effets \cle{indirects} (incendies, glissements de terrain, raz-de-marée ou tsunami si le séisme se produit en mer).
\item \textbf{Conclure} : les manifestations d'un séisme sont d'autant plus importantes que le séisme est fort et que les constructions sont fragiles.
\end{enumerate}
\textbf{Réflexe} : toujours citer le document (témoignage, photo) pour justifier chaque manifestation relevée.
\end{methode}

\begin{exoresolu}[Les manifestations d'un séisme]
\textbf{Énoncé.} Après un séisme survenu dans une localité, un témoin déclare : « Le sol a tremblé pendant quelques secondes, des fissures sont apparues sur les murs des maisons, plusieurs cases se sont effondrées, une fissure s'est ouverte dans le sol et un incendie s'est déclaré à la suite de la rupture d'une canalisation de gaz. »
\begin{enumerate}
\item Définis le terme séisme.
\item Relève dans ce témoignage trois manifestations directes du séisme.
\item Cite une manifestation indirecte et explique pourquoi elle est dite indirecte.
\end{enumerate}
\tcblower
\textbf{Solution détaillée.}
\begin{enumerate}
\item Un \cle{séisme} (ou tremblement de terre) est une série de vibrations brèves et brusques du sol, provoquées par la libération soudaine d'énergie accumulée dans les roches de la lithosphère.
\item Les manifestations directes relevées dans le témoignage sont :
\begin{itemize}
\item les \textbf{vibrations du sol} (« le sol a tremblé pendant quelques secondes ») ;
\item les \textbf{fissures dans les murs} des maisons ;
\item l'\textbf{effondrement des cases} ;
\item l'\textbf{ouverture d'une fissure dans le sol}.
\end{itemize}
Ces manifestations sont dites \emph{directes} car elles résultent immédiatement des vibrations du sol provoquées par les ondes sismiques.
\item La manifestation indirecte est \textbf{l'incendie}. Elle est dite indirecte car elle ne résulte pas directement des vibrations du sol : elle est une \emph{conséquence} d'un dégât direct (la rupture de la canalisation de gaz provoquée par les secousses).
\end{enumerate}
\textbf{Conclusion.} Un séisme se manifeste par des vibrations du sol, des fissures et des effondrements (effets directs), qui peuvent entraîner des effets indirects comme les incendies.
\end{exoresolu}

\courssub{Méthode 2 : Distinguer foyer et épicentre sur un schéma}

\begin{methode}[Identifier le foyer et l'épicentre d'un séisme]
Pour exploiter un schéma de coupe de la lithosphère montrant l'origine d'un séisme :
\begin{enumerate}
\item \textbf{Repérer le foyer} : c'est la zone \emph{en profondeur}, à l'intérieur de la lithosphère, où le séisme prend naissance (rupture brutale des roches le long d'une faille).
\item \textbf{Repérer l'épicentre} : c'est le point de la \emph{surface} du sol situé à la verticale du foyer.
\item \textbf{Traduire la propagation} : les ondes sismiques partent du foyer dans toutes les directions ; on peut les représenter par des cercles concentriques ou des flèches.
\item \textbf{Raisonner sur les dégâts} : les dégâts sont maximaux à l'épicentre et diminuent quand on s'en éloigne.
\end{enumerate}
\textbf{À retenir} : foyer = en profondeur ; épicentre = en surface, à la verticale du foyer. Ne jamais confondre les deux dans la rédaction.
\end{methode}

\begin{exoresolu}[Foyer, épicentre et propagation des ondes]
\textbf{Énoncé.} Un séisme se produit à \SI{10}{km} de profondeur sous une région habitée.
\begin{enumerate}
\item Définis les termes foyer et épicentre.
\item Schématise par un dessin légendé la position du foyer, de l'épicentre et la propagation des ondes sismiques.
\item Deux villages A et B sont situés respectivement à \SI{5}{km} et à \SI{60}{km} de l'épicentre. Dans quel village les dégâts seront-ils les plus importants ? Justifie.
\end{enumerate}
\tcblower
\textbf{Solution détaillée.}
\begin{enumerate}
\item Le \cle{foyer} est la zone située en profondeur dans la lithosphère où le séisme prend naissance (rupture brutale des roches). L'\cle{épicentre} est le point de la surface du sol situé à la verticale du foyer.
\item Schéma :

\begin{popfigure}
\begin{center}
\begin{tikzpicture}[scale=1]
\draw[thick, popdark] (-4,0) -- (4,0) node[right] {\small Surface du sol};
\fill[popink] (0,-2) circle (0.12);
\node[below, popink] at (0,-2.15) {\small Foyer (10 km de profondeur)};
\fill[poporange] (0,0) circle (0.12);
\node[above, poporange] at (0,0.1) {\small Épicentre};
\draw[dashed, popmuted] (0,0) -- (0,-2);
\draw[popblue, thick] (0,-2) circle (0.7);
\draw[popblue] (0,-2) circle (1.3);
\draw[popblue!60] (0,-2) circle (1.9);
\node[popblue] at (2.6,-2.6) {\small Ondes sismiques};
\end{tikzpicture}
\end{center}
\legende{Schéma de l'origine d'un séisme : le foyer en profondeur, l'épicentre à la surface à sa verticale, et les ondes sismiques qui se propagent dans toutes les directions (cercles concentriques).}
\end{popfigure}

\item Les dégâts seront les plus importants au \textbf{village A}. En effet, les ondes sismiques perdent de leur énergie en se propageant : plus on s'éloigne de l'épicentre, plus les vibrations s'atténuent. Le village A, situé à seulement \SI{5}{km} de l'épicentre, recevra des vibrations beaucoup plus fortes que le village B situé à \SI{60}{km}.
\end{enumerate}
\textbf{Conclusion.} L'épicentre est la zone de surface où les dégâts sont maximaux ; leur intensité diminue avec la distance à l'épicentre.
\end{exoresolu}

\courssub{Méthode 3 : Expliquer l'origine d'un séisme (faille et déplacement des roches)}

\begin{methode}[Rédiger une explication de l'origine d'un séisme]
Pour expliquer comment naît un séisme, rédiger une démarche en étapes :
\begin{enumerate}
\item \textbf{Situer le contexte} : les roches de la lithosphère sont soumises à des forces (contraintes) liées aux mouvements internes de la Terre.
\item \textbf{Décrire l'accumulation d'énergie} : les roches se déforment lentement et accumulent de l'énergie, comme un bâton qu'on plie.
\item \textbf{Décrire la rupture} : lorsque la contrainte dépasse la résistance des roches, celles-ci se rompent brutalement le long d'une \cle{faille} (fracture de la roche avec déplacement des deux compartiments).
\item \textbf{Conclure} : l'énergie accumulée est libérée soudainement sous forme d'\cle{ondes sismiques} qui se propagent et font vibrer le sol : c'est le séisme.
\end{enumerate}
\textbf{Réflexe} : utiliser le vocabulaire exact : contrainte, déformation, rupture, faille, libération d'énergie, ondes sismiques.
\end{methode}

\begin{experience}[Modélisation de l'origine d'un séisme : la règle pliée]
\textbf{Matériel} : une règle en plastique souple.\\
\textbf{Protocole} : maintenir la règle par ses deux extrémités et la plier progressivement.\\
\textbf{Observations} : la règle se déforme d'abord lentement ; au-delà d'une certaine courbure, elle se rompt brutalement en émettant un bruit sec et en projetant de petits fragments.\\
\textbf{Interprétation} : la règle joue le rôle des roches de la lithosphère ; les forces exercées jouent le rôle des contraintes ; la rupture brutale avec libération d'énergie (bruit, fragments projetés) correspond à la rupture des roches le long d'une faille, qui libère l'énergie accumulée sous forme d'ondes sismiques.\\
\textbf{Conclusion} : un séisme naît de la rupture brutale de roches préalablement déformées par des contraintes.
\end{experience}

\begin{exoresolu}[L'expérience du bâton plié et l'origine des séismes]
\textbf{Énoncé.} Au laboratoire, on plie progressivement une règle en plastique maintenue par ses deux extrémités. La règle se déforme puis, au-delà d'une certaine courbure, elle se rompt brutalement en émettant un bruit sec et en projetant de petits fragments.
\begin{enumerate}
\item À quelle partie de la Terre cette expérience fait-elle penser ?
\item En utilisant cette analogie, explique l'origine d'un séisme.
\item Quel nom donne-t-on à la rupture des roches accompagnée d'un déplacement ?
\end{enumerate}
\tcblower
\textbf{Solution détaillée.}
\begin{enumerate}
\item Cette expérience est un modèle qui fait penser aux \textbf{roches de la lithosphère} (la partie externe et rigide de la Terre) soumises à des forces.
\item Explication de l'origine d'un séisme par analogie :
\begin{itemize}
\item les forces exercées sur la règle correspondent aux \textbf{contraintes} que subissent les roches de la lithosphère à cause des mouvements internes de la Terre ;
\item la déformation progressive de la règle correspond à la \textbf{déformation lente des roches}, qui accumulent de l'énergie ;
\item la rupture brutale de la règle, avec bruit et projection de fragments, correspond à la \textbf{rupture soudaine des roches} lorsque la contrainte dépasse leur résistance : l'énergie accumulée est libérée sous forme d'\textbf{ondes sismiques} qui font vibrer le sol.
\end{itemize}
\item La rupture des roches accompagnée d'un déplacement des deux compartiments s'appelle une \cle{faille}.
\end{enumerate}
\textbf{Conclusion.} Un séisme résulte de la rupture brutale des roches de la lithosphère, déformées par des contraintes, le long d'une faille : l'énergie libérée se propage sous forme d'ondes sismiques.
\end{exoresolu}

\courssub{Méthode 4 : Exploiter les échelles d'évaluation d'un séisme}

\begin{methode}[Distinguer intensité et magnitude]
Pour évaluer un séisme, il faut bien distinguer les deux notions :
\begin{enumerate}
\item \textbf{Identifier ce qu'on mesure} : les effets et dégâts observés (intensité) ou l'énergie libérée au foyer (magnitude).
\item \textbf{Intensité} : elle se mesure avec l'\cle{échelle de Mercalli} (ou MSK), graduée en degrés de I à XII, à partir de l'observation des dégâts et des effets ressentis par la population. Elle \emph{varie} d'un lieu à l'autre pour un même séisme.
\item \textbf{Magnitude} : elle se mesure avec l'\cle{échelle de Richter}, à partir des enregistrements des sismographes. Elle a \emph{une seule valeur} pour un séisme donné.
\item \textbf{Conclure} en comparant : l'intensité dépend du lieu d'observation ; la magnitude caractérise le séisme lui-même.
\end{enumerate}
\textbf{Piège à éviter} : ne jamais écrire « degrés de l'échelle de Richter » : l'échelle de Richter mesure la magnitude, pas des degrés de dégâts.
\end{methode}

\begin{aretenir}[Intensité ou magnitude ?]
\begin{center}
\begin{tabularx}{\linewidth}{|l|X|X|}
\hline
\rowcolor{popblueL} \textbf{Critère} & \textbf{Intensité} & \textbf{Magnitude} \\
\hline
Ce qui est mesuré & Effets et dégâts observés en un lieu & Énergie libérée au foyer \\
\hline
Échelle utilisée & Échelle de Mercalli (MSK), degrés I à XII & Échelle de Richter \\
\hline
Moyen d'obtention & Observation des dégâts, enquêtes & Enregistrements du sismographe (sismogramme) \\
\hline
Nombre de valeurs & Variable d'un lieu à l'autre & Une seule valeur par séisme \\
\hline
\end{tabularx}
\end{center}
\end{aretenir}

\begin{exoresolu}[Intensité et magnitude d'un même séisme]
\textbf{Énoncé.} Un séisme de magnitude 6,5 sur l'échelle de Richter frappe une région. Les observations recueillies sont les suivantes :
\begin{center}
\begin{tabularx}{\linewidth}{|l|X|}
\hline
\rowcolor{popblueL} \textbf{Lieu} & \textbf{Effets observés} \\
\hline
Ville A (proche de l'épicentre) & Nombreuses maisons détruites, fissures dans le sol \\
\hline
Ville B (à \SI{80}{km}) & Objets renversés, quelques fissures dans les murs \\
\hline
Ville C (à \SI{200}{km}) & Vibrations ressenties, aucun dégât \\
\hline
\end{tabularx}
\end{center}
\begin{enumerate}
\item Que signifie la valeur 6,5 annoncée ? Avec quel appareil est-elle obtenue ?
\item L'intensité du séisme est-elle la même dans les trois villes ? Justifie.
\item Quelle échelle utilise-t-on pour mesurer l'intensité ?
\end{enumerate}
\tcblower
\textbf{Solution détaillée.}
\begin{enumerate}
\item La valeur 6,5 est la \cle{magnitude} du séisme : elle exprime l'importance de l'énergie libérée au foyer. Elle est déterminée à partir des enregistrements des vibrations du sol réalisés par un \cle{sismographe} (le document obtenu s'appelle un \cle{sismogramme}).
\item Non, l'intensité n'est pas la même dans les trois villes. L'\cle{intensité} mesure les effets et les dégâts en un lieu donné : elle est \textbf{maximale près de l'épicentre} (ville A : maisons détruites) et \textbf{décroît quand on s'en éloigne} (ville B : fissures ; ville C : simples vibrations sans dégâts).
\item L'intensité se mesure avec l'\cle{échelle de Mercalli} (échelle MSK), graduée en degrés de I (séisme non ressenti) à XII (destruction totale), à partir de l'observation des dégâts.
\end{enumerate}
\textbf{Conclusion.} Un même séisme possède une seule magnitude (énergie libérée, mesurée au sismographe) mais des intensités différentes selon les lieux (dégâts observés, échelle de Mercalli).
\end{exoresolu}

\courssub{Méthode 5 : Exploiter une carte de répartition mondiale des séismes}

\begin{methode}[Analyser la localisation mondiale des séismes]
Pour exploiter une carte montrant la répartition des épicentres des séismes dans le monde :
\begin{enumerate}
\item \textbf{Observer} la carte : les épicentres ne sont pas répartis au hasard ; ils forment des \textbf{alignements} (zones étroites et allongées).
\item \textbf{Localiser} ces alignements : le long des \textbf{côtes de l'océan Pacifique} (« ceinture de feu »), au niveau des \textbf{dorsales océaniques} (reliefs sous-marins au milieu des océans) et dans certaines chaînes de montagnes (Alpes, Himalaya).
\item \textbf{Comparer} avec la carte des volcans : les séismes et les volcans se situent souvent dans les \emph{mêmes zones}.
\item \textbf{Conclure} : les séismes se produisent dans des zones bien précises du globe, appelées \cle{zones sismiques}, qui correspondent aux limites des grandes plaques rigides de la lithosphère.
\end{enumerate}
\textbf{Réflexe} : toujours appuyer la conclusion sur des exemples précis lus sur la carte (Pacifique, dorsales, Himalaya).
\end{methode}

\begin{exoresolu}[Lecture d'une carte de répartition des séismes]
\textbf{Énoncé.} Sur une carte mondiale, on a reporté par des points les épicentres des séismes enregistrés pendant plusieurs années. On constate que ces points forment des alignements : le long des côtes de l'océan Pacifique (Japon, Chili, Californie), au milieu de l'océan Atlantique et dans la région de l'Himalaya.
\begin{enumerate}
\item Que constate-t-on concernant la répartition des séismes à la surface du globe ?
\item Cite deux régions du monde où les séismes sont fréquents.
\item Les séismes et les éruptions volcaniques se produisent-ils dans les mêmes zones ? Que peut-on en conclure ?
\end{enumerate}
\tcblower
\textbf{Solution détaillée.}
\begin{enumerate}
\item On constate que les séismes ne sont \textbf{pas répartis uniformément} à la surface du globe : leurs épicentres se concentrent dans des \textbf{zones étroites et allongées}, formant des alignements appelés \cle{zones sismiques}.
\item Deux régions où les séismes sont fréquents : les \textbf{bordures de l'océan Pacifique} (Japon, Chili, Californie — la « ceinture de feu du Pacifique ») et la région de l'\textbf{Himalaya}. On peut aussi citer le milieu de l'océan Atlantique (dorsale médio-atlantique).
\item Oui, la comparaison des cartes montre que les séismes et les volcans se situent souvent dans les \textbf{mêmes zones}. On en conclut que ces phénomènes ont une \textbf{origine commune} : ils sont liés aux mouvements internes de la Terre et se produisent aux limites des grandes plaques rigides de la lithosphère.
\end{enumerate}
\textbf{Conclusion.} Les séismes sont localisés dans des zones précises du globe (ceinture de feu du Pacifique, dorsales océaniques, chaînes de montagnes), qui correspondent aux zones de fragilité de la lithosphère.
\end{exoresolu}

\begin{savaistu}[Et au Cameroun ?]
Le Cameroun n'est pas à l'abri des séismes : la région du \textbf{Mont Cameroun}, volcan actif situé dans la ligne volcanique du Cameroun, connaît régulièrement de faibles secousses enregistrées par des sismographes. Des séismes de faible magnitude ont aussi été ressentis dans d'autres régions du pays. C'est pourquoi les normes de construction et l'information des populations concernent aussi notre pays.
\end{savaistu}

\courssub{Méthode 6 : Proposer des mesures de prévision et de prévention}

\begin{methode}[Répondre à une situation concrète de prévention des séismes]
Face à une situation de la vie courante (construction d'un village, aménagement d'une ville en zone sismique) :
\begin{enumerate}
\item \textbf{Rappeler la limite de la prévision} : on sait localiser les zones sismiques, mais on ne peut \emph{pas} prévoir la date exacte d'un séisme. La prévision repose sur la surveillance (sismographes, mesure des déformations du sol, observation des signes annonciateurs comme la variation du niveau des nappes d'eau souterraines ou le comportement anormal des animaux).
\item \textbf{Proposer des mesures de prévention} :
\begin{itemize}
\item construire des bâtiments \cle{parasismiques} (fondations profondes, matériaux souples et légers, structures renforcées) ;
\item respecter les normes de construction en zone sismique ;
\item éviter de construire sur les failles connues et sur les pentes instables ;
\item informer et éduquer la population (conduites à tenir : s'éloigner des bâtiments, se mettre sous une table, couper le gaz et l'électricité).
\end{itemize}
\item \textbf{Justifier} chaque mesure par le risque qu'elle réduit.
\end{enumerate}
\textbf{Réflexe} : bien distinguer \emph{prévision} (anticiper le séisme) et \emph{prévention} (limiter les dégâts).
\end{methode}

\begin{exoresolu}[Aménager une ville en zone sismique]
\textbf{Énoncé.} Une ville se développe dans une région traversée par une faille active, où plusieurs séismes ont déjà été enregistrés. Le maire te consulte pour l'aménagement d'un nouveau quartier.
\begin{enumerate}
\item Peut-on prévoir la date exacte du prochain séisme dans cette région ? Justifie ta réponse.
\item Propose trois conseils à donner au maire pour limiter les dégâts d'un séisme futur, en justifiant chacun d'eux.
\item Indique deux comportements à adopter par les habitants pendant une secousse.
\end{enumerate}
\tcblower
\textbf{Solution détaillée.}
\begin{enumerate}
\item Non, on ne peut pas prévoir la date exacte d'un séisme. La science permet seulement d'\textbf{identifier les zones sismiques} (ici, la faille active) et de \textbf{surveiller} certains signes (déformations du sol, petites secousses, variations du niveau des eaux souterraines), mais l'instant de la rupture des roches reste imprévisible.
\item Trois conseils de prévention :
\begin{itemize}
\item \textbf{Construire des bâtiments parasismiques} (fondations profondes, matériaux souples et légers) : ils résistent mieux aux vibrations et s'effondrent moins facilement ;
\item \textbf{ne pas construire sur la faille ni sur les pentes instables} : cela évite les destructions directes par la rupture du sol et les glissements de terrain ;
\item \textbf{informer et entraîner la population} (exercices d'évacuation, affichage des consignes) : une population préparée réagit correctement et limite le nombre de victimes.
\end{itemize}
\item Pendant une secousse, les habitants doivent : \textbf{s'éloigner des bâtiments} et des lignes électriques s'ils sont dehors ; \textbf{se mettre sous une table solide} ou contre un mur porteur s'ils sont à l'intérieur, et \textbf{couper le gaz et l'électricité} dès que possible pour éviter les incendies.
\end{enumerate}
\textbf{Conclusion.} Si la prévision exacte des séismes est impossible, la prévention (constructions parasismiques, choix des sites, éducation des populations) permet de réduire considérablement les dégâts et les pertes humaines.
\end{exoresolu}

\begin{attention}[Les 5 erreurs qui coûtent des points au BEPC]
\begin{enumerate}
\item \textbf{Confondre foyer et épicentre.} Le foyer est \emph{en profondeur} (lieu de la rupture des roches) ; l'épicentre est \emph{à la surface}, à la verticale du foyer. Sur un schéma, place-les correctement et légende-les.
\item \textbf{Confondre intensité et magnitude.} L'intensité (échelle de Mercalli, degrés I à XII) mesure les \emph{dégâts observés} et varie d'un lieu à l'autre ; la magnitude (échelle de Richter) mesure l'\emph{énergie libérée} au foyer et a une seule valeur par séisme. Ne dis jamais « degrés de Richter ».
\item \textbf{Oublier l'appareil de mesure.} Les vibrations d'un séisme s'enregistrent avec un \cle{sismographe} (le document obtenu est un sismogramme). Cite-le chaque fois qu'on te demande comment on mesure un séisme.
\item \textbf{Affirmer qu'on peut prévoir la date d'un séisme.} C'est faux : on peut seulement localiser les zones sismiques et surveiller des signes annonciateurs. Écris toujours que la prévision exacte (date et heure) est \emph{impossible} actuellement.
\item \textbf{Rédiger sans justifier à partir des documents.} Quand un énoncé fournit un témoignage, un tableau ou une carte, cite les données précises (distances, dégâts observés, régions) pour appuyer ta réponse : une affirmation non justifiée par le document perd des points.
\end{enumerate}
\end{attention}

\newpage

\coursec{Exercices}

\courssub{Je vérifie mes connaissances}

\begin{exercice}[QCM : les bons réflexes]
Pour chaque question, une seule réponse est exacte. Recopie la lettre correspondante sur ta copie.

\textbf{1.} Le point situé à la surface de la Terre, à la verticale du foyer, s'appelle :
\begin{itemize}
\item[a)] l'hypocentre ;
\item[b)] l'épicentre ;
\item[c)] le sismographe.
\end{itemize}

\textbf{2.} Les ondes sismiques qui arrivent en premier à une station d'enregistrement sont :
\begin{itemize}
\item[a)] les ondes S ;
\item[b)] les ondes de surface ;
\item[c)] les ondes P.
\end{itemize}

\textbf{3.} L'échelle de Mercalli mesure :
\begin{itemize}
\item[a)] la magnitude d'un séisme ;
\item[b)] l'intensité d'un séisme, d'après les dégâts observés ;
\item[c)] la profondeur du foyer.
\end{itemize}

\textbf{4.} La plupart des séismes se produisent :
\begin{itemize}
\item[a)] au centre des plaques lithosphériques ;
\item[b)] uniquement sous les océans ;
\item[c)] le long des limites des plaques lithosphériques.
\end{itemize}
\end{exercice}

\begin{exercice}[Vrai ou faux ?]
Pour chaque affirmation, indique si elle est vraie ou fausse, puis justifie ta réponse en une ou deux phrases.

\textbf{1.} Un séisme est une secousse brusque du sol provoquée par une libération brutale d'énergie dans la lithosphère.

\textbf{2.} L'intensité d'un séisme est la même en tous les points de la région touchée.

\textbf{3.} Les ondes S se propagent plus vite que les ondes P.

\textbf{4.} On peut prévoir la date et l'heure exactes d'un séisme grâce aux sismographes.

\textbf{5.} Le Cameroun possède des zones exposées aux séismes, notamment le long de la ligne volcanique du Mont Cameroun.
\end{exercice}

\begin{exercice}[Définitions à compléter]
Recopie et complète les définitions suivantes avec les mots ou groupes de mots qui conviennent.

\textbf{1.} Un \cle{séisme} est un \ldots\ldots\ldots du sol, provoqué par la rupture brutale des roches en profondeur.

\textbf{2.} Le \cle{foyer} est la zone \ldots\ldots\ldots où naît le séisme.

\textbf{3.} Un \cle{sismographe} est un appareil qui \ldots\ldots\ldots les vibrations du sol ; le document obtenu s'appelle un \ldots\ldots\ldots .

\textbf{4.} La \cle{magnitude} mesure \ldots\ldots\ldots libérée au foyer, sur l'échelle de \ldots\ldots\ldots .
\end{exercice}

\begin{exercice}[Calcul direct : distance à l'épicentre]
Une station sismique enregistre l'arrivée des ondes P à 14 h 02 min 10 s et celle des ondes S à 14 h 02 min 50 s.

\textbf{1.} Calcule l'écart S--P (durée séparant l'arrivée des deux ondes).

\textbf{2.} Sachant que, dans la région étudiée, un écart S--P de 10 secondes correspond à une distance d'environ \SI{80}{km} entre la station et l'épicentre, calcule la distance de la station à l'épicentre.

\textbf{3.} Cette seule station suffit-elle à localiser précisément l'épicentre ? Justifie.
\end{exercice}

\courssub{Je m'entraîne}

\begin{exercice}[Lecture d'un sismogramme]
Le document ci-dessous représente le sismogramme enregistré à la station de Kribi lors d'un séisme. On y distingue deux trains d'ondes : le premier arrive à $t_1 = 25$ s, le second à $t_2 = 65$ s.

\textbf{1.} Identifie, en justifiant, le train d'ondes correspondant aux ondes P et celui correspondant aux ondes S.

\textbf{2.} Mesure l'écart S--P sur ce sismogramme.

\textbf{3.} À l'aide de l'abaque suivante (écart S--P en fonction de la distance à l'épicentre), détermine la distance entre la station de Kribi et l'épicentre du séisme.

\begin{center}
\begin{tabular}{|l|c|c|c|c|c|}
\hline
\rowcolor{popblueL} Écart S--P (s) & 10 & 20 & 30 & 40 & 50 \\
\hline
Distance à l'épicentre (km) & 80 & 160 & 240 & 320 & 400 \\
\hline
\end{tabular}
\end{center}

\textbf{4.} Explique pourquoi les ondes P et les ondes S n'arrivent pas en même temps à la station.
\end{exercice}

\begin{exercice}[Intensité ou magnitude ?]
Voici des affirmations entendues après un séisme. Classe-les dans un tableau à deux colonnes selon qu'elles concernent l'\cle{intensité} (échelle de Mercalli) ou la \cle{magnitude} (échelle de Richter). Corrige ensuite les affirmations contenant une erreur de vocabulaire.

\textbf{a)} « À Douala, les murs de plusieurs maisons se sont fissurés : on estime le séisme au degré VI. »

\textbf{b)} « Le séisme a libéré une énergie correspondant à une magnitude de 6,2. »

\textbf{c)} « La magnitude du séisme était plus forte en ville qu'en brousse, car il y a eu plus de dégâts. »

\textbf{d)} « L'intensité du séisme, mesurée au sismographe, vaut 5,8. »

\textbf{e)} « Dans le village le plus proche de l'épicentre, presque tous les bâtiments ont été détruits : degré IX. »
\end{exercice}

\begin{exercice}[Localisation d'un épicentre par triangulation]
Un séisme a été enregistré par trois stations camerounaises. Les distances à l'épicentre, déduites des écarts S--P, sont les suivantes :

\begin{center}
\begin{tabular}{|l|c|}
\hline
\rowcolor{popblueL} Station & Distance à l'épicentre (km) \\
\hline
Yaoundé & 200 \\
\hline
Bafoussam & 150 \\
\hline
Ngaoundéré & 300 \\
\hline
\end{tabular}
\end{center}

\textbf{1.} Sur une carte du Cameroun à l'échelle $1$ cm pour $100$ km, convertis chaque distance en centimètres sur la carte.

\textbf{2.} Décris la construction à réaliser (instruments, tracés) pour localiser l'épicentre.

\textbf{3.} Pourquoi faut-il au moins trois stations pour obtenir une position précise de l'épicentre ?
\end{exercice}

\begin{exercice}[Le rebond élastique et les failles]
On réalise l'expérience suivante : on plie lentement une règle en plastique maintenue aux deux extrémités, jusqu'à ce qu'elle se rompe brutalement en vibrant.

\textbf{1.} À quelle structure géologique la règle pliée est-elle comparée ? À quoi correspond sa rupture brutale ?

\textbf{2.} À l'aide du modèle du \cle{rebond élastique}, explique pourquoi les séismes se produisent le long des \cle{failles actives}.

\textbf{3.} La faille de Foumban (région de l'Ouest camerounais) est une faille active. Quelle précaution d'aménagement cela impose-t-il aux habitants de cette région ?
\end{exercice}

\courssub{Je me prépare au BEPC}

\begin{exercice}[Type BEPC : le séisme de la région du Mont Cameroun]
Un séisme est enregistré par la station sismologique de Buéa, située sur les flancs du Mont Cameroun. Le sismogramme obtenu montre l'arrivée des ondes P à 9 h 15 min 12 s et celle des ondes S à 9 h 15 min 36 s. On donne l'abaque de la région : un écart S--P de 8 s correspond à une distance épicentrale de \SI{64}{km}.

\textbf{1.} Définis les termes suivants : séisme, foyer, épicentre.

\textbf{2.} Calcule l'écart S--P enregistré à la station de Buéa.

\textbf{3.} Détermine, à l'aide de l'abaque, la distance entre la station et l'épicentre.

\textbf{4.} Deux autres stations donnent les distances suivantes : Limbé : \SI{40}{km} ; Kumba : \SI{90}{km}. Explique, en décrivant la construction, comment localiser l'épicentre sur une carte à l'échelle $1$ cm pour $20$ km.

\textbf{5.} La magnitude du séisme est de 5,1 sur l'échelle de Richter. À Buéa, les habitants ont ressenti une forte secousse et quelques fissures sont apparues sur les murs (degré V de Mercalli). Distingue clairement, à partir de cet exemple, les notions de magnitude et d'intensité.

\textbf{6.} Cite deux mesures de prévention que les autorités de la région du Sud-Ouest peuvent prendre pour limiter les dégâts des séismes futurs.
\end{exercice}

\begin{exercice}[Type BEPC : la répartition mondiale des séismes]
On dispose d'une carte mondiale indiquant la position des épicentres des séismes enregistrés pendant dix ans, ainsi que d'une carte des limites des plaques lithosphériques (dorsales océaniques, fosses océaniques, grandes failles). On constate que les épicentres forment des alignements étroits : le long de la ceinture de feu du Pacifique, le long des dorsales médio-océaniques et dans la chaîne alpine et himalayenne.

\textbf{1.} Rappelle ce qu'est un séisme et où naît un séisme.

\textbf{2.} Décris, en deux ou trois phrases, la répartition des épicentres à la surface du globe.

\textbf{3.} En comparant les deux cartes, dégage la relation qui existe entre la répartition des séismes et les limites des plaques lithosphériques.

\textbf{4.} Propose une explication : pourquoi les limites de plaques sont-elles des zones où les roches se déforment puis se rompent ?

\textbf{5.} Rédige une conclusion de cinq à huit lignes, justifiée par les documents, répondant à la question : « Les séismes se produisent-ils au hasard à la surface de la Terre ? »
\end{exercice}

\begin{exercice}[Type BEPC : protéger une commune exposée aux séismes]
La commune de Wum, dans la région du Nord-Ouest du Cameroun, est traversée par une faille active. Le maire te consulte, toi qui as suivi le cours sur les séismes, pour préparer un plan de protection de la population.

\textbf{1.} Explique au maire, en quelques lignes, pourquoi sa commune est exposée au risque sismique (utilise les notions de faille active et de rebond élastique).

\textbf{2.} Peut-on prévoir la date exacte du prochain séisme ? Justifie ta réponse et explique la différence entre \cle{prévision} et \cle{prévention}.

\textbf{3.} Rédige un texte d'environ dix lignes, adressé au maire, présentant au moins quatre mesures concrètes de prévention : règles de construction (bâtiments parasismiques), choix des zones d'habitation, information et entraînement de la population, organisation des secours.

\textbf{4.} Termine par deux conseils de comportement à adopter pendant une secousse, que tu proposeras d'afficher dans les écoles de la commune.
\end{exercice}



\newpage

\coursec{Corrigés des exercices}

\begin{corrige}[Exercice 1 — QCM : les bons réflexes]
\textbf{1.} Réponse b) : l'\cle{épicentre} est le point de la surface terrestre situé à la verticale du foyer. Le foyer (ou hypocentre) est en profondeur ; le sismographe est un appareil d'enregistrement, pas un point.

\textbf{2.} Réponse c) : les ondes P (ondes primaires) sont les plus rapides ; elles arrivent donc en premier à la station, avant les ondes S et les ondes de surface.

\textbf{3.} Réponse b) : l'échelle de Mercalli (de I à XII) évalue l'\cle{intensité} d'un séisme à partir des effets observés sur les personnes, les constructions et le paysage. La magnitude se mesure sur l'échelle de Richter.

\textbf{4.} Réponse c) : les séismes se concentrent le long des limites des plaques lithosphériques (dorsales, fosses, grandes failles), où les roches se déforment puis se rompent brutalement.
\end{corrige}

\begin{corrige}[Exercice 2 — Vrai ou faux ?]
\textbf{1.} Vrai. C'est la définition même du séisme : la rupture brutale des roches en profondeur libère l'énergie accumulée, qui se propage sous forme d'ondes et fait trembler le sol.

\textbf{2.} Faux. L'intensité est maximale près de l'épicentre et diminue généralement quand on s'en éloigne ; elle dépend aussi de la nature du sol et de la qualité des constructions.

\textbf{3.} Faux. Les ondes P sont plus rapides que les ondes S : c'est d'ailleurs l'écart entre leurs temps d'arrivée qui permet de calculer la distance à l'épicentre.

\textbf{4.} Faux. Le sismographe enregistre les vibrations pendant et après le séisme ; aucun appareil ne permet aujourd'hui de prévoir la date, l'heure et le lieu exacts d'un séisme. On peut seulement identifier les zones à risque et s'y préparer (prévention).

\textbf{5.} Vrai. La ligne volcanique camerounaise (Mont Cameroun, région du Sud-Ouest) et certaines failles actives (par exemple la faille de Foumban) sont des zones où des séismes ont été enregistrés.
\end{corrige}

\begin{corrige}[Exercice 3 — Définitions à compléter]
\textbf{1.} Un séisme est un \textbf{tremblement (une secousse brusque)} du sol, provoqué par la rupture brutale des roches en profondeur.

\textbf{2.} Le foyer est la zone \textbf{souterraine (en profondeur, dans la lithosphère)} où naît le séisme.

\textbf{3.} Un sismographe est un appareil qui \textbf{enregistre} les vibrations du sol ; le document obtenu s'appelle un \textbf{sismogramme}.

\textbf{4.} La magnitude mesure \textbf{l'énergie} libérée au foyer, sur l'échelle de \textbf{Richter}.

\begin{attention}[Points de vigilance]
Ne pas confondre sismographe (l'appareil) et sismogramme (le document enregistré) ; ne pas confondre magnitude (énergie, échelle de Richter) et intensité (dégâts, échelle de Mercalli).
\end{attention}
\end{corrige}

\begin{corrige}[Exercice 4 — Calcul direct : distance à l'épicentre]
\textbf{1.} Écart S--P $= \SI{50}{s} - \SI{10}{s} = \SI{40}{s}$.

\textbf{2.} La distance est proportionnelle à l'écart S--P (loi empirique ou abaque fourni) :
\[ d = \frac{\SI{40}{s}}{\SI{10}{s}} \times \SI{80}{km} = \SI{320}{km}. \]

\textbf{3.} Non. Une seule station donne seulement la distance : l'épicentre peut se trouver n'importe où sur le cercle de rayon \SI{320}{km} centré sur la station. Il faut au moins trois stations pour obtenir un point unique (triangulation).

\begin{attention}[Point de vigilance]
Bien convertir les horaires en secondes avant de soustraire, et ne pas oublier l'unité (km) dans le résultat final.
\end{attention}
\end{corrige}

\begin{corrige}[Exercice 5 — Lecture d'un sismogramme]
\textbf{1.} Le premier train d'ondes ($t_1 = \SI{25}{s}$) correspond aux ondes P, les plus rapides ; le second ($t_2 = \SI{65}{s}$) correspond aux ondes S, plus lentes.

\textbf{2.} Écart S--P $= t_2 - t_1 = \SI{65}{s} - \SI{25}{s} = \SI{40}{s}$.

\textbf{3.} D'après l'abaque, un écart de \SI{40}{s} correspond à une distance de \SI{320}{km} entre la station de Kribi et l'épicentre.

\textbf{4.} Les ondes P et les ondes S naissent en même temps au foyer, mais elles ne se propagent pas à la même vitesse dans les roches : les ondes P sont plus rapides que les ondes S. Plus la station est éloignée de l'épicentre, plus l'écart entre leurs temps d'arrivée est grand.
\end{corrige}

\begin{corrige}[Exercice 6 — Intensité ou magnitude ?]
\begin{center}
\begin{tabularx}{\linewidth}{|X|X|}
\hline
\rowcolor{popblueL} Intensité (Mercalli) & Magnitude (Richter) \\
\hline
a) degré VI à Douala (dégâts observés) & b) magnitude 6,2 (énergie libérée) \\
\hline
e) degré IX près de l'épicentre (destructions) & \\
\hline
\end{tabularx}
\end{center}

Corrections des erreurs de vocabulaire :

\textbf{c)} Faux : c'est l'\textbf{intensité} (les effets et les dégâts) qui varie d'un lieu à l'autre ; la magnitude est une valeur unique pour un séisme donné, car elle mesure l'énergie libérée au foyer.

\textbf{d)} Faux : la valeur 5,8 mesurée au sismographe est la \textbf{magnitude} ; l'intensité ne se mesure pas avec un appareil, elle s'estime d'après les dégâts observés sur le terrain.
\end{corrige}

\begin{corrige}[Exercice 7 — Localisation d'un épicentre par triangulation]
\textbf{1.} À l'échelle 1 cm pour 100 km : Yaoundé : 2 cm ; Bafoussam : 1,5 cm ; Ngaoundéré : 3 cm.

\textbf{2.} Matériel : carte, compas, règle, crayon. Pour chaque station, on place la pointe sèche du compas sur la station et on trace le cercle dont le rayon est la distance convertie à l'échelle. L'épicentre se trouve au point d'intersection des trois cercles.

\textbf{3.} Avec une seule station, l'épicentre peut être n'importe où sur un cercle ; avec deux stations, les deux cercles se coupent en deux points possibles ; il faut une troisième station pour lever le doute et obtenir un point unique.

\begin{attention}[Point de vigilance]
Vérifier la conversion à l'échelle avant de tracer ; un rayon mal converti décale l'intersection et fausse toute la localisation.
\end{attention}
\end{corrige}

\begin{corrige}[Exercice 8 — Le rebond élastique et les failles]
\textbf{1.} La règle pliée est comparée aux roches de part et d'autre d'une faille, déformées lentement par les mouvements des plaques lithosphériques. Sa rupture brutale correspond à la rupture des roches le long de la faille, c'est-à-dire au séisme.

\textbf{2.} Le long d'une faille active, les blocs rocheux se déforment lentement et accumulent de l'énergie élastique, comme la règle pliée. Quand la contrainte dépasse la résistance des roches, celles-ci se rompent brutalement : les blocs « rebondissent » vers leur position d'équilibre et l'énergie accumulée est libérée en quelques secondes sous forme d'ondes sismiques.

\textbf{3.} Il faut éviter de construire des habitations et des bâtiments publics (écoles, hôpitaux) directement sur le tracé de la faille, et appliquer des règles de construction parasismique (fondations solides, murs chaînés et renforcés, bâtiments peu élevés).
\end{corrige}

\begin{corrige}[Exercice 9 — Type BEPC : le séisme de la région du Mont Cameroun]
\textbf{1.} Un \cle{séisme} est un tremblement brusque du sol provoqué par la rupture brutale des roches en profondeur. Le \cle{foyer} est la zone souterraine où naît le séisme. L'\cle{épicentre} est le point de la surface terrestre situé à la verticale du foyer.

\textbf{2.} Écart S--P $= \SI{36}{s} - \SI{12}{s} = \SI{24}{s}$.

\textbf{3.} Distance $= \dfrac{\SI{24}{s}}{\SI{8}{s}} \times \SI{64}{km} = \SI{192}{km}$.

\textbf{4.} Conversion à l'échelle 1 cm pour 20 km : Buéa : 9,6 cm ; Limbé : 2 cm ; Kumba : 4,5 cm. Au compas, on trace trois cercles centrés sur chaque station avec ces rayons ; le point d'intersection des trois cercles est l'épicentre.

\textbf{5.} La magnitude (5,1) mesure l'énergie libérée au foyer : c'est une valeur unique pour ce séisme, calculée à partir des sismogrammes. L'intensité (degré V à Buéa) décrit les effets ressentis et les dégâts observés en un lieu donné : elle varie d'un endroit à l'autre et diminue en général quand on s'éloigne de l'épicentre.

\textbf{6.} Par exemple : imposer des normes de construction parasismique pour les bâtiments neufs ; informer et entraîner la population (exercices d'évacuation dans les écoles) ; interdire les constructions sur les failles actives connues ; organiser les secours (plans d'urgence, centres d'accueil).

\textbf{Barème indicatif (sur 10) :} définitions 2 pts ; calcul de l'écart S--P 1 pt ; distance avec unité 1,5 pt ; construction de triangulation 2 pts ; distinction magnitude/intensité 2 pts ; mesures de prévention 1,5 pt.

\begin{attention}[Points de vigilance]
Citer l'unité à chaque résultat ; ne pas écrire que l'intensité est « mesurée au sismographe » ; pour la question 4, décrire explicitement l'usage du compas et l'intersection des trois cercles.
\end{attention}
\end{corrige}

\begin{corrige}[Exercice 10 — Type BEPC : la répartition mondiale des séismes]
\textbf{1.} Un séisme est un tremblement brusque du sol dû à la rupture brutale des roches ; il naît au foyer, en profondeur dans la lithosphère.

\textbf{2.} Les épicentres ne sont pas répartis uniformément : ils forment des alignements étroits (ceinture de feu du Pacifique, dorsales médio-océaniques, chaînes alpine et himalayenne), tandis que de vastes régions en sont presque dépourvues.

\textbf{3.} Les alignements d'épicentres coïncident avec les limites des plaques lithosphériques : les séismes se produisent essentiellement aux frontières des plaques.

\textbf{4.} Aux limites des plaques, celles-ci s'écartent, se rapprochent ou coulissent l'une contre l'autre. Ces mouvements déforment lentement les roches, qui accumulent de l'énergie élastique ; lorsque la contrainte dépasse leur résistance, les roches se rompent brutalement le long des failles : c'est le séisme (rebond élastique).

\textbf{5.} Conclusion attendue : les séismes ne se produisent pas au hasard. La carte montre que les épicentres dessinent des zones étroites et régulières qui correspondent aux limites des plaques lithosphériques. Là, les mouvements des plaques déforment les roches jusqu'à leur rupture brutale, qui libère l'énergie accumulée. Au contraire, l'intérieur des plaques, loin des limites, est beaucoup moins touché. Connaître les limites des plaques permet donc d'identifier les zones à risque sismique et d'y organiser la prévention.

\textbf{Barème indicatif (sur 10) :} définition 1,5 pt ; description de la carte 2 pts ; relation séismes--limites de plaques 2 pts ; explication par le rebond élastique 2 pts ; conclusion rédigée et justifiée 2,5 pts.

\begin{attention}[Points de vigilance]
La conclusion doit citer les documents (alignements observés) et répondre explicitement à la question posée ; éviter d'affirmer que les séismes ne se produisent « que » dans les océans.
\end{attention}
\end{corrige}

\begin{corrige}[Exercice 11 — Type BEPC : protéger une commune exposée aux séismes]
\textbf{1.} La commune est traversée par une faille active : les blocs rocheux de part et d'autre se déforment lentement sous l'effet des mouvements de la lithosphère et accumulent de l'énergie élastique. Lorsque la contrainte dépasse la résistance des roches, celles-ci se rompent brutalement (rebond élastique) : l'énergie libérée se propage sous forme d'ondes sismiques qui font trembler le sol de la commune.

\textbf{2.} Non : aucune méthode ne permet aujourd'hui de prévoir la date, l'heure et le lieu exacts d'un séisme. La \cle{prévision} consiste à identifier les zones à risque (failles actives) et à estimer la probabilité d'un séisme ; la \cle{prévention} consiste à agir avant le séisme (constructions adaptées, information, secours) pour limiter les dégâts, puisque le séisme lui-même ne peut être ni empêché ni daté.

\textbf{3.} Exemple de texte : « Monsieur le Maire, votre commune étant traversée par une faille active, je vous propose quatre mesures. Premièrement, imposez des règles de construction parasismique : fondations profondes et solidaires, murs chaînés et renforcés, bâtiments peu élevés et légers. Deuxièmement, réservez le tracé de la faille et ses abords à des espaces verts ou des terrains de sport, et interdisez-y les habitations, les écoles et les hôpitaux. Troisièmement, informez régulièrement la population et organisez des exercices d'évacuation dans les écoles et les marchés, afin que chacun connaisse les bons gestes. Quatrièmement, préparez l'organisation des secours : un plan d'urgence, des équipes formées, du matériel de premiers soins et des lieux d'accueil identifiés à l'avance. Ces mesures n'empêcheront pas le séisme, mais elles sauveront de nombreuses vies. »

\textbf{4.} Conseils à afficher : « Pendant la secousse, reste calme, abrite-toi sous une table solide ou contre un mur porteur, loin des fenêtres et des objets qui peuvent tomber. » — « Si tu es dehors, éloigne-toi des bâtiments, des poteaux électriques et des arbres, et reste à découvert jusqu'à la fin des secousses. »

\textbf{Barème indicatif (sur 10) :} explication du risque (faille active, rebond élastique) 2 pts ; distinction prévision/prévention 2 pts ; texte avec au moins quatre mesures concrètes 4 pts ; deux conseils de comportement 2 pts.

\begin{attention}[Points de vigilance]
Ne pas écrire qu'on peut « empêcher » un séisme ; les mesures doivent être concrètes et adaptées à une commune camerounaise ; le texte doit être rédigé en phrases complètes, pas en simple liste.
\end{attention}
\end{corrige}

\newpage

\coursec{Fiche bilan}

\begin{popfigure}
\begin{tikzpicture}[
  mindmap,
  grow cyclic,
  every node/.style={concept, execute at begin node=\hskip0pt},
  concept color=popblue!80,
  level 1/.style={level distance=4.5cm, sibling angle=72},
  level 2/.style={level distance=3cm, sibling angle=45},
  branche1/.style={concept color=popblue, font=\small\bfseries},
  branche2/.style={concept color=poporange, font=\small\bfseries},
  branche3/.style={concept color=popgreen, font=\small\bfseries},
  branche4/.style={concept color=poppurple, font=\small\bfseries},
  branche5/.style={concept color=poppink, font=\small\bfseries},
  sousbranche/.style={concept color=popmuted!60, font=\scriptsize, level distance=2.5cm}
]
\node[concept, text width=2.5cm, align=center, font=\large\bfseries] {Séismes}
  child[branche1] { node[concept, text width=2.2cm, align=center] {Manifestations}
    child[sousbranche] { node[concept, text width=1.8cm, align=center] {Secousses du sol} }
    child[sousbranche] { node[concept, text width=1.8cm, align=center] {Dégâts matériels} }
    child[sousbranche] { node[concept, text width=1.8cm, align=center] {Tsunamis (si mer)} }
  }
  child[branche2] { node[concept, text width=2.2cm, align=center] {Origine}
    child[sousbranche] { node[concept, text width=1.8cm, align=center] {Foyer (hypocentre)} }
    child[sousbranche] { node[concept, text width=1.8cm, align=center] {Épicentre} }
    child[sousbranche] { node[concept, text width=1.8cm, align=center] {Ondes P \& S} }
    child[sousbranche] { node[concept, text width=1.8cm, align=center] {Rupture rocheuse} }
  }
  child[branche3] { node[concept, text width=2.2cm, align=center] {Évaluation}
    child[sousbranche] { node[concept, text width=1.8cm, align=center] {Magnitude (Richter)} }
    child[sousbranche] { node[concept, text width=1.8cm, align=center] {Intensité (Mercalli)} }
    child[sousbranche] { node[concept, text width=1.8cm, align=center] {Sismogramme} }
    child[sousbranche] { node[concept, text width=1.8cm, align=center] {Échelle logarithmique} }
  }
  child[branche4] { node[concept, text width=2.2cm, align=center] {Répartition mondiale}
    child[sousbranche] { node[concept, text width=1.8cm, align=center] {Limites de plaques} }
    child[sousbranche] { node[concept, text width=1.8cm, align=center] {Ceinture de feu Pacifique} }
    child[sousbranche] { node[concept, text width=1.8cm, align=center] {Dorsales médio-océaniques} }
    child[sousbranche] { node[concept, text width=1.8cm, align=center] {Cameroun : Mt Cameroun} }
  }
  child[branche5] { node[concept, text width=2.2cm, align=center] {Prévention}
    child[sousbranche] { node[concept, text width=1.8cm, align=center] {Construction parasismique} }
    child[sousbranche] { node[concept, text width=1.8cm, align=center] {Plan d'urgence familial} }
    child[sousbranche] { node[concept, text width=1.8cm, align=center] {Surveillance sismique} }
    child[sousbranche] { node[concept, text width=1.8cm, align=center] {Prévision improbable} }
  };
\end{tikzpicture}
\legende{Carte mentale récapitulative de la leçon sur les séismes.}
\end{popfigure}

\begin{bilan}
\courssub{Définitions clés}
\begin{itemize}
  \item \cle{Séisme} : Secousse brutale du sol due à une libération soudaine d'énergie à l'intérieur de la Terre.
  \item \cle{Foyer (hypocentre)} : Point profond à l'intérieur de la Terre où naît la rupture rocheuse (profondeur $< 700\,\text{km}$).
  \item \cle{Épicentre} : Projection verticale du foyer à la surface du sol ; zone des dégâts maximaux.
  \item \cle{Ondes sismiques} : Ondes P (primaires, compression, rapides $\approx 6-8\,\text{km/s}$) et ondes S (secondaires, cisaillement, lentes $\approx 3-5\,\text{km/s}$). Les ondes de surface (L, R) causent les dégâts.
  \item \cle{Magnitude} : Mesure de l'énergie libérée à la source (échelle de Richter, logarithmique : 1 unité = 32 fois plus d'énergie).
  \item \cle{Intensité} : Mesure des effets observés en surface (échelle de Mercalli modifiée, I à XII, subjective).
\end{itemize}

\courssub{Repères numériques \& Formules}
\begin{center}
\begin{tabularx}{\linewidth}{|l|X|}\hline
\textbf{Notion} & \textbf{Valeur / Relation} \\ \hline
Vitesse onde P ($v_P$) & $5,5$ à $8,0\,\text{km/s}$ (selon roche) \\ \hline
Vitesse onde S ($v_S$) & $3,0$ à $4,5\,\text{km/s}$ ($v_S \approx 0,6\,v_P$) \\ \hline
Différence temps $\Delta t = t_S - t_P$ & $\Delta t \approx \frac{d}{v_S} - \frac{d}{v_P} \approx \frac{d}{8}$ (avec $d$ en km, $\Delta t$ en s) \\ \hline
Magnitude $M$ (Richter) & $M = \log_{10}(A) - \log_{10}(A_0(\Delta))$ ; $M+1 \Leftrightarrow E \times 32$ \\ \hline
Profondeur foyers & Superficiels ($<70\,\text{km}$), Intermédiaires ($70-300\,\text{km}$), Profonds ($>300\,\text{km}$) \\ \hline
\end{tabularx}
\end{center}

\courssub{Méthodes express}
\begin{enumerate}
  \item \textbf{Localiser un épicentre} : Sur 3 stations minimum, mesurer $\Delta t = t_S-t_P$ $\rightarrow$ calculer distance $d$ (abaque ou formule) $\rightarrow$ tracer 3 cercles (rayon $d$) $\rightarrow$ intersection = épicentre.
  \item \textbf{Lire un sismogramme} : Identifier premier arrivage (onde P) et second arrivage (onde S) $\rightarrow$ mesurer $\Delta t$ $\rightarrow$ déduire distance épicentrale.
  \item \textbf{Comparer Magnitude vs Intensité} : Magnitude = une seule valeur par séisme (énergie) ; Intensité = plusieurs valeurs selon la distance à l'épicentre et la nature du sol (amplification site).
\end{enumerate}

\courssub{Contexte camerounais}
\begin{itemize}
  \item Zone active : Ligne volcanique du Cameroun (Mont Cameroun, Manengouba, lac Nyos).
  \item Séismes historiques : 1986 (lac Nyos, éruption limnique déclenchée par séisme ?), séismes réguliers de faible magnitude sur le Mont Cameroun.
  \item Surveillance : Réseau sismologique national (IRGM, stations de Buéa, Yaoundé, Ngaoundéré).
\end{itemize}
\end{bilan}

\begin{aretenir}[Checklist avant le BEPC]
\begin{itemize}
  \item $\square$ Je définis correctement foyer, épicentre, ondes P et S.
  \item $\square$ J'explique pourquoi les ondes P arrivent avant les ondes S.
  \item $\square$ Je calcule la distance épicentrale $d$ à partir de $\Delta t = t_S-t_P$ (formule ou abaque).
  \item $\square$ Je construis l'épicentre par intersection de 3 cercles (méthode des 3 stations).
  \item $\square$ Je distingue magnitude (énergie, logarithmique, unique) et intensité (effets, subjective, variable selon lieu).
  \item $\square$ Je localise les grandes zones sismiques mondiales (limites de plaques : convergence, divergence, transformante).
  \item $\square$ Je cite la « Ceinture de feu du Pacifique » comme zone la plus active.
  \item $\square$ J'identifie le Mont Cameroun comme zone sismique active au Cameroun.
  \item $\square$ Je sais que la prévision précise (date, lieu, magnitude) est impossible à ce jour.
  \item $\square$ Je cite 3 mesures de prévention : construction parasismique, plan d'urgence, éducation des populations.
  \item $\square$ J'explique l'amplification des dégâts sur sols meubles (alluvions) vs roches dures.
  \item $\square$ Je relie séisme sous-marin $\rightarrow$ tsunami (ex. Japon 2011, Sumatra 2004).
\end{itemize}
\end{aretenir}

\findecours

\end{document}
